\documentclass[11pt]{article}
\usepackage[T1]{fontenc}
\usepackage[utf8]{inputenc}
\usepackage{lmodern}
\usepackage{microtype}
\usepackage[letterpaper,margin=1in]{geometry}
\usepackage{amsmath,amssymb,amsthm,mathtools}
\usepackage{enumitem}
\usepackage[numbers,sort&compress]{natbib}
\usepackage{xurl}
\usepackage[hidelinks]{hyperref}
\usepackage{bookmark}
\input{glyphtounicode}
\input{glyphtounicode-cmex}
\pdfglyphtounicode{mapsto}{007C}
\pdfglyphtounicode{parenleftbig}{0028}
\pdfglyphtounicode{parenrightbig}{0029}
\pdfglyphtounicode{bracketleftbig}{005B}
\pdfglyphtounicode{bracketrightbig}{005D}
\pdfglyphtounicode{braceleftbig}{007B}
\pdfglyphtounicode{bracerightbig}{007D}
\pdfglyphtounicode{vextendsingle}{23D0}
\pdfglyphtounicode{vextenddouble}{2016}
\pdfglyphtounicode{parenleftbigg}{0028}
\pdfglyphtounicode{parenrightbigg}{0029}
\pdfglyphtounicode{braceleftbigg}{007B}
\pdfglyphtounicode{bracerightbigg}{007D}
\pdfglyphtounicode{parenleftBigg}{0028}
\pdfglyphtounicode{parenrightBigg}{0029}
\pdfglyphtounicode{braceleftBigg}{007B}
\pdfglyphtounicode{bracerightBigg}{007D}
\pdfglyphtounicode{circleplustext}{2A01}
\pdfglyphtounicode{circleplusdisplay}{2A01}
\pdfglyphtounicode{circlemultiplydisplay}{2A02}
\pdfglyphtounicode{summationtext}{2211}
\pdfglyphtounicode{producttext}{220F}
\pdfglyphtounicode{summationdisplay}{2211}
\pdfglyphtounicode{productdisplay}{220F}
\pdfglyphtounicode{integraldisplay}{222B}
\pdfglyphtounicode{tildewide}{02DC}
\pdfglyphtounicode{tildewider}{02DC}
\pdfglyphtounicode{hatwide}{02C6}
\pdfglyphtounicode{radicalbig}{221A}
\pdfglyphtounicode{radicalBig}{221A}
\pdfglyphtounicode{radicalbigg}{221A}
\pdfgentounicode=1
\ifdefined\pdfinfoomitdate\pdfinfoomitdate=1\fi
\ifdefined\pdftrailerid\pdftrailerid{}\fi
\ifdefined\pdfsuppressptexinfo\pdfsuppressptexinfo=15\fi
\hypersetup{
  pdftitle={Weak pure infiniteness and O-infinity absorption},
  pdfsubject={Pure infiniteness, exactness, and tensorial absorption of the Cuntz algebra},
  pdfauthor={OpenAI},
  pdfcreator={pdfLaTeX},
  bookmarksnumbered=true,
  pdfdisplaydoctitle=true
}

\newcommand{\C}{\mathbb C}
\newcommand{\R}{\mathbb R}
\newcommand{\N}{\mathbb N}
\newcommand{\Z}{\mathbb Z}
\newcommand{\F}{\mathbb F}
\newcommand{\T}{\mathbb T}
\newcommand{\U}{\mathcal U}
\newcommand{\Ad}{\operatorname{Ad}}
\newcommand{\tr}{\operatorname{tr}}
\newcommand{\id}{\operatorname{id}}
\newcommand{\op}{\mathrm{op}}
\newcommand{\IIone}{\mathrm{II}_1}
\newcommand{\Wstar}{\mathrm{W}^{*}}
\newcommand{\Her}{\operatorname{Her}}
\newcommand{\Id}{\operatorname{Id}}
\newcommand{\Prim}{\operatorname{Prim}}
\newcommand{\CuEq}{\sim_{\mathrm{Cu}}}
\newcommand{\Oinf}{\mathcal O_\infty}
\newcommand{\Kc}{\mathcal K}
\newcommand{\K}{\mathcal K}
\newcommand{\fin}{\operatorname{Fin}}
\newcommand{\ran}{\operatorname{ran}}
\newcommand{\dist}{\operatorname{dist}}
\newcommand{\propP}{\textup{(P)}}
\newcommand{\propF}{\textup{(F\ensuremath{_n})}}
\newcommand{\doi}[1]{\href{https://doi.org/#1}{doi:\,\nolinkurl{#1}}}

\theoremstyle{plain}
\newtheorem{theorem}{Theorem}[section]
\newtheorem{proposition}[theorem]{Proposition}
\newtheorem{lemma}[theorem]{Lemma}
\newtheorem{corollary}[theorem]{Corollary}
\theoremstyle{definition}
\newtheorem{definition}[theorem]{Definition}
\theoremstyle{remark}
\newtheorem{remark}[theorem]{Remark}
\numberwithin{equation}{section}
\setlist[enumerate]{itemsep=3pt,topsep=6pt}
\setlist[itemize]{itemsep=3pt,topsep=6pt}
\setlength{\emergencystretch}{1.5em}
\allowdisplaybreaks[1]
\widowpenalty=8000
\clubpenalty=8000

\title{Weak pure infiniteness and \texorpdfstring{\(\Oinf\)}{O-infinity} absorption}
\author{OpenAI}
\date{September 25, 2026}

\begin{document}
\maketitle

\begin{abstract}
We prove that every complex \(C^*\)-algebra in which each positive
element is properly infinite is strongly purely infinite. This answers
the ordinary-to-strong part of Kirchberg--R{\o}rdam's comparison question.
For exact algebras, we also prove that proper infiniteness of one fixed
finite amplification of every positive element implies proper infiniteness
of each positive element. Consequently, every separable nuclear algebra
with this fixed-amplification property absorbs \(\Oinf\), without
assumptions of unitality or simplicity.
\end{abstract}

\section{Introduction}\label{sec:introduction}

Pure infiniteness expresses the ability of positive elements of a
\(C^*\)-algebra to contain several copies of themselves in Cuntz
comparison.  For a nonsimple algebra, an individual comparison and
the simultaneous comparison of two elements are different demands:
the latter must also control a mixed entry.  We prove that proper
infiniteness of every positive element already provides this
compatibility.  We then give two ways to remove a fixed finite
amplification, one for exact algebras and an independent construction
for nuclear algebras.  In the separable nuclear case the result yields
absorption of the Cuntz algebra \(\Oinf\).

\subsection{Infiniteness notions and main results}

All \(C^*\)-algebras are complex and may be nonunital.  Let \(\Kc\)
denote the compact operators on an infinite-dimensional separable
Hilbert space.  For \(u,v\in(C\otimes\Kc)_+\), write \(u\precsim v\)
if some sequence \(x_j\in C\otimes\Kc\) satisfies
\(x_j^*vx_j\to u\) in norm, and write \(u\CuEq v\) if comparison holds
in both directions.  This is stable Cuntz comparison
\cite{Cuntz1978}.  Finite matrices over \(C\) occupy the standard
corners of \(C\otimes\Kc\); comparisons between different sizes use
zero padding in the other coordinates.  Write \(h^{\oplus n}\) for
the direct sum of \(n\) copies of \(h\).  A positive element \(h\) is
\emph{properly infinite} if \(h\oplus h\precsim h\).  We include zero
as properly infinite.  For \(h\in C_+\), let
\[
 \Her(h)=\overline{hCh},\qquad
 \Id(h)=\text{the closed two-sided ideal of \(C\) generated by \(h\)}.
\]

We will use the following two comparison conditions:
\begin{align}
 &h\text{ is properly infinite for every }h\in C_+,
 \tag{P}\label{eq:property-P}\\
 &h^{\oplus n}\text{ is properly infinite for every }h\in C_+,
 \quad\text{for one fixed integer }n\ge1.
 \tag{\(F_n\)}\label{eq:property-F}
\end{align}
In \(\propF\), the same \(n\) is chosen before all positive elements
are tested.  It therefore applies to every scalar spectral cutoff
used later.  The premise concerns scalar positives in \(C\), with
their indicated amplifications.

Condition \(\propP\) is equivalent to pure infiniteness in the
nonsimple sense of Kirchberg--R{\o}rdam
\cite[Theorem~4.16]{KR2000}; the equivalence is also recalled in
\cite[Section~21]{STW99problems}.  The existence of an
\(n\) satisfying \(\propF\) is \emph{weak pure infiniteness} in the
fixed-amplification sense.  We use the following formulation of
\emph{strong pure infiniteness}, valid for arbitrary \(C^*\)-algebras
\cite[Definition~3.1]{KS}: for every \(x,y\in C_+\) and
\(\varepsilon>0\), there are \(s,t\in C\) such that
\begin{equation}\label{eq:SPI-definition}
 \|s^*x^2s-x^2\|<\varepsilon,\qquad
 \|t^*y^2t-y^2\|<\varepsilon,\qquad
 \|s^*xy t\|<\varepsilon.
\end{equation}
The coefficients in this formulation have no prescribed norm bound.

Kirchberg--R{\o}rdam introduced the weak and strong notions in
\cite{KR2002}.  Strong pure infiniteness implies ordinary pure
infiniteness, and \(\propP\) implies \(\propF\) with \(n=1\).
The reverse implications form the comparison question in
Kirchberg--R{\o}rdam's Question~9.5, as recalled in
\cite[Section~21, Problem~LXXII]{STW99problems}.  In version~2 of that problem collection, the question asks whether
the three notions agree in general and in the nuclear setting.
Our first theorem proves the implication from ordinary to strong
pure infiniteness for arbitrary complex \(C^*\)-algebras, in the
following more flexible form.

\begin{theorem}\label{thm:individual-to-strong}
Let \(C\) be a complex \(C^*\)-algebra satisfying \(\propP\).
For all \(a,b\in C_+\), \(c\in C\), and \(\varepsilon>0\), there are
\(s,t\in C\) such that
\begin{equation}\label{eq:main-diagonalization}
 \|s^*as-a\|<\varepsilon,\qquad
 \|t^*bt-b\|<\varepsilon,\qquad
 \|s^*ct\|<\varepsilon.
\end{equation}
Consequently \(C\) is strongly purely infinite.
\end{theorem}

The cross term \(c\) is arbitrary and need not be a product of the
diagonal inputs.  Taking \((a,b,c)=(x^2,y^2,xy)\) gives the stated
consequence.  Arbitrary-entry diagonalization is a known equivalent
form of strong pure infiniteness
\cite[Remarks~3.2(i)--(ii)]{KS}; our proof obtains it directly from
\(\propP\).  It also gives coefficient bounds for positive
contraction inputs that are independent of the requested mixed-entry
accuracy; see Proposition~\ref{prop:bounded-diagonalization}.

Recall that a \(C^*\)-algebra is exact if its minimal tensor product
preserves short exact sequences.  Exactness supplies the additional
uniformity needed to remove the fixed amplification.

\begin{theorem}\label{thm:exact-infiniteness}
Let \(C\) be an exact complex \(C^*\)-algebra.  Suppose that there is
one integer \(n\ge1\) such that \(h^{\oplus n}\) is properly infinite
for every \(h\in C_+\).  Then every positive element of \(C\) is
properly infinite.  Consequently \(C\) is strongly purely infinite.
\end{theorem}

These theorems give a positive answer to the ordinary-to-strong part
of Problem~LXXII in general and to the full comparison question for
exact algebras, hence also for nuclear algebras.  The weak-to-ordinary implication
for arbitrary nonexact algebras remains outside these conclusions.
The statements include nonseparable and nonunital algebras.  Section~\ref{sec:nuclear}
gives an independent proof of the first conclusion of
Theorem~\ref{thm:exact-infiniteness} under nuclearity.  Its statement
is a restriction of that theorem, while its finite-tuple obstruction
and merger of hereditary witnesses provide a different method.
The same method extends to exact algebras by the deeper theorem that
exactness passes to quotients; see Remark~\ref{rem:tuple-exact}.
The direct exact proof avoids that permanence theorem.

Let \(\Oinf\) be the Cuntz algebra generated by countably many
isometries with mutually orthogonal ranges.  Kirchberg--R{\o}rdam
proved the equivalence between strong pure infiniteness and
\(\Oinf\)-absorption for separable nuclear algebras that are stable
or have an approximate unit of projections \cite{KR2002}.
Kirchberg's later central-sequence results cover general nonunital
separable nuclear algebras
\cite[Proposition~4.4 and Remark~4.6]{Kirchberg2006}; see also the
paragraph following \cite[Definition~3.1]{KS}.  At that scope they
give the following tensorial consequence.

\begin{corollary}\label{cor:absorption}
Let \(A\) be a separable nuclear complex \(C^*\)-algebra.  If there is
one integer \(n\ge1\) such that \(h^{\oplus n}\) is properly infinite
for every \(h\in A_+\), then
\[
 A\cong A\otimes_{\min}\Oinf.
\]
\end{corollary}

For separable nuclear algebras, therefore, weak pure infiniteness,
pure infiniteness, strong pure infiniteness, and \(\Oinf\)-absorption
are equivalent.  The reverse implication from absorption to strong
pure infiniteness is due to Kirchberg--R{\o}rdam
\cite{KR2002}; see also the discussion after
\cite[Definition~3.1]{KS}.  Thus a single finite amplification test
on positive elements characterizes this tensorial regularity even
when the algebra is nonsimple or nonunital.

For example, the cone \(C_0((0,1])\otimes\Oinf\) is separable
nuclear and absorbs \(\Oinf\) by the self-absorption of \(\Oinf\).
It is nonsimple, nonunital, and has no nonzero projections: the norm
of a continuous projection-valued function takes only the values
zero and one and must tend to zero at the missing endpoint.

\subsection{Earlier results and the role of ideal structure}

The comparison question has been settled under several structural
conditions.  Kirchberg--R{\o}rdam treat simple, real-rank-zero, and
approximately divisible cases \cite{KR2002}; the real-rank-zero
conclusion is also recalled in \cite[Section~21]{STW99problems}.
Blanchard--Kirchberg prove that ordinary pure infiniteness implies
strong pure infiniteness when the primitive ideal space is Hausdorff
\cite[Theorem~5.8]{BlanchardKirchberg2004}.  R{\o}rdam's proof of
\cite[Theorem~5.2]{Rordam2004Z} gives the same implication for
\(\mathcal Z\)-absorbing purely infinite algebras, where
\(\mathcal Z\) is the Jiang--Su algebra.  These two implications do
not require exactness or nuclearity.

Pasnicu--R{\o}rdam establish equivalence of the three infiniteness
notions for algebras with the ideal property, meaning that projections
separate ideals \cite[Proposition~2.14]{PasnicuRordam2007}.
For algebras of topological dimension zero, Pasnicu--Rouzbehani prove
the implication from ordinary to strong pure infiniteness
\cite[Theorems~3.10(i) and~3.13]{PasnicuRouzbehani2025}.
Here topological dimension zero means that the primitive ideal space
has a basis of compact open sets.  Theorem~\ref{thm:individual-to-strong}
requires no restriction on this space or on the ideal structure.

The weak-to-ordinary implication is closely related to the problem of
finding sufficiently many orthogonal equivalent elements in a
hereditary subalgebra.  Kirchberg--R{\o}rdam showed that the Global
Glimm Halving property permits this implication
\cite[Proposition~4.15]{KR2002}; see
\cite[Theorem~39]{STW99problems}.  An equivalent formulation of the
Global Glimm Property is that, for every \(a\in C_+\) and
\(\varepsilon>0\), there is \(x\in\Her(a)\) with \(x^2=0\) such
that \((a-\varepsilon)_+\) belongs to the closed two-sided ideal
generated by \(x\) \cite[Remark~3.2 and Theorem~3.6(3)]{TV2023}.
Elliott--Rouzbehani prove weak-to-ordinary pure infiniteness in
topological dimension zero \cite{ER2023}.  Ng--Thiel--Vilalta recover
that implication through the Global Glimm Property in a proof that
includes nonseparable algebras
\cite[Theorem~2.3 and Corollary~2.5]{NTV2026}.
The initial and range squares of a square-zero element are
orthogonal and Cuntz equivalent.  Our exact argument constructs such
elements while retaining a nonzero image throughout a prescribed
set of primitive quotients.  The relevant sets are the norm slices
\(\{P\in\Prim(C):\|h+P\|\ge r\}\), where
\(h\in C_+\) and \(r>0\); they are quasi-compact even when
\(\Prim(C)\) is not Hausdorff.  Repeated halving on the same slice
then supplies enough copies to remove the fixed amplification.  The nuclear argument instead combines
representation-dependent positive elements into one element that
covers a prescribed cutoff.

\subsection{Proof strategy and the independent routes}

Two tasks drive the argument.  To pass from \(\propP\) to strong
pure infiniteness, one must choose two hereditary corners that retain
the relevant parts of their inputs while making the mixed term small.
To remove \(n\) from \(\propF\), one must first produce sufficiently
many equivalent orthogonal pieces inside a scalar hereditary
subalgebra.  Figure~\ref{fig:proof-routes} separates these tasks and
the two methods for the second one.

\begin{figure}[htbp]
\centering
\small
\fbox{\strut Comparison and cutoff transport (Section~\ref{sec:foundations})}
\par\medskip
\newcommand{\routebox}[1]{\fbox{\parbox[c][13mm][c]{0.27\textwidth}{\centering #1}}}
\begin{tabular}{@{}c@{\hspace{3mm}}c@{\hspace{3mm}}c@{}}
\textbf{Arbitrary algebra, \(\propP\)}&
\textbf{Nuclear algebra, \(\propF\)}&
\textbf{Exact algebra, \(\propF\)}\\[3pt]
\routebox{Graph and four entries\\Sections~\ref{sec:graph}--\ref{sec:four-variables}}&
\routebox{Graph and finite tuples\\Sections~\ref{sec:graph}, \ref{sec:nuclear}}&
\routebox{Uniform perturbation\\Section~\ref{sec:exactness}}\\
\(\Downarrow\)&\(\Downarrow\)&\(\Downarrow\)\\
\routebox{Products and iteration\\Section~\ref{sec:separation}}&
\routebox{Local elements and merger\\Section~\ref{sec:nuclear}}&
\routebox{Square-zero localization\\Section~\ref{sec:localization}}\\
\(\Downarrow\)&\(\Downarrow\)&\(\Downarrow\)\\
Strong pure infiniteness&\(\propP\)&\(\propP\)
\end{tabular}
\par\medskip
\(\propP\ \xRightarrow{\ \text{general route}\ }\
  \text{strong pure infiniteness}
  \ \xRightarrow{\ \text{separable nuclear}\ }\
  \Oinf\text{-absorption}\)
\caption{Proof routes and their hypotheses.  Comparison transport
supports all three routes.  Either output \(\propP\) can enter the
left-hand route.  Remark~\ref{rem:tuple-exact} extends the middle
route to exact algebras using quotient permanence.  Under exactness, Section~\ref{sec:quantitative}
gives an alternative passage from \(\propP\) to strong pure
infiniteness.  The final arrow is the external absorption criterion
used in Corollary~\ref{cor:absorption}.}
\label{fig:proof-routes}
\end{figure}


\paragraph{Comparison transport.}
Section~\ref{sec:foundations} supplies the common tool.  A Cuntz
comparison, after a spectral cutoff, can be implemented by a
rectangular partial isometry in the bidual.  Its products with the
cutoff hereditary algebra lie in the original algebra, and conjugation
preserves norms in every representation, also at matrix levels.
Thus one can place several orthogonal copies in a hereditary
subalgebra and transport an entire finite tuple through the same map.
These are stronger implementation data than a comparison of classes
alone.

\paragraph{From individual to simultaneous comparison.}
For \(h\ge0\), let \(e_h\) be its support projection in the universal
bidual.  Starting from positive contractions \(a,b\) and a contraction
\(c\), the first goal is to find positive contractions
\(p\in\Her(a)\), \(q\in\Her(b)\) with
\[
 \|e_pce_q\|<1,
\]
while requiring represented norm one of \(p\) or \(q\) whenever the
corresponding input has represented norm greater than \(1/2\).
This retention of represented norms is what later gives ideal
comparison.

Section~\ref{sec:graph} builds one corner as a graph inside several
orthogonal copies of a cutoff of \(a\).  The graph has finitely many
freely chosen algebra entries.  If the positive defect \(b(1-c^*p^2c)b\) vanishes in a
representation that detects the relevant part of \(b\), equality in
contraction inequalities forces all of
those entries to equal fixed measurements, computed before the entries
were chosen.  Under \(\propP\), Section~\ref{sec:four-variables}
excludes these equalities with four entries.  The first pair forces
an isometry with a nonzero defect, using the scaling-element
construction of Blackadar--Cuntz
\cite{BlackadarCuntz1982,PasnicuRordam2007}; the second pair is forced
both into and away from that defect.  Spectral cutoffs place the
relevant supports in actual operator ranges, and a directed family
of finite tests makes the contradiction work without separability.
Two ideal cutoffs then give the strict support estimate above.

Section~\ref{sec:separation} upgrades strict contraction to arbitrarily
small compression.  Uniformly bounded comparison rows show that
\(\propP\) passes to bounded products.  A product contradiction
therefore supplies one contraction factor below one, and iteration
preserves both represented norm signals while reducing the mixed
compression.  Finally, high spectral parts of \(a,b\) are compared
into the separated corners.  The coefficients have bounds determined
by the diagonal cutoffs, independent of the requested mixed-entry
accuracy.  This order of choices completes
Theorem~\ref{thm:individual-to-strong}.

\paragraph{A separate nuclear proof of amplification removal.}
Section~\ref{sec:nuclear} reuses the graph construction with a
different obstruction.  A finite tuple witnessing nonexactness has
linear matrix norms that no tuple in an exact algebra can realize
at every matrix level.  Finite-dimensional compressions nevertheless
approximate each fixed norm.  Under \(\propF\), one rectangular
carrier transports each entire compressed tuple into \(M_n(C)\).
If every choice failed to exclude the graph relations, their limiting
norms would occur in a quotient of \(M_n(C)\).  Nuclearity makes this
quotient nuclear, hence exact, giving a contradiction.  The matrix
norm and fixed-element quotient arguments are related to Pisier's
operator-space localization methods \cite[pp.~162--165]{Pisier1995}.

It remains to obtain the positive inputs that the graph needs and
combine them.  The fixed-amplification hypothesis and Kadison
transitivity \cite{Kadison1957} give, at each irreducible
representation detecting a positive \(u\), a positive contraction in \(\Her(u)\)
whose every spectral cutoff is properly infinite and whose norm in
that representation is one.  This element may
depend on the representation.  Repeated graph shrinking and a polar
deformation merge finitely many such elements while retaining all
their norm signals.  A finite ideal argument covers each cutoff of
\(u\) by one merged element and then compares two copies of that
cutoff back to \(u\).  This proves \(\propP\) under nuclearity,
independently of the exact argument.  The use of hereditary elements
to cover ideals is related to the filling-family criteria of
\cite[Theorem~1.1 and Definitions~4.2,~4.7]{KS}; here the merger supplies
the simultaneous norm control.

\paragraph{The exact proof and its quantitative completion.}
Section~\ref{sec:exactness} starts from the nonexact universal unital
\(C^*\)-algebra of one contraction.  Its finite matrix tests and its
embedded copy of \(C^*(\mathbb F_2)\) were studied by
Courtney--Sherman \cite{CS}, with Wassermann's nonexact free-group
example as a predecessor \cite{Wassermann1976}.  Finite-word
compression at all matrix levels turns nonexactness into one positive
error threshold: no contraction in a unital exact algebra can contain
approximately reducing copies of every finite matrix contraction.
The same copy must approximately intertwine both the contraction and
its adjoint.

Combining this obstruction with \(\propF\) gives a perturbation that
prevents cancellation simultaneously in every primitive quotient of a
prescribed norm slice.  Its lower norm bound is proportional to the
perturbation size, with a constant independent of the slice and of
\(n\).  Section~\ref{sec:localization} applies this estimate through
three finite coupling steps.  For a hereditary subalgebra nonzero
on a prescribed quasi-compact subset of \(\Prim(C)\), it constructs
a square-zero element nonzero at every point of that subset.  Repeated
halving on one fixed norm slice puts \(n\) copies of a positive
element below the input in Cuntz comparison; ideal comparison and a
cutoff limit prove
\(\propP\).  The general theorem then completes
Theorem~\ref{thm:exact-infiniteness}.

Section~\ref{sec:quantitative} gives an alternative final step under
exactness.  It constructs positives in two prescribed hereditary
subalgebras with arbitrarily small product, while retaining a quotient
norm floor fixed before the product error is chosen.  Two compatible
inverse perturbations preserve the two floors.  Clipping the later
comparison coefficients bounds their norms, so the small product
survives implementation and gives strong pure infiniteness directly.
Finally, for a separable nuclear algebra, the central-sequence
absorption criterion applied in Section~\ref{sec:nuclear} gives
Corollary~\ref{cor:absorption}.

\section{Comparison and transport}\label{sec:foundations}

The diagonalization argument repeatedly places several copies of a
positive element inside one hereditary subalgebra.  Comparison of
Cuntz classes alone does not record enough information for this use:
the copies must have a common implementation, their products must
remain in the algebra, and their norms must agree in every
representation.  We establish these facts here, together with the
ideal and cutoff rules used to retain norm signals.

Cuntz introduced the comparison relation in \cite{Cuntz1978}.
The proper-target ideal comparison in
Lemma~\ref{lem:comparison-rules}(iii) is the standard consequence
recorded in \cite[Proposition~3.5(ii)]{KR2000}; it is also used with
that reference in
\cite[proof of Proposition~2.38]{KwasniewskiMeyer2023}.
The cutoff implementation below is related to R{\o}rdam's
comparison lemma \cite[Proposition~2.4]{Rordam1992}, in the form
recalled in \cite[Section~2.1]{AntoinePereraThiel2020}.
We give the rectangular and representation statements needed here
with their proofs.

All comparisons below are stable: finite positive matrices are
placed in a common stabilization by padding with zeros.  For a
\(C^*\)-algebra \(D\) and \(h\in D_+\), let \(\Id_D(h)\)
denote the closed two-sided ideal generated by \(h\), and put
\[
 \Her_D(h)=\overline{hDh},\qquad
 e_h=\mathbf 1_{(0,\infty)}(h)\in D^{**}.
\]
We omit the subscript when the ambient algebra is clear.  Every
support projection is taken in the universal bidual.  Rectangular
matrices over \(C^{**}\) are viewed in the corresponding corners of
a matrix algebra over \(C^{**}\).  The normal extension of a
representation sends \(e_h\) to \(e_{\pi(h)}\): apply normality to
\(h(h+1/j)^{-1}\uparrow e_h\).  When using a unit or a support
projection in a degenerate representation, we restrict to its
essential subspace.  This leaves the norm of every represented
algebra element unchanged.

\begin{lemma}[Elementary comparison rules]\label{lem:comparison-rules}
Let \(D\) be a \(C^*\)-algebra.
\begin{enumerate}[label=\textup{(\roman*)}]
\item For \(h\in D_+\),
\[
 \Her_D(h)=D\cap e_hD^{**}e_h.
\]
If \(0\le u\le K v\) for \(u,v\in D_+\) and \(K\ge0\), then
\(u\in\Her(v)\) and \(u\precsim v\).  More generally, every
positive element of \(\Her(v)\) is Cuntz subequivalent to \(v\).
\item For a rectangular algebra matrix \(z\), one has
\(z^*z\CuEq zz^*\).  Cuntz subequivalence is transitive, is
preserved by block sums, and is closed under norm limits in its
first variable when the second variable is fixed.  Every
\(*\)-homomorphism preserves Cuntz subequivalence and equivalence,
after amplification.  For integers \(m,n\ge1\), \(d\in M_n(D)_+\),
and \(v\in M_m(D)_+\), stable comparison \(d\precsim v\) is
equivalent to the existence of \(r_j\in M_{m,n}(D)\) with
\(r_j^*vr_j\to d\) in norm.
\item If \(d,v\) are positive elements in a common stabilization,
\(d\in\Id(v)\), and \(v\) is properly infinite, then
\(d\precsim v\).
\end{enumerate}
\end{lemma}

\begin{proof}
We first prove the bidual description in (i), including its norm
closure assertion.  One inclusion follows on \(hDh\) and then by
closure.  Conversely, suppose \(u\in D_+\) and \(e_hu=u\).
Take increasing continuous positive contractions \(f_j\) with
\(f_j(0)=0\) and \(f_j(t)\uparrow1\) for \(t>0\).  On the compact
state space of the unitization, the continuous functions
\[
 \varphi\longmapsto
 \varphi\bigl(u(1-f_j(h))u\bigr)
\]
decrease pointwise to zero.  Indeed, the normal extension of each
state evaluates the strong limit \(f_j(h)\uparrow e_h\).
Dini's theorem makes this convergence uniform.  The positive norm
formula and \((1-f_j(h))^2\le1-f_j(h)\) then give
\(\|(1-f_j(h))u\|\to0\).  Thus
\(u=\lim_j f_j(h)u f_j(h)\in\Her(h)\).
Here \(f_j(h)\in\Her(h)\) follows from continuous functional
calculus, by first cutting \(f_j\) off near zero.  For a general
element supported on both sides by \(e_h\), apply the positive case
to the positive and negative parts of its real and imaginary parts.

If \(0\le u\le Kv\), compression by \(1-e_v\) gives
\(u^{1/2}(1-e_v)=0\).  Hence \(u\) is supported by \(e_v\) and
belongs to \(\Her(v)\).  If \(u\in\Her(v)_+\), take a cutoff
approximate identity \(g_j(v)\) for \(\Her(v)\), with every \(g_j\)
zero near zero.  The quotient function below is continuous when
extended by zero at zero, and gives elements
\[
 x_j=\frac{g_j(v)}{\sqrt v}\,u^{1/2}\in D,\qquad
 x_j^*v x_j=u^{1/2}g_j(v)^2u^{1/2}\longrightarrow u.
\]
This proves the comparison assertion.

For (ii), the rectangular coefficient
\(z(z^*z+1/j)^{-1/2}\) compares \(z^*z\) to \(zz^*\), since its
conjugation of \(zz^*\) equals
\[
 (z^*z)^2(z^*z+1/j)^{-1}\longrightarrow z^*z.
\]
Apply the same calculation to \(z^*\) for the reverse comparison.
The inverses lie in unitizations and are multiplied by algebra
matrices.  Block-diagonal coefficients prove preservation under
block sums.  For transitivity, first approximate \(u\) by \(x^*vx\),
then approximate \(v\) by \(y^*wy\) within the requested error
divided by \(\max\{1,\|x\|^2\}\); the coefficient \(yx\) compares
\(u\) to \(w\).  A diagonal choice of coefficients proves closure
under norm limits in the first variable.

If \(\phi:D\to E\) is a \(*\)-homomorphism, apply
\(\phi\otimes\id_{\Kc}\) to a comparison sequence; continuity
preserves its norm limit.  Applying this in both directions preserves
equivalence.  For finite \(d,v\), let \(P_n,P_m\) be the scalar
coordinate projections in the multiplier algebra of \(D\otimes\Kc\).
If \(x_j^*vx_j\to d\) in the stabilization, then
\[
 r_j=P_mx_jP_n\in M_{m,n}(D),\qquad
 r_j^*vr_j=P_nx_j^*vx_jP_n\longrightarrow d.
\]
Conversely, rectangular coefficients embed in the stabilization and
give the same convergence there.  The multiplier compressions retain
algebra membership without a unit, and the argument includes zero
matrices.

For (iii), work in the common stabilization \(D\).  The closed span
of \(Dv^{1/2}D\) is \(\Id_D(v)\), so finite sums
\(\sum_{i=1}^r x_i v^{1/2}y_i\) approximate \(d^{1/2}\).
For such a sum, define
\[
 X=[x_1\ \cdots\ x_r]\in M_{1,r}(D),\qquad
 Y=[y_1\ \cdots\ y_r]^T\in M_{r,1}(D),
\]
and put \(V=\operatorname{diag}(v,\ldots,v)\in M_r(D)\).  Then
\[
 z=\sum_{i=1}^r x_i v^{1/2}y_i=X V^{1/2}Y.
\]
The order and rectangular rules already proved give
\[
 z^*z\le\|X\|^2Y^*VY,\qquad
 Y^*VY\CuEq V^{1/2}YY^*V^{1/2}\le\|Y\|^2V.
\]
It follows that \(z^*z\precsim V\).  Proper infiniteness, block
addition, and induction give \(V\precsim v\) for every finite \(r\).
Now \(z^*z\to d\), and the norm-limit rule yields \(d\precsim v\).
This proof includes \(v=0\).
\end{proof}

In particular, Cuntz equivalence preserves proper infiniteness, and
a block sum of properly infinite positives is properly infinite.
These statements follow from block addition and transitivity.
Also, \(u\precsim v\) implies \(\Id(u)\subset\Id(v)\), since
each comparison approximant belongs to \(\Id(v)\).

\begin{lemma}[Transport after a cutoff]\label{lem:transport}
Let \(C\) be a \(C^*\)-algebra, let \(m,n\ge1\) be integers,
and let \(d\in M_n(C)_+\) and \(v\in M_m(C)_+\).
Suppose \(d\precsim v\) in \(C\otimes\Kc\).  For
\(\varepsilon>0\), put
\(h=(d-\varepsilon)_+\).  There is \(w\in M_{m,n}(C^{**})\)
such that
\begin{equation}\label{eq:transport-supports}
 w^*w=e_h,\qquad ww^*\le e_v,
\end{equation}
and
\begin{equation}\label{eq:transport-membership}
 w\Her_{M_n(C)}(h)\subset M_{m,n}(C),\qquad
 w\Her_{M_n(C)}(h)w^*\subset\Her_{M_m(C)}(v).
\end{equation}
Conjugation by \(w\) is an injective \(*\)-homomorphism on
\(\Her(h)\).  For \(l\in\Her(h)_+\),
\begin{equation}\label{eq:transport-cutoffs}
 (wlw^*-\tau)_+=w(l-\tau)_+w^*,\qquad
 w(l-\tau)_+w^*\CuEq(l-\tau)_+\quad(\tau\ge0),
\end{equation}
and
\begin{equation}\label{eq:transport-hereditary-image}
 w\Her(l)w^*=\Her(wlw^*).
\end{equation}
Every representation, with matching matrix amplifications, satisfies
\begin{equation}\label{eq:transport-represented-norm}
 \|\pi^{(m)}(wkw^*)\|=\|\pi^{(n)}(k)\|
 \qquad(k\in\Her(h)).
\end{equation}
The products in these formulas are algebra elements; \(w\) itself
is only asserted to lie in the bidual.
\end{lemma}

\begin{proof}
Compress a stabilization comparison coefficient to the first \(m\)
target and first \(n\) source coordinates.  These scalar multiplier
projections give \(s\in M_{m,n}(C)\) with
\(\|d-s^*vs\|<\varepsilon/2\).  Put \(z=v^{1/2}s h\).
Order in the unitization and functional calculus for \(d\) give
\begin{equation}\label{eq:transport-order}
 \frac{\varepsilon}{2}h^2
 \le z^*z=h s^*vs h
 \le(\|d\|+\varepsilon/2)h^2.
\end{equation}
For the lower bound, use \(h d h\ge\varepsilon h^2\).
The two inequalities and Lemma~\ref{lem:comparison-rules}(i) show
\(\Her(z^*z)=\Her(h)\).  Consequently
\(\Her(|z|)=\Her(h)\) and \(e_{z^*z}=e_h\).
Let \(z=w|z|\) be the rectangular polar decomposition in the
bidual.  Then \(w^*w=e_h\); the factorization through \(v^{1/2}\)
gives \(ww^*=e_{zz^*}\le e_v\).

For \(k\in\Her(h)=\Her(|z|)\), take cutoff functions \(g_j\) for
\(|z|\), zero near zero, with \(g_j(|z|)k\to k\).  Since
\[
 w g_j(|z|)=z\,\frac{g_j(|z|)}{|z|}\in M_{m,n}(C),
\]
norm closure gives \(wk\in M_{m,n}(C)\).  If \(k\ge0\), then
\(wkw^*=(wk^{1/2})(wk^{1/2})^*\) belongs to \(M_m(C)\).
It is supported under \(e_v\), hence belongs to \(\Her(v)\).
Linear decomposition handles arbitrary \(k\).
The identity \(w^*w=e_h\) proves multiplicativity and injectivity.
It also proves \eqref{eq:transport-cutoffs}: functional calculus
vanishing at zero commutes with the homomorphism, and the two
squares of \(w(l-\tau)_+^{1/2}\) are the asserted positives.

For \eqref{eq:transport-hereditary-image}, normal functional
calculus gives \(e_{wlw^*}=we_lw^*\).  Thus the image of
\(\Her(l)\) is supported under \(e_{wlw^*}\), which gives one
inclusion by Lemma~\ref{lem:comparison-rules}(i).  A sandwich
by the image of \(l\) has the form
\[
 (wlw^*)b(wlw^*)
 =w\bigl((wl)^*b(wl)\bigr)w^*.
\]
The inner algebra element is supported on both sides by \(e_l\),
so lies in \(\Her(l)\).  Since conjugation is an injective
\(*\)-homomorphism on \(\Her(l)\), its image is norm closed;
taking closures gives the reverse inclusion.  Finally, the normal
extension of \(\pi\) carries
\eqref{eq:transport-supports} to the same identities with
represented supports.  Conjugating back by the represented \(w^*\)
recovers \(\pi^{(n)}(k)\).  The two contractive conjugations give
both inequalities in \eqref{eq:transport-represented-norm}.
\end{proof}

If \(d_0\in M_\ell(C)_+\) and \(d=d_0^{\oplus r}\in
M_{r\ell}(C)_+\), split \(w\in M_{m,r\ell}(C^{**})\) into
block columns \(w_i\in M_{m,\ell}(C^{**})\).  Then
\[
 w_i^*w_j=\delta_{ij}e_{(d_0-\varepsilon)_+}.
\]
The \(r\) copies have orthogonal ranges and the same initial support.
The single rectangular map also transports every entry of a matrix
over the common hereditary algebra.  Returning to the notation of
Lemma~\ref{lem:transport}, for \(j\ge1\) the same support
identities apply to
\(1_j\otimes w\in M_{jm,jn}(C^{**})\).  Thus, for every
\(K\in M_j(\Her_{M_n(C)}(h))\), entrywise conjugation gives an
algebra matrix and preserves its norm in every representation, with
the matching \(jm\) and \(jn\) amplifications.  This is the joint
norm assertion needed when a finite tuple is transported.

For orthogonal \(c_1,\ldots,c_r\in C_+\), the row
\(R=[c_1^{1/2},\ldots,c_r^{1/2}]\in M_{1,r}(C)\) satisfies
\[
 R^*R=\operatorname{diag}(c_1,\ldots,c_r),\qquad
 RR^*=\sum_i c_i.
\]
Thus their block sum and their sum in \(C\) are Cuntz equivalent.
Moreover, \(f(\sum_i c_i)=\sum_i f(c_i)\) for every continuous
\(f\) with \(f(0)=0\), first for polynomials with zero constant term
and then by continuous functional calculus.  Apply this calculation
in \(M_m(C)\) to the orthogonal transported copies
\(\theta_i(y)=w_i y w_i^*\) of
\(y\in\Her_{M_\ell(C)}((d_0-\varepsilon)_+)_+\).
The cutoff and comparison assertions in
Lemma~\ref{lem:transport} give
\begin{equation}\label{eq:orthogonal-copy-cutoffs}
 \left(\sum_{i=1}^r\theta_i(y)-\tau\right)_+
 \CuEq ((y-\tau)_+)^{\oplus r}\qquad(\tau\ge0).
\end{equation}
In particular this includes the zero cutoff.

\begin{lemma}[Ideals and spectral cutoffs]\label{lem:ideals-cutoffs}
Let \(D\) be a \(C^*\)-algebra.
\begin{enumerate}[label=\textup{(\roman*)}]
\item If \(J=\overline{\bigcup_\lambda J_\lambda}\) for an
increasing directed family of closed ideals and \(z\in J_+\), then
for every \(\varepsilon>0\), \((z-\varepsilon)_+\) lies in some
\(J_\lambda\).  In particular, the ideals
\(\Id((z-\delta)_+)\), directed as \(\delta\downarrow0\), densely
exhaust \(\Id(z)\).
\item For \(u,v\in D_+\), one has \(u\in\Id(v)\) if and only if
every irreducible representation \(\pi\) with \(\pi(v)=0\)
also satisfies \(\pi(u)=0\).
\item If \(f,g\) are continuous nonnegative functions of a
positive element \(a\), vanish at zero, and have the same open
positivity set in \(\operatorname{sp}(a)\), then
\(\Her(f(a))=\Her(g(a))\) and \(f(a)\CuEq g(a)\).
\end{enumerate}
\end{lemma}

\begin{proof}
For (i), choose \(b\in J_\lambda\) with \(\|z-b\|<\varepsilon\).
Then \(\|z+J_\lambda\|<\varepsilon\), and functional calculus in
the quotient gives \((z-\varepsilon)_+\in J_\lambda\).
Since \((z-\delta)_+\to z\), the particular exhaustion follows.
For (ii), apply separation of points by irreducible representations
to the quotient \(D/\Id(v)\).  For (iii), the two support
projections coincide.  Lemma~\ref{lem:comparison-rules}(i) gives
equality of the hereditary algebras and mutual comparison.
\end{proof}

\begin{definition}\label{def:rich}
A positive contraction \(a\) is \emph{rich} if
\((a-\varepsilon)_+\) is properly infinite for every
\(\varepsilon\ge0\).
\end{definition}

Zero is rich, and under \(\propP\) every positive contraction is
rich.  A \emph{ramp} is a continuous nondecreasing function with
values in \([0,1]\), zero on a neighborhood of zero and constant
at \(1\) beyond a specified larger threshold.  A ramp of a rich
contraction is rich.  Indeed, for \(0\le\tau<1\), the open
positivity set of \((f(a)-\tau)_+\) is either empty or the
positivity set of \((a-\rho)_+\) for some \(\rho\ge0\).
Lemma~\ref{lem:ideals-cutoffs}(iii) and richness apply; for
\(\tau\ge1\) the cutoff is zero.  Lemma~\ref{lem:transport}
gives the same conclusion for transported ramps.

\section{Graph measurements and conditional shrinking}\label{sec:graph}

For a \(C^*\)-algebra \(C\), write \(\widetilde C\) for its
minimal unitization when \(C\) is nonunital, and set
\(\widetilde C=C\) when \(C\) is unital.  We use its canonical
action on \(C^{**}\) in products with support projections.
Given positive contractions \(a,y\in C\), with \(a\) rich,
and a contraction \(c\in\widetilde C\), we seek positive
contractions \(p_1\in\Her(a)\) and \(q_1\in\Her(y)\) with two
properties: each has represented norm one wherever its corresponding
input has norm greater than \(1/2\), and
\(\|e_{p_1}ce_{q_1}\|<1\).  The norm signals will later give
ideal comparison.

The construction has two parts.  A graph of orthogonal copies
first turns possible failure into finitely many explicit equalities.
Once those equalities have been excluded, ideal cutoffs convert
the resulting nonvanishing into one global support estimate.
We prove both parts here while leaving the finite exclusion
problem to the later arguments.  Throughout this section put
\[
 s=\frac14,\qquad \alpha_0=\frac1{32}.
\]
When the integer \(k\ge1\) has been chosen, fix
\(\gamma_g,\sigma_g>0\) with
\begin{equation}\label{eq:graph-constants}
 \gamma_g^2+k\sigma_g^2<1.
\end{equation}
For example, \(\gamma_g=1/2\) and
\(\sigma_g=1/(2\sqrt{k+1})\) satisfy this inequality.  Both the
number of coordinates and these coefficients are fixed before
the measurements below are computed.

\begin{proposition}[Fixed measurements of a graph]\label{prop:graph}
Let \(C\) be a \(C^*\)-algebra, let \(a,y\in C_+\) be
contractions with \(a\) rich, and let \(c\in\widetilde C\) be a
contraction.  For an integer \(k\ge1\) and fixed positive
\(\gamma_g,\sigma_g\) satisfying \eqref{eq:graph-constants},
there are
\[
 X\in\Her((a-1/8)_+)_+,\qquad
 B_1,\ldots,B_k\in\Her((a-1/8)_+)
\]
with the following property.  For any contractions
\[
 S_i\in\Her((X-\alpha_0)_+),\qquad 1\le i\le k,
\]
one can form a rich positive contraction \(p\in\Her(a)\) such
that, for every representation \(\pi\) of \(C\),
\(\|\pi(a)\|>1/2\) implies \(\|\pi(p)\|=1\).
If a representation \(\pi\) of \(C\) satisfies
\[
 \pi\bigl(y(1-c^*p^2c)y\bigr)=0,\qquad \|\pi(y)\|>s,
\]
then, in that same representation,
\begin{equation}\label{eq:graph-failure}
 \|\pi(X)\|>s,\qquad \pi(B_i)=\pi(S_i)\quad(1\le i\le k).
\end{equation}
The elements \(X,B_1,\ldots,B_k\) are fixed before the
\(S_i\) are chosen.
\end{proposition}

\begin{proof}
Let \(e=e_{(a-1/8)_+}\) and \(H=\Her((a-1/8)_+)\).
Since \(a\) is properly infinite,
\(a^{\oplus(k+2)}\precsim a\).  Apply
Lemma~\ref{lem:transport} with cutoff \(1/8\), and split the
resulting row into slots labelled \(0,1,\ldots,k,*\).  They
satisfy
\begin{equation}\label{eq:graph-slots}
 v_i^*v_j=\delta_{ij}e,\qquad v_i v_i^*\le e_a,\qquad
 v_iH\subset C.
\end{equation}
Choose a ramp \(t=t(a)\in H\) that is zero on \([0,1/4]\)
and one on \([1/2,1]\).  It is rich, and its represented norm
is one whenever the represented norm of \(a\) exceeds \(1/2\).

We now compute the measurements, before making any choice of
\(S_i\).  Put \(F=cyc^*\in C_+\).  The block calculation behind
their definition is as follows.  Suppose \(L\in C^{**}\) satisfies
\(L^*L=e\) and has slot coordinates
\(v_0^*L=\gamma_g e\), \(v_i^*L=\sigma_g S_i\), and suppose
\(F=LH_0L^*\) for a positive \(H_0\in eC^{**}e\).  Then
\[
 v_0^*Fv_0=\gamma_g^2H_0,\qquad
 v_i^*Fv_0=\sigma_g\gamma_g S_iH_0.
\]
Vanishing of the defect in a representation will force this form
there, with \(tH_0=H_0\) and with \(t\) acting identically on the
prescribed supports of
\(S_i\).  The first measurement therefore recovers \(H_0\), and
an inverse on its spectral cutoff recovers each \(S_i\) from the
second.  These are conditional identities used to design the
measurements; their verification follows after the graph is built.

Choose a continuous nonnegative function \(h_{\rm inv}\) on
\([0,\infty)\), zero
at zero, equal to \(\lambda^{-1}\) for \(\lambda\ge\alpha_0\),
and bounded by \(\alpha_0^{-1}\).  Define
\begin{equation}\label{eq:graph-measurements}
 X=\gamma_g^{-2}t v_0^*Fv_0t,\qquad
 B_i=(\sigma_g\gamma_g)^{-1}
       t v_i^*Fv_0t\,h_{\rm inv}(X).
\end{equation}
The sandwiches are \((v_it)^*F(v_0t)\), hence belong to \(C\).
They are supported on both sides by \(e\), so they lie in \(H\);
the same is true after multiplication by \(h_{\rm inv}(X)\).
Thus \(X\ge0\) and \(X,B_i\) have the asserted algebra
membership.  The slots, \(t,c,y\) determine them without using
any \(S_i\).

For the chosen \(S_i\), define in the corner \(eC^{**}e\)
\[
 J=\left((1-\gamma_g^2)e
       -\sigma_g^2\sum_{i=1}^kS_i^*S_i\right)^{1/2}.
\]
Then define in \(C^{**}\)
\[
 L=\gamma_g v_0+\sigma_g\sum_{i=1}^k v_iS_i+v_*J,
 \qquad p=LtL^*.
\]
Here \(S_i\in H\), since \(X\in H\).  The radicand is at least
\((1-\gamma_g^2-k\sigma_g^2)e\).  Its square root belongs to
\(H+\mathbb C e\) and multiplies \(H\) into itself.  Hence
\(Lh\in C\) for every \(h\in H\), term by term in
\eqref{eq:graph-slots}.  The slot identities also give
\[
 L^*L=e,\qquad e_aL=L.
\]
It follows, by the membership and support argument of
Lemma~\ref{lem:transport}, that \(h\mapsto LhL^*\) is an
injective \(*\)-homomorphism from \(H\) into \(\Her(a)\).
Indeed, for \(h\ge0\) its value is
\((Lh^{1/2})(Lh^{1/2})^*\in C\), supported under \(e_a\);
and \(L^*L=e\) gives multiplicativity and recovers \(h\) on
conjugating back.  The same identities give preservation of
cutoffs, their Cuntz classes, and represented norms, exactly
as in that lemma.  Thus \(p\) is a rich positive contraction
with the stated representation signal.  If \(e=0\), all the
operators on this corner are zero and these assertions still
hold.

It remains to see what vanishing of the defect forces.  Work
on the essential subspace of a representation in which
\[
 d=y(1-c^*p^2c)y=0,\qquad \|y\|>s,
\]
and omit the representation symbols.  The normal extensions
preserve the preceding support identities.  For every vector
\(\xi\),
\[
 \|y\xi\|=\|pc\,y\xi\|
       \le\|c\,y\xi\|\le\|y\xi\|.
\]
Equality first on \(\ran y\) and then on its closure shows
that \(c\) is isometric on \(e_y\mathcal H\).  For a positive
contraction \(p\), equality \(\|p\zeta\|=\|\zeta\|\) implies
\(p\zeta=\zeta\), by applying \(1-p^2\).  Consequently
\(ce_y\mathcal H\) lies in the \(p=1\) subspace, and the
positive operator \(F=cyc^*\) is supported on that subspace.
Since \(p=LtL^*\) and \(L^*L=e\), the \(p=1\) subspace lies
in the range of \(L\), and its pullback lies in the \(t=1\)
subspace.  With \(H_0=L^*FL\), we obtain
\[
 F=LH_0L^*,\qquad tH_0=H_0=H_0t,\qquad
 \|H_0\|=\|F\|=\|y\|>s.
\]
The last equality uses the isometric action of \(c\) on
\(e_y\mathcal H\).

Now \(v_0^*L=\gamma_g e\), \(v_i^*L=\sigma_g S_i\), and
\(L^*v_0=\gamma_g e\).  Substituting these identities into
\eqref{eq:graph-measurements} yields
\[
 X=tH_0t=H_0,\qquad
 B_i=tS_iH_0t\,h_{\rm inv}(X).
\]
Since \(tX=X\), \(t\) is the identity on
\(\overline{\ran X}=e_X\mathcal H\).  Each represented \(S_i\)
has both supports under \(e_{(X-\alpha_0)_+}\).  Therefore
\(tS_i=S_i\) and \(S_iXh_{\rm inv}(X)=S_i\).  The last display
gives \(B_i=S_i\), while \(\|X\|=\|H_0\|>s\).
This proves \eqref{eq:graph-failure}.
\end{proof}

Thus it is enough to choose the \(S_i\) in their prescribed
spectral corner so that the simultaneous relations
\eqref{eq:graph-failure} cannot hold.  The graph then supplies
the nonvanishing implication assumed in the next lemma.  The
lemma uses the copied element only through its richness and
its representation signal, so it is independent of how the
finite equalities were excluded.

\begin{lemma}[From nonvanishing to shrinking]\label{lem:conditional-shrinking}
Let \(C\) be a \(C^*\)-algebra, let \(a,y\in C_+\) be
contractions with \(a\) rich, and let \(c\in\widetilde C\) be
a contraction.  Suppose \(p\in\Her(a)\) is a rich positive
contraction satisfying
\[
 \|\pi(a)\|>1/2\quad\Longrightarrow\quad\|\pi(p)\|=1
 \qquad\text{for every representation }\pi\text{ of }C.
\]
Put \(d=y(1-c^*p^2c)y\in C_+\), and assume
\begin{equation}\label{eq:common-nonvanishing}
 \pi(d)=0\quad\Longrightarrow\quad\|\pi(y)\|\le\frac14
 \qquad\text{for every representation }\pi\text{ of }C.
\end{equation}
Then there are positive contractions
\[
 p_1\in\Her(a),\qquad q_1\in\Her(y)
\]
such that \(p_1\) is rich,
\begin{equation}\label{eq:shrinking-signals}
 \begin{split}
  \|\pi(a)\|>1/2&\ \Longrightarrow\ \|\pi(p_1)\|=1,\\
  \|\pi(y)\|>1/2&\ \Longrightarrow\ \|\pi(q_1)\|=1
 \end{split}
 \qquad\text{for every representation }\pi\text{ of }C,
\end{equation}
and
\begin{equation}\label{eq:shrinking-support}
 \|e_{p_1}c e_{q_1}\|_{C^{**}}<1.
\end{equation}
The last inequality is a single norm bound in the universal
bidual.
\end{lemma}

\begin{proof}
The contraction assumptions give \(0\le d\le y^2\le1\).
By \eqref{eq:common-nonvanishing}, every irreducible
representation vanishing on \(d\) vanishes on \((y-s)_+\).
The primitive ideal test in Lemma~\ref{lem:ideals-cutoffs}(ii)
therefore gives \((y-s)_+\in\Id(d)\).  Set \(s'=3/8\).
Apply cutoff exhaustion in Lemma~\ref{lem:ideals-cutoffs}(i)
to \((y-s)_+\) with cutoff \(s'-s\).  After decreasing the
resulting parameter if necessary, there is \(0<\delta<1\) with
\begin{equation}\label{eq:shrinking-first-ideal}
 (y-s')_+\in\Id(h_d),\qquad h_d=(d-\delta)_+.
\end{equation}
Define
\[
 z=yh_d^{1/2}\in C,\qquad q=zz^*=y h_d y\in\Her(y).
\]
The order \(d\le y^2\le1\) gives
\begin{equation}\label{eq:shrinking-ideal-order}
 \delta h_d
 \le h_d^{1/2}d h_d^{1/2}
 \le z^*z=h_d^{1/2}y^2h_d^{1/2}
 \le h_d.
\end{equation}
Hence \(z^*z\) and \(h_d\) have the same hereditary algebra,
and \(z^*z\CuEq zz^*=q\).  It follows that
\(\Id(q)=\Id(h_d)\).  A second application of cutoff
exhaustion, to \((y-s')_+\) with cutoff \(1/2-s'\), gives
\(\delta'>0\) such that
\begin{equation}\label{eq:shrinking-second-ideal}
 (y-1/2)_+\in\Id((q-\delta')_+).
\end{equation}
Let \(q_1\) be a ramp of \(q\), zero on \([0,\delta'/2]\)
and one on \([\delta',\infty)\).  It is a positive contraction
in \(\Her(q)\subset\Her(y)\).  If \(\|\pi(y)\|>1/2\),
then \(\pi((y-1/2)_+)\ne0\).  The ideal inclusion
\eqref{eq:shrinking-second-ideal} forces
\(\pi((q-\delta')_+)\ne0\), so \(\|\pi(q)\|>\delta'\)
and \(\|\pi(q_1)\|=1\).

To bound the universal-bidual norm, work in the universal
nondegenerate representation.  For
\(\xi\in e_{h_d}\mathcal H
=\mathbf1_{(\delta,\infty)}(d)\mathcal H\),
\begin{align}
 \|pc\,y\xi\|^2
 &=\|y\xi\|^2-\langle d\xi,\xi\rangle\notag\\
 &\le\|y\xi\|^2-\delta\|\xi\|^2
 \le(1-\delta)\|y\xi\|^2.
 \label{eq:shrinking-norm-loss}
\end{align}
The last inequality uses that \(y\) is a contraction.
Apply this estimate to vectors in \(\ran(h_d^{1/2})\).
The closure of \(\ran(yh_d^{1/2})\) is \(e_q\mathcal H\),
so continuity gives
\[
 \|pc e_q\|\le\sqrt{1-\delta}.
\]
Choose \(\sqrt{1-\delta}<\tau<1\), and take \(p_1\) to be a
ramp of \(p\), zero on \([0,\tau]\) and one at \(1\).
Functional calculus gives \(e_{p_1}\le\tau^{-2}p^2\),
while \(e_{q_1}\le e_q\).  Consequently
\begin{equation}\label{eq:shrinking-final-bound}
 \|e_{p_1}c e_{q_1}\|
 \le\tau^{-1}\|pc e_{q_1}\|
 \le\frac{\sqrt{1-\delta}}{\tau}<1.
\end{equation}
The ramp of the rich \(p\) is rich and remains in \(\Her(a)\).
Spectral mapping at its plateau shows that it has norm one in
every representation where \(p\) has norm one.  This proves
the first signal in \eqref{eq:shrinking-signals}; the second
was proved above.  The proof also covers zero inputs.
\end{proof}

\section{Four entries under individual proper infiniteness}
\label{sec:four-variables}

Assume in this section that \(C\) satisfies \(\propP\).
We choose four entries that rule out the fixed measurement relations
in Proposition~\ref{prop:graph}.  This will supply the nonvanishing
hypothesis of Lemma~\ref{lem:conditional-shrinking} and hence a
strict contraction between hereditary supports.  A spectral cutoff
is essential: it places a support in an actual operator range, where
the following positivity argument applies.

\begin{lemma}[A spectral cutoff in an actual range]\label{lem:actual-range}
Let \(x,Q\) be bounded positive operators on a Hilbert space
\(\mathcal H\), and put \(U=xQ^2x\).  For every \(\eta>0\),
\begin{equation}\label{eq:actual-range}
 e_{(U-\eta)_+}\mathcal H\subset\ran(xQ),
 \qquad
 \ker Q\cap\ran(xQ)=\{0\}.
\end{equation}
\end{lemma}

\begin{proof}
Let \(f_\eta\) be the bounded Borel function equal to
\(\lambda^{-1}\) for \(\lambda>\eta\) and zero otherwise.
Its norm on the spectrum is at most \(\eta^{-1}\).  If
\(\zeta\in e_{(U-\eta)_+}\mathcal H\), then
\[
 \zeta=Uf_\eta(U)\zeta
      =xQ\bigl(Qx f_\eta(U)\zeta\bigr),
\]
which proves the actual-range inclusion.  If
\(\zeta=xQ\xi\) and \(Q\zeta=0\), then
\[
 0=\langle QxQ\xi,\xi\rangle=\|x^{1/2}Q\xi\|^2,
\]
so \(\zeta=xQ\xi=0\).  The assertion concerns the actual range;
no closed-range conclusion is being used.
\end{proof}

\begin{proposition}[Four entries avoid the fixed measurements]
\label{prop:four-avoidance}
Suppose \(C\) satisfies \(\propP\).  Let \(X\in C_+\) and
\(B_1,B_2,B_3,B_4\in C\).  With
\[
 \alpha_0=\frac1{32}<\alpha=\frac1{16}
 <\beta=\frac18<s=\frac14,\qquad h=\frac12,
\]
there exist contractions \(T_1,\ldots,T_4\) in
\(\Her((X-\alpha_0)_+)\) for which no representation \(\pi\)
of \(C\) satisfies
\begin{equation}\label{eq:four-forbidden}
 \|\pi(X)\|>s,\qquad \pi(B_i)=\pi(T_i)\quad(1\le i\le4).
\end{equation}
\end{proposition}

\begin{proof}
We argue by contradiction.  If every four-entry choice could be
matched in a representation, the fixed elements could be passed to
a norm ultraproduct.  In an irreducible representation of the
resulting image of \(C\), the first pair would give an isometry with
nonzero defect, while the last pair would rule out that defect.

All formulas described as fixed below depend only on \(X,B_i\)
and the displayed constants.  Define
\begin{equation}\label{eq:four-first-fixed}
 x=(X-\alpha)_+,\qquad V=(X-\beta)_+,\qquad
 Q=(h1-B_1B_1^*)_+,\qquad U=xQ^2x.
\end{equation}
Here \(Q\in\widetilde C_+\), whereas \(x,V,U\in C_+\).
Let \(\fin(C)\) denote the set of finite subsets of \(C\), and order
\[
 \Lambda=\mathbb N_{\ge1}\times\fin(C),
\]
by the usual order in the first coordinate and inclusion in the
second.  At \(\lambda=(m,E)\), put \(\eta=1/m\), so
\(\eta\to0\) along this directed set, and every fixed \(w\in C\)
eventually belongs to \(E\).

\medskip\noindent\emph{The first two entries.}
The first pair will force \(UwV=0\) for every fixed \(w\in C\)
in a limiting image of \(C\).  At the index \((m,E)\), we test
only the finitely many \(w\in E\).  A positive sum lets one norm
bound control all of these tests: set
\begin{equation}\label{eq:four-D}
 A_\eta=(U-\eta)_+,\qquad
 D_0=A_\eta\left(\sum_{w\in E}wVw^*\right)A_\eta,\qquad
 D=g_\eta(D_0),
\end{equation}
where \(g_\eta\) is a ramp that is zero on \([0,\eta]\) and
one on \([2\eta,\infty)\).  In fact,
\[
 D_0=\sum_{w\in E}
 (A_\eta wV^{1/2})(A_\eta wV^{1/2})^*.
\]
Thus a small norm for \(D_0\) makes each
\(A_\eta wV^{1/2}\) small.  We have
\[
 D_0\in\Her(A_\eta)\cap\Id(V),\qquad
 e_D\le e_{(D_0-\eta)_+}\le e_{A_\eta}\le e_x.
\]
Since \(V\) is properly infinite, ideal comparison gives
\(D_0\oplus D_0\precsim V\).  Apply
Lemma~\ref{lem:transport} with cutoff \(\eta\), split the
carrier into two columns, and multiply these columns by
\(D^{1/2}\).  Taking adjoints of the resulting two elements gives
contractions \(T_1,T_2\in C\) with
\begin{equation}\label{eq:four-first-entries}
 T_iT_j^*=\delta_{ij}D,\qquad e_VT_i^*=T_i^*
 \quad(i,j=1,2).
\end{equation}
Their left supports lie under \(e_D\) and their right supports
under \(e_V\).  Both are under \(e_x\le e_{(X-\alpha_0)_+}\);
they therefore belong to the required hereditary algebra.

Consider a representation satisfying
\eqref{eq:four-forbidden} for the entries at this index, and
omit its symbols.  Then \(B_1B_1^*=D\) and
\(Q=(h1-D)_+\).  The spectral subspace of \(D\) for
\((h,\infty)\) is killed by \(Q\) and lies in
\(e_D\mathcal H\subset e_{(U-\eta)_+}\mathcal H\).
Lemma~\ref{lem:actual-range} shows that this subspace is zero.
Thus
\begin{equation}\label{eq:four-first-bounds}
 B_1B_1^*=D\le h1,\qquad \|D_0\|\le2\eta
 \quad\text{in this representation}.
\end{equation}
For the second conclusion, if the norm of \(D_0\) were at
least \(2\eta\), spectral mapping and the plateau of \(g_\eta\)
would give \(\|D\|=1>h\).

\medskip\noindent\emph{A fixed prospective defect and the last two entries.}
Choose ramps
\[
 H=h_0(X),\qquad K=k_0(X)
\]
so that \(h_0\) vanishes on a neighborhood of
\([0,\alpha]\) and equals one on \([\beta,\infty)\), and
\(k_0\) vanishes on a neighborhood of \([0,\alpha_0]\) and
equals one on \([\alpha,\infty)\).  In particular \(KH=H\).
Set
\begin{equation}\label{eq:four-r}
 r_i=h^{-1/2}B_i^*H\quad(i=1,2),\qquad
 b_0=(1-H^2)^{1/2}.
\end{equation}
The reason for these definitions is the following conditional
calculation.  In the limiting irreducible representation below, the
first pair will give
\[
 r_i^*r_j=\delta_{ij}H^2,\qquad Hr_i=r_i\quad(i,j=1,2),
 \qquad H\ne0.
\]
These relations are not being asserted in \(C\) at this stage.
They imply \(b_0r_i=0\), and hence
\[
 v=r_1+b_0,\qquad v^*v=H^2+b_0^2=1,\qquad v^*r_2=0.
\]
Since \(r_2^*r_2=H^2\ne0\), the isometry \(v\) has a nonzero
defect projection \(1-vv^*\).  This is the scaling-element
construction of Blackadar--Cuntz \cite{BlackadarCuntz1982}, in the
form recalled in \cite[Remark~2.4]{PasnicuRordam2007}; the second
entry supplies the additional nonzero orthogonal range.

To choose the last two entries before taking any limit, we encode
this prospective defect by one fixed element of \(C\).  Compression
by \(K\) puts the formula in the required hereditary algebra, and
clipping makes it a positive contraction even before the preceding
relations hold.  Let
\(\chi(t)=\min\{1,\max\{0,t\}\}\) on \(\mathbb R\), and define
\begin{equation}\label{eq:four-P0}
 P_0=\chi\!\left(
 K\bigl[1-(r_1+b_0)(r_1+b_0)^*\bigr]K
 \right)\in C.
\end{equation}
The compressed selfadjoint element is in \(C\) and supported
under \(e_K\); since \(\chi(0)=0\), continuous functional
calculus in \(\Her(K)\) gives
\[
 e_{P_0}\le e_K\le e_{(X-\alpha_0)_+}.
\]
This membership and support statement uses no relation from a
representation.

Define the following additional fixed elements:
\begin{equation}\label{eq:four-second-fixed}
 x'=(P_0-1/2)_+,\qquad
 Q'=(h1-B_3B_3^*)_+,\qquad R=x'(Q')^2x'.
\end{equation}
At the same index \(\eta\), property \(\propP\) gives
\[
 ((R-\eta)_+)^{\oplus2}\precsim(R-\eta)_+.
\]
Apply Lemma~\ref{lem:transport} to this comparison with cutoff
\(\eta\).  If \(t_R\) is a ramp of \(R\),
zero on \([0,2\eta]\) and one on \([4\eta,\infty)\), the two
carrier columns times \(t_R\) give contractions \(T_3,T_4\in C\)
such that
\begin{equation}\label{eq:four-last-entries}
 T_3^*T_4=0,\qquad T_3^*T_3=T_4^*T_4=t_R^2.
\end{equation}
Their left supports lie under \(e_{(R-\eta)_+}\) and their
right supports under \(e_{(R-2\eta)_+}\).  Both are under
\(e_{x'}\le e_{(X-\alpha_0)_+}\).  Transport preserves the
represented norms, so
\begin{equation}\label{eq:four-last-norm}
 \|\pi(R)\|>4\eta\quad\Longrightarrow\quad
 \|\pi(T_3)\|=1.
\end{equation}

In a representation satisfying \eqref{eq:four-forbidden},
\(B_3=T_3\).  Apply Lemma~\ref{lem:actual-range} to \(x',Q',R\).
The spectral subspace of \(B_3B_3^*\) above \(h\) is killed by
\(Q'\) and lies under \(e_{(R-\eta)_+}\), so it is zero.
Equation~\eqref{eq:four-last-norm} then gives
\begin{equation}\label{eq:four-last-bounds}
 B_3B_3^*\le h1,\qquad \|R\|\le4\eta
 \quad\text{in this representation}.
\end{equation}

\medskip\noindent\emph{Passing only fixed elements to an ultraproduct.}
Suppose, for a contradiction, that every index in \(\Lambda\)
admits a representation satisfying \eqref{eq:four-forbidden}.
Choose one, denoted \(\pi_\lambda\), on its nondegenerate
essential subspace.  There is no size issue in this choice: from
any such representation, choose a vector detecting
\(\|\pi_\lambda(X)\|>s\) and restrict to its reducing cyclic
subrepresentation.  A set of cyclic representations suffices.
Choose an ultrafilter \(\omega\) containing all tails of
\(\Lambda\), and form the norm ultraproduct
\[
 \mathcal B=\prod_{\omega}B(\mathcal H_\lambda)
 =\frac{\prod_\lambda B(\mathcal H_\lambda)}
 {\{(z_\lambda):\lim_{\omega}\|z_\lambda\|=0\}}.
\]
The norm of a class is \(\lim_{\omega}\|z_\lambda\|\).
Let \(\Phi:C\to\mathcal B\) be
\(\Phi(z)=[(\pi_\lambda(z))]\), and write \(B=\Phi(C)\).
The image is a closed \(C^*\)-subalgebra.  The unital extensions
carry the unitized formulas into \(C^*(B,1_{\mathcal B})\).
We suppress \(\Phi\) on fixed elements from now on.

For any fixed \(w\in C\), eventually \(w\in E\).  Positivity
of the summands in \eqref{eq:four-D} and
\eqref{eq:four-first-bounds} yield
\[
 \|\pi_\lambda((U-\eta)_+wV^{1/2})\|^2\le2\eta
\]
on that tail.  Since \(\|U-(U-\eta)_+\|\le\eta\),
\begin{equation}\label{eq:four-UCV-estimate}
 \|\pi_\lambda(UwV)\|
 \le\sqrt{2\eta}\,\|V\|^{1/2}
      +\eta\|w\|\|V\|\longrightarrow0.
\end{equation}
It follows that \(UBV=0\) in the image algebra.  The other exact
relations and bounds at each index give
\begin{equation}\label{eq:four-ultraproduct}
 \begin{gathered}
 \|X\|\ge s,\qquad B_1B_1^*\le h1,\qquad
 B_iB_j^*=\delta_{ij}B_1B_1^*,\quad HB_i^*=B_i^*
       \quad(i,j=1,2),\\
 UBV=0,\qquad R=0,\qquad B_3B_3^*\le h1,\qquad
 B_3^*B_4=0,\quad B_3^*B_3=B_4^*B_4.
 \end{gathered}
\end{equation}
For \(HB_i^*=B_i^*\), use \(e_VT_i^*=T_i^*\) and
\(H=1\) on the support of \(V\).
Fix also a continuous \(h_P:[0,1]\to[0,1]\) with
\(h_P(0)=0\) and \(h_P=1\) on \([1/2,1]\).  The left
support of \(T_4\) lies under \(e_{x'}\), so we obtain
\begin{equation}\label{eq:four-B4-support}
 h_P(P_0)B_4=B_4.
\end{equation}
Every element in \eqref{eq:four-ultraproduct} and
\eqref{eq:four-B4-support} is the image of a fixed algebra
element or of its fixed unitized functional calculus.  We have
not placed the varying \(T_i,D,D_0,t_R\) in \(B\), and we do
not require proper infiniteness of \(B\) or of \(\mathcal B\).
The directed finite sets in \eqref{eq:four-UCV-estimate} handle
all \(w\in C\) without a separability assumption.

\medskip\noindent\emph{A nonzero defect gives the contradiction.}
Since \(\|X\|\ge s>\beta\), the image \(V=(X-\beta)_+\)
is nonzero.  Choose an irreducible representation of \(B\)
in which \(V\ne0\).  Extend this nondegenerate representation
to \(C^*(B,1_{\mathcal B})=B+\mathbb C1_{\mathcal B}\) by
sending \(1_{\mathcal B}\) to the identity operator; this remains
valid if \(B\) has a unit different from \(1_{\mathcal B}\).
Thus the unitized expressions \(Q,b_0,Q'\) and the defect below
have their stated represented meanings.  Again suppress the
representation symbols.  The closed span of \(BV\mathcal H\) is a nonzero
reducing subspace, and hence all of \(\mathcal H\).
The relation \(UBV=0\) therefore gives \(U=0\).
Put \(D_1=B_1B_1^*\).  Since \(D_1\le h1\), we have
\(Q=h1-D_1\), and
\[
 0=U=(Qx)^*(Qx)
 \quad\Longrightarrow\quad D_1x=hx.
\]
By closure \(D_1=h1\) on \(e_x\mathcal H\).  The support of
\(H\) lies there.  Equations~\eqref{eq:four-r} and
\eqref{eq:four-ultraproduct} give
\begin{equation}\label{eq:four-isometry-relations}
 r_i^*r_j=\delta_{ij}H^2,\qquad Hr_i=r_i
 \quad(i,j=1,2).
\end{equation}
Also \(HV=V\ne0\), so \(H\ne0\) and \(r_2\ne0\).

The conditional calculation preceding \eqref{eq:four-P0} now
applies: \(v=r_1+b_0\) is an isometry with \(v^*r_2=0\).
Its defect projection
\[
 P=1-vv^*
\]
is nonzero, since \(Pr_2=r_2\ne0\).

We check that this projection is the represented fixed element
\(P_0\), a point for which both sides of the cutoff matter.
The scalar ramps give \(KH=H\).  From \(Hr_1=r_1\) follows
\(Kr_1=r_1\), and the definition \(r_1=h^{-1/2}B_1^*H\)
gives \(r_1K=r_1\).  Also \(K\) commutes with \(b_0\).
Expanding the defect gives
\[
 P=H^2-r_1r_1^*-r_1b_0-b_0r_1^*.
\]
Each term is fixed by multiplication by \(K\) on both sides.
Therefore \(KPK=P\), and \(\chi(P)=P\) because \(P\) is a
projection.  Formula~\eqref{eq:four-P0} yields
\begin{equation}\label{eq:four-P0-is-P}
 P_0=P.
\end{equation}

Now \(x'=(P_0-1/2)_+=P/2\).  Put \(D_3=B_3B_3^*\le h1\);
then \(Q'=h1-D_3\).  The relation \(R=0\) gives
\(Q'x'=0\), so
\begin{equation}\label{eq:four-last-contradiction}
 D_3P=hP.
\end{equation}
Equations~\eqref{eq:four-B4-support} and
\eqref{eq:four-P0-is-P} say \(B_4=PB_4\).
But \(B_3^*B_4=0\), so \(D_3B_4=0\), whereas
\eqref{eq:four-last-contradiction} gives \(D_3B_4=hB_4\).
Thus \(B_4=0\), and
\(B_3^*B_3=B_4^*B_4\) gives \(B_3=0\).
This contradicts \eqref{eq:four-last-contradiction}, since
\(h>0\) and \(P\ne0\).

The contradiction shows that some index has no representation
satisfying \eqref{eq:four-forbidden}.  The four entries chosen
at that index prove the proposition.
\end{proof}

We now apply the avoidance proposition to the graph measurements.
This gives the form of shrinking that will be used throughout the
next section.

\begin{proposition}[Shrinking under individual proper infiniteness]
\label{prop:shrinking-P}
Suppose \(C\) satisfies \(\propP\).  Let \(a,y\in C\) be positive
contractions and let \(c\in\widetilde C\) be a contraction.  There
are positive contractions \(p_1\in\Her(a)\) and
\(q_1\in\Her(y)\), with \(p_1\) rich, such that for every
representation \(\pi\) of \(C\),
\begin{equation}\label{eq:shrinking-P-signals}
 \|\pi(a)\|>1/2\ \Longrightarrow\ \|\pi(p_1)\|=1,
 \qquad
 \|\pi(y)\|>1/2\ \Longrightarrow\ \|\pi(q_1)\|=1,
\end{equation}
and, in the universal bidual,
\begin{equation}\label{eq:shrinking-P-support}
 \|e_{p_1}c e_{q_1}\|<1.
\end{equation}
\end{proposition}

\begin{proof}
Every positive contraction of \(C\) is rich under \(\propP\), since
each of its scalar cutoffs is properly infinite.  In
Proposition~\ref{prop:graph}, fix \(k=4\) before the measurements
are computed, and choose, for example,
\(\gamma_g=1/2\) and \(\sigma_g=1/(2\sqrt5)\).  Then
\(\gamma_g^2+4\sigma_g^2<1\).  The graph proposition gives fixed
\(X\) and \(B_1,\ldots,B_4\) for the inputs \(a,y,c\).
Apply Proposition~\ref{prop:four-avoidance} to these elements and
use its contractions \(T_i\) as the graph choices \(S_i\).
The resulting rich positive contraction \(p\in\Her(a)\) has the
required norm-one signal for \(a\).

Put \(d=y(1-c^*p^2c)y\).  If a representation \(\pi\) had
\(\pi(d)=0\) and \(\|\pi(y)\|>1/4\), the graph implication would
give \(\|\pi(X)\|>1/4\) and \(\pi(B_i)=\pi(T_i)\) for every
\(i\).  These are exactly the relations excluded by
Proposition~\ref{prop:four-avoidance}.  Thus
\begin{equation}\label{eq:nonvanishing-P}
 \pi\bigl(y(1-c^*p^2c)y\bigr)=0
 \quad\Longrightarrow\quad \|\pi(y)\|\le1/4
\end{equation}
for every representation.  The same \(p\) is rich and carries the
\(a\)-signal, so Lemma~\ref{lem:conditional-shrinking} applies to
\(a,y,c,p\) and proves the stated conclusions.  This argument also
covers zero inputs.
\end{proof}

\section{Uniform separation and general diagonalization}
\label{sec:separation}

We now prove Theorem~\ref{thm:individual-to-strong}.  The strict
support bound in Proposition~\ref{prop:shrinking-P} initially
depends on the inputs.  Products turn it into a universal constant
below one; iteration then makes the support compression arbitrarily
small.  A final cutoff comparison gives coefficients whose norms
are bounded independently of that requested smallness.

\begin{lemma}[Products preserve individual proper infiniteness]
\label{lem:product-P}
Let \((C_i)_{i\in I}\) be any family of \(C^*\)-algebras
satisfying \(\propP\).  Their bounded product
\(D=\prod_{i\in I}C_i\) satisfies \(\propP\).
More precisely, for a positive contraction \(d=(d_i)\in D\)
and \(0<\delta<1/2\), there is a row
\(s=(s_i)\in M_{1,2}(D)\) such that
\begin{equation}\label{eq:product-row}
 s^*ds=(d-2\delta)_+\oplus(d-2\delta)_+,\qquad
 \|s\|\le\delta^{-1/2}.
\end{equation}
\end{lemma}

\begin{proof}
For an empty index set, \(D=0\) and the zero row has the asserted
properties.  Otherwise, for each \(i\), put
\[
 h_i=(d_i-\delta)_+,\qquad
 k_i=(d_i-2\delta)_+\oplus(d_i-2\delta)_+.
\]
The positive \(h_i\) is properly infinite.  Apply
Lemma~\ref{lem:transport} to \(h_i\oplus h_i\precsim h_i\)
with cutoff \(\delta\).  If \(v_i\) is the resulting
rectangular partial isometry, then
\[
 r_i=v_i k_i^{1/2}\in M_{1,2}(C_i),\qquad
 r_i^*r_i=k_i,\qquad e_{h_i}r_i=r_i,\qquad \|r_i\|\le1.
\]
The algebra membership follows from transport regularity; no
membership of \(v_i\) itself is needed.  Define
\[
 \phi_\delta(t)=
 \begin{cases}
  t/\delta^{3/2},&0\le t\le\delta,\\
  t^{-1/2},&t\ge\delta.
 \end{cases}
\]
This continuous function vanishes at zero, is bounded by
\(\delta^{-1/2}\), and satisfies
\(t^{1/2}\phi_\delta(t)=1\) for \(t>\delta\).  Consequently
\[
 s_i=\phi_\delta(d_i)r_i\in M_{1,2}(C_i),\qquad
 d_i^{1/2}s_i=r_i,\qquad
 \|s_i\|\le\delta^{-1/2}.
\]
The uniform bound places \((s_i)\) in \(M_{1,2}(D)\).
Coordinate functional calculus gives
\eqref{eq:product-row}.  Its right-hand side converges to
\(d\oplus d\) in norm as \(\delta\downarrow0\), proving
proper infiniteness of \(d\).  Scaling covers every positive
element of \(D\).
\end{proof}

The general product-comparison criterion of
Antoine--Perera--Thiel \cite[Lemma~5.7]{AntoinePereraThiel2020}
also implies this permanence result: for each cutoff
\(\varepsilon\), one can take a target cutoff
\(\delta=\varepsilon/2\) and use proper infiniteness in every
coordinate.  The proof above records the particular bounded
rows that make this inference possible.  Coordinatewise
comparison without a common cutoff bound would not suffice.

\begin{lemma}[A uniform support contraction]\label{lem:uniform-shrinking}
There exists \(\theta\in(0,1)\) such that for every
\(C^*\)-algebra \(C\) satisfying \(\propP\), every pair of positive
contractions \(a,y\in C\), and every contraction \(c\in C\),
there are positive contractions \(p\in\Her(a)\),
\(q\in\Her(y)\) such that, for every representation \(\pi\) of
\(C\),
\[
 \|\pi(a)\|>1/2\ \Longrightarrow\ \|\pi(p)\|=1,
 \qquad
 \|\pi(y)\|>1/2\ \Longrightarrow\ \|\pi(q)\|=1,
\]
and
\begin{equation}\label{eq:uniform-support}
 \|e_p c e_q\|\le\theta.
\end{equation}
\end{lemma}

\begin{proof}
If no such bound existed, for each \(m\ge2\) we could choose
an algebra \(C_m\) satisfying \(\propP\) and contractions
\(a_m,y_m,c_m\), with the first two positive, such that every
pair with the stated membership and simultaneous signals has
support norm greater than \(1-1/m\).  Form their bounded
product \(D\) and the three elements \(a=(a_m)\),
\(y=(y_m)\), \(c=(c_m)\in D\).  Lemma~\ref{lem:product-P}
gives \(\propP\) for \(D\), so
Proposition~\ref{prop:shrinking-P} produces \(p,q\) with the
signals and \(\kappa=\|e_pce_q\|<1\).

For each \(m\), the coordinate homomorphism \(D\to C_m\)
preserves hereditary membership, and its normal extension sends
the two support projections to the coordinate supports.
Thus the coordinate support norm is at most \(\kappa\).
If \(\pi\) is any representation of \(C_m\), its composition
with the coordinate homomorphism is a representation of \(D\),
so both signals hold for the coordinate pair as well.  Choosing
\(m\) with \(1-1/m>\kappa\) contradicts its selection.
The resulting constant can be increased, if necessary, to lie
in \((0,1)\).
\end{proof}

This proof uses the normal coordinate maps, not an identification
of \(D^{**}\) with a product of biduals.  The input \(c\) in
Lemma~\ref{lem:uniform-shrinking} belongs to \(C\).
This ensures that the sequence \((c_m)\) belongs to \(D\);
a sequence of arbitrary individual unitizations need not belong
to \(D\) or its unitization.  We need the uniform statement
only in the stated form.

\begin{proposition}[Arbitrarily small support compression]
\label{prop:iterated-separation}
Suppose \(C\) satisfies \(\propP\).  For positive contractions
\(a_*,b_*\in C\), a contraction \(c_*\in C\), and
\(\zeta>0\), there are positive contractions
\(p\in\Her(a_*)\), \(q\in\Her(b_*)\) such that
\[
 \|e_p c_*e_q\|<\zeta
\]
and, for every representation \(\pi\),
\[
 \|\pi(a_*)\|>1/2\ \Longrightarrow\ \|\pi(p)\|=1,
 \qquad
 \|\pi(b_*)\|>1/2\ \Longrightarrow\ \|\pi(q)\|=1.
\]
\end{proposition}

\begin{proof}
Fix \(\theta\) from Lemma~\ref{lem:uniform-shrinking}.
On \([0,1]\) set
\[
 f(t)=(2t-1)_+,\qquad g(t)=\min\{1,(4t-1)_+\}.
\]
Both vanish near zero, \(f(1)=1\), and \(g=1\) on
\([1/2,1]\).  For a positive contraction \(h\),
\begin{equation}\label{eq:iteration-cutoffs}
 g(h)=e_hg(h),\qquad g(h)e_{f(h)}=e_{f(h)}.
\end{equation}
Apply Lemma~\ref{lem:uniform-shrinking} once to
\(a_*,b_*,c_*\), obtaining \(p_1,q_1\) with support norm
at most \(\theta\).  Suppose inductively that
\(\|e_{p_k}c_*e_{q_k}\|\le\theta^k\).  The next normalized
input is
\begin{equation}\label{eq:iteration-normalization}
 c^{(k)}=\theta^{-k}g(p_k)c_*g(q_k)\in C,\qquad
 \|c^{(k)}\|\le1.
\end{equation}
Apply the uniform lemma to \(f(p_k),f(q_k),c^{(k)}\).
The new supports lie under those of \(f(p_k)\) and \(f(q_k)\);
by \eqref{eq:iteration-cutoffs},
\[
 e_{p_{k+1}}c_*e_{q_{k+1}}
 =\theta^k e_{p_{k+1}}c^{(k)}e_{q_{k+1}},
 \qquad
 \|e_{p_{k+1}}c_*e_{q_{k+1}}\|\le\theta^{k+1}.
\]
The hereditary algebras are nested.  In every representation,
norm one of \(p_k\) implies norm one of \(f(p_k)\), by
spectral mapping and \(f(1)=1\), and similarly for \(q_k\).
Thus the two original signals persist at every step.
All positive contractions needed as first inputs are rich under \(\propP\).
Choose \(k\) with \(\theta^k<\zeta\).
\end{proof}

\begin{proposition}[Comparison with bounded coefficients]
\label{prop:bounded-diagonalization}
Suppose \(C\) satisfies \(\propP\).  Let \(a,b\in C_+\) be
contractions and let \(c\in C\).  Given
\[
 0<\delta<1/2,\qquad 0<\rho\le2\delta,\qquad \zeta>0,
\]
there exist \(s,t\in C\) such that
\begin{align}
 \|s^*as-a\|,\ \|t^*bt-b\|&<2\delta+\rho,
 \label{eq:bounded-diagonals}\\
 \|s^*ct\|&<\zeta,\qquad
 \|s\|,\ \|t\|\le\delta^{-1/2}.
 \label{eq:bounded-cross}
\end{align}
\end{proposition}

\begin{proof}
Let \(a_*,b_*\) be ramps of \(a,b\), zero on
\([0,\delta]\) and one on \([2\delta,1]\).
Put \(M_c=\max\{1,\|c\|\}\).  Choose a finite iteration in
Proposition~\ref{prop:iterated-separation} so that its outputs
satisfy the signals and
\begin{equation}\label{eq:bounded-chosen-separation}
 \|e_p c e_q\|<\delta\zeta.
\end{equation}
For example, apply the construction to \(c/M_c\) and take
\(k\) with \(M_c\theta^k<\delta\zeta\).
This choice precedes the comparison coefficients below.

It remains to realize the diagonal cutoffs through \(p\) and \(q\)
with a coefficient bound depending only on \(\delta\) and \(\rho\).

Put \(d_a=(a-2\delta)_+\).  If an irreducible representation
detects \(d_a\), its norm on \(a\) is greater than \(2\delta\).
It then has norm one on \(a_*\), and the signal gives norm
one on \(p\).  Lemma~\ref{lem:ideals-cutoffs}(ii) yields
\(d_a\in\Id(p)\).  Since
\[
 e_p\le e_{a_*}\le\mathbf1_{(\delta,\infty)}(a),
 \qquad pap\ge\delta p^2,
\]
we have \(\Id(pap)=\Id(p)\): one inclusion is immediate, and
the displayed order gives the other.  The positive \(pap\)
is properly infinite under \(\propP\), so
Lemma~\ref{lem:comparison-rules}(iii) gives
\(d_a\precsim pap\).

Compressing a stabilization comparison coefficient to the scalar
source and target corners supplies \(w\in C\) with
\[
 \|w^*papw-d_a\|<\rho.
\]
Set \(s=pw\).  Then \(e_ps=s\), \(s^*as=w^*papw\), and
\begin{equation}\label{eq:final-coefficient-bound}
 \delta s^*s\le s^*as,\qquad
 \|s\|^2
 \le\frac{\|d_a\|+\rho}{\delta}
 \le\frac{1-2\delta+\rho}{\delta}
 \le\frac1\delta.
\end{equation}
In particular the more general estimate
\(\|pw\|^2\le(1+\rho)/\delta\) holds without the
restriction \(\rho\le2\delta\).  The same argument for \(b,q\)
gives \(t=qw'\), \(e_qt=t\), and the identical bounds.
Because \(\|a-d_a\|\le2\delta\), these comparisons prove
\eqref{eq:bounded-diagonals}.  Finally
\[
 \|s^*ct\|
 \le\|s\|\|t\|\|e_p c e_q\|
 \le\delta^{-1}\|e_p c e_q\|<\zeta.
\]
The bound in \eqref{eq:final-coefficient-bound} depends on
\(\delta,\rho\), not on \(p,q\) or the separation accuracy.
This is why the preceding choice
\eqref{eq:bounded-chosen-separation} remains valid after the
comparison coefficients are found.
\end{proof}

\begin{proof}[Proof of Theorem~\ref{thm:individual-to-strong}]
For arbitrary \(a,b\in C_+\), \(c\in C\), and
\(\varepsilon>0\), put \(M=\max\{1,\|a\|,\|b\|\}\).
Choose \(0<\delta<1/2\) and \(0<\rho\le2\delta\)
with \(M(2\delta+\rho)<\varepsilon\).  Apply
Proposition~\ref{prop:bounded-diagonalization} to \(a/M,b/M,c\)
with \(\zeta=\varepsilon\).  Multiplying the two diagonal
inequalities by \(M\) gives all three inequalities in
\eqref{eq:main-diagonalization}.

For positive \(x,y\in C\), specialize these inequalities to
\((a,b,c)=(x^2,y^2,xy)\).  They are precisely
\eqref{eq:SPI-definition}, so \(C\) is strongly purely infinite.
This uses the unrestricted definition directly.  If one instead
writes a block matrix
\(Z=\left(\begin{smallmatrix}a&c\\c^*&b\end{smallmatrix}\right)\),
the norm of
\(\operatorname{diag}(s,t)^*Z\operatorname{diag}(s,t)-(a\oplus b)\)
is at most the maximum of the two diagonal errors plus
\(\|s^*ct\|\).  Hence the same estimates also give diagonal
approximation of every such block matrix, in particular of every
positive one.  The argument includes zero elements and the zero
algebra.
\end{proof}

\section{Removing amplification under nuclearity}\label{sec:nuclear}

We now remove the fixed amplification under nuclearity.  This proof
uses the graph construction of Section~\ref{sec:graph}: a finite tuple
of matrix entries rules out the graph's failure relations, and the
resulting shrinking operation allows us to join rich positive
elements detected by different representations.  Its second output
need not be rich; the merger retains \(n\) orthogonal copies so that
\(\propF\) makes their sum rich.  A finite ideal argument then reaches every positive
element.  The proof of the following theorem is
independent of the exact localization argument in
Sections~\ref{sec:exactness}--\ref{sec:localization}.

\begin{theorem}[Removing amplification under nuclearity]
\label{thm:nuclear-to-individual}
Let \(C\) be a nuclear complex \(C^*\)-algebra.  Suppose that one
integer \(n\ge1\), fixed for the whole algebra, makes
\(h^{\oplus n}\) properly infinite for every \(h\in C_+\).
Then every positive element of \(C\) is properly infinite.
\end{theorem}

Theorem~\ref{thm:nuclear-to-individual} assumes neither separability
nor a unit.  We keep the fixed integer \(n\) throughout this section.
In the tuple construction it controls the number of graph entries,
while the dimensions of the finite matrix approximations may vary.

\subsection{A finite tuple that exact algebras cannot realize}

The graph measurements are fixed before their entries are chosen.
To choose those entries under nuclearity, we use a finite tuple whose
linear matrix norms cannot all occur in an exact algebra, even though
each fixed such norm is a limit of norms of finite matrix tuples.
Nonexact examples originate in work of Wassermann
\cite{Wassermann1976}.  Pisier develops the operator-space formulation
and finite compression mechanism in
\cite[Theorems~1 and~7, Lemma~2 and Equation~(7)]{Pisier1995}.
We use the existence of a nonexact algebra from this theory and give
the tensor norm and compression arguments needed here.  They retain
the quantifiers across all matrix levels required for the graph entries.

Recall that a \(C^*\)-algebra \(D\) is \emph{exact} if tensoring every
short exact sequence with \(D\) in the minimal tensor norm preserves
exactness.  Thus exactness identifies the quotient norm on
\((D\otimes_{\min}J)/(D\otimes_{\min}I)\) with the spatial norm on
\(D\otimes_{\min}(J/I)\), for every closed ideal \(I\triangleleft J\).
The next proof starts from a strict inequality between these norms.

\begin{lemma}[A finite tuple and finite matrix approximations]
\label{lem:forbidden-tuple}
There are \(\ell\ge1\), a \(C^*\)-algebra \(W\), contractions
\(z_1,\ldots,z_\ell\in W\), and contractions
\(z_j^{(r)}\in M_{N_r}(\C)\), where \(N_r\ge1\), \(r\ge1\), and
\(1\le j\le\ell\), with the following properties.
\begin{enumerate}[label=\textup{(\roman*)}]
\item No tuple \(b_1,\ldots,b_\ell\) in any exact \(C^*\)-algebra
satisfies
\begin{equation}\label{eq:forbidden-tuple}
 \left\|\sum_{j=1}^{\ell} A_j\otimes b_j\right\|
 =\left\|\sum_{j=1}^{\ell} A_j\otimes z_j\right\|
 \quad\text{for every }m\ge1\text{ and every }A_j\in M_m(\C).
\end{equation}
\item For every fixed \(m\ge1\) and \(A_j\in M_m(\C)\),
\begin{equation}\label{eq:tuple-approximations}
 \lim_{r\to\infty}
 \left\|\sum_{j=1}^{\ell} A_j\otimes z_j^{(r)}\right\|
 =\left\|\sum_{j=1}^{\ell} A_j\otimes z_j\right\|.
\end{equation}
\end{enumerate}
\end{lemma}

\begin{proof}
Take a nonexact \(C^*\)-algebra \(W\), whose existence is recorded
in \cite[Remark~(i) following Theorem~1, pp.~163--164]{Pisier1995}.
There are
\(I\triangleleft J\) such that the canonical surjection
\[
 (W\otimes_{\min}J)/(W\otimes_{\min}I)
 \longrightarrow W\otimes_{\min}(J/I)
\]
is not injective.  The map is surjective because its algebraic image
is dense and the range of a \(C^*\)-homomorphism is closed.  Choose a
representative \(T_0\in W\otimes_{\min}J\) of a nonzero class in its
kernel.  Its norm modulo \(W\otimes_{\min}I\) is positive, whereas
its image in \(W\otimes_{\min}(J/I)\) is zero.  A sufficiently close
finite algebraic approximation to \(T_0\) retains a strict gap between
these two norms.  After normalizing its nonzero first factors and
absorbing their norms into the second factors, we obtain contractions
\(z_j\in W\) and elements \(h_j\in J\) for which
\begin{equation}\label{eq:tuple-gap}
 \left\|\sum_j z_j\otimes h_j+W\otimes_{\min}I\right\|
 >\left\|\sum_j z_j\otimes(h_j+I)\right\|_{\min}.
\end{equation}
For a tuple \(\mathbf v=(v_j)_{j=1}^{\ell}\) in a \(C^*\)-algebra
\(D\), write
\[
 \mu_{\mathbf v,m}(A_1,\ldots,A_\ell)
 =\left\|\sum_{j=1}^{\ell}A_j\otimes v_j\right\|_{M_m(D)}.
\]
We show that these functions determine both norms in
\eqref{eq:tuple-gap}, with the coefficient algebra and its ideal
held fixed.

For arbitrary coefficients \(d_j\in J\), represent the fixed
nonzero algebra \(J\) faithfully on a Hilbert space \(\mathcal H\).
Then
\begin{equation}\label{eq:tuple-compression-formula}
 \left\|\sum_j v_j\otimes d_j\right\|_{\min}
 =\sup_{m,V}
 \mu_{\mathbf v,m}(V^*d_1V,\ldots,V^*d_\ell V),
\end{equation}
where \(V:\C^m\to\mathcal H\) ranges over finite-dimensional
isometries.  Indeed, also represent \(D\) faithfully.  The operators
on the right are the two-sided compressions in the second Hilbert
factor, so their norms are at most the left-hand norm.  Finite-rank
projections, directed by inclusion, converge strongly to the
identity.  Their two-sided compressions converge strongly to the
tensor sum; evaluation at a fixed vector whose image nearly attains
its norm gives the reverse inequality.  This applies to arbitrary,
including nonselfadjoint, tensor sums.

For the quotient norm, choose a positive contractive approximate
identity \((f_\lambda)\) of the fixed ideal \(I\).  Put
\(J_D=D\otimes_{\min}I\).  For \(T\in D\otimes_{\min}J\), define
\[
 R_\lambda(T)
 =(1\otimes(1-f_\lambda))T(1\otimes(1-f_\lambda)),
\]
using the unitized multiplier operations.  Elementary tensors and
density give
\[
 T-R_\lambda(T)\in J_D,\qquad
 \|R_\lambda(K)\|\longrightarrow0\quad\text{for each fixed }K\in J_D.
\]
The cutdowns are contractions.  Hence, for every fixed \(K\in J_D\),
\[
 \dist(T,J_D)\le\|R_\lambda(T)\|
 \le\|T-K\|+\|R_\lambda(K)\|.
\]
Taking lower and upper limits and then the infimum over \(K\) yields
the net limit
\begin{equation}\label{eq:tuple-quotient-formula}
 \left\|\sum_j v_j\otimes d_j+J_D\right\|
 =\lim_\lambda
 \left\|\sum_j v_j\otimes
 (1-f_\lambda)d_j(1-f_\lambda)\right\|_{\min}.
\end{equation}
The cutdown coefficients lie in \(J\), and no quasicentrality is
needed.  This argument uses convergence on each fixed element of
\(J_D\).  A related fixed-element cutdown mechanism appears in
\cite[proof of Lemma~3, p.~165]{Pisier1995}; the calculation above
establishes the stated formula for general \(I\triangleleft J\) and
its approximate-identity net.  Equations~\eqref{eq:tuple-compression-formula} and
\eqref{eq:tuple-quotient-formula} show that the displayed quotient
norm depends only on the functions \(\mu_{\mathbf v,m}\).  Applying
the same compression argument in \(J/I\) shows the same for
\(\|\sum_j v_j\otimes(d_j+I)\|_{\min}\).

If a tuple in an exact algebra had the same functions
\(\mu_{\mathbf z,m}\) at every matrix level, the two formulas above
would transfer this strict gap to that algebra.  Exactness makes its
canonical quotient map isometric, a contradiction.  This proves
(i).  The conclusion excludes simultaneous equality for all finite
matrix tests; it does not assert that a fixed finite list of tests
suffices.

For (ii), set \(W_0=C^*(z_1,\ldots,z_\ell)\).  This algebra is
separable, and its inclusion in \(W\) preserves matrix norms.  It
has a faithful representation on a separable Hilbert space.  Indeed,
choose a dense sequence in the positive unit sphere and, for each
member, a state attaining its norm.  Their GNS representations have
faithful direct sum: for every nonzero positive element, approximation
by that sequence shows that some chosen state is positive on it.
Choose increasing nonzero
finite-rank projections \(P_r\to1\) strongly in that representation,
and set
\[
 N_r=\dim P_r\mathcal H,\qquad
 z_j^{(r)}=P_rz_jP_r\big|_{P_r\mathcal H}.
\]
These are contractions.  For each fixed \(m\) and \((A_j)\), the
corresponding two-sided compressions of \(\sum_j A_j\otimes z_j\)
converge strongly to that operator and have norms at most its norm.
Evaluation at a fixed near-norming vector gives the matching lower
limit, proving \eqref{eq:tuple-approximations}.  The same sequence
works for every fixed test.  The assertion is pointwise in the test,
and uses no multiplicativity of the compressions.
\end{proof}

\subsection{Choosing graph entries under nuclearity}

We use the finite tuples from Lemma~\ref{lem:forbidden-tuple} to
choose entries in a spectral hereditary subalgebra.  The central
point is that the transport lemma provides one rectangular map for
the whole tuple.  Its conjugation preserves the norm of every joint linear
matrix expression, not merely the norms of the individual entries.

\begin{proposition}[Nuclear avoidance of the fixed measurements]
\label{prop:nuclear-avoidance}
Let \(C\) be nuclear and satisfy \(\propF\) for a fixed \(n\ge1\),
and let \(\ell\) be as in Lemma~\ref{lem:forbidden-tuple}.  For
\(X\in C_+\) and fixed elements
\(B_{j,pq}\in C\), where \(1\le j\le\ell\) and \(1\le p,q\le n\),
there are contractions
\[
 S_{j,pq}\in\Her((X-\alpha_0)_+)
\]
such that no representation \(\pi\) satisfies simultaneously
\begin{equation}\label{eq:nuclear-forbidden}
 \|\pi(X)\|>s,\qquad
 \pi(B_{j,pq})=\pi(S_{j,pq})
 \quad(1\le j\le\ell,\ 1\le p,q\le n).
\end{equation}
Here \(\alpha_0=1/32\) and \(s=1/4\), as in
Proposition~\ref{prop:graph}.
\end{proposition}

\begin{proof}
Put \(g=(X-\alpha_0)_+\), choose \(0<\nu<s-\alpha_0\), and put
\(h_g=(g-\nu)_+\).  Choose a ramp \(t'=t'(X)\) that is zero on
\([0,\alpha_0+\nu]\) and one on \([s,\infty)\).  Then
\(t'\in\Her(h_g)\), and
\begin{equation}\label{eq:nuclear-tuple-signal}
 \|\pi(X)\|>s\quad\Longrightarrow\quad\|\pi(t')\|=1.
\end{equation}
The element \(v=g^{\oplus n}\in M_n(C)\) is properly infinite by
\(\propF\), applied to the scalar positive \(g\).  For the dimensions
\(N_r\) of Lemma~\ref{lem:forbidden-tuple},
\[
 d_r=g^{\oplus N_r}\precsim v^{\oplus N_r}\precsim v.
\]
The first comparison uses \(g\) as a summand of \(v\); the second
uses proper infiniteness of \(v\).  Lemma~\ref{lem:transport}, with
cutoff \(\nu\), gives one rectangular carrier
\[
 w_r\in M_{n,N_r}(C^{**}),\qquad
 w_r^*w_r=1_{N_r}\otimes e_{h_g},\qquad
 w_rw_r^*\le1_n\otimes e_g.
\]
Its conjugation sends \(M_{N_r}(\Her(h_g))\) into
\(M_n(\Her(g))=\Her_{M_n(C)}(v)\).  Use this same carrier for every
\(j\), defining
\begin{equation}\label{eq:nuclear-tuple-choices}
 S_j^{(r)}=w_r(z_j^{(r)}\otimes t')w_r^*
 \in M_n(\Her(g)),\qquad
 S_{j,pq}^{(r)}=(S_j^{(r)})_{pq}.
\end{equation}
Each matrix has norm at most one, so its entries are contractions in
the required hereditary subalgebra.

Suppose that every \(r\) admits a representation \(\pi_r\) satisfying
\eqref{eq:nuclear-forbidden} for these entries.  Restricting to a
reducing cyclic subrepresentation that still detects
\(\|\pi_r(X)\|>s\), we may make these choices from a set of
representations.  Write
\(B_j=[B_{j,pq}]\in M_n(C)\); these elements are fixed independently
of \(r\).  At every matrix level \(m\), the normal amplification of
the single \(w_r\) is isometric on the supported source algebra, by
Lemma~\ref{lem:transport}.  The represented entry equalities and
\eqref{eq:nuclear-tuple-signal} therefore give, for every
\(A_j\in M_m(\C)\),
\begin{align}
 \left\|\sum_j A_j\otimes\pi_r^{(n)}(B_j)\right\|
 &=\left\|\left(\sum_j A_j\otimes z_j^{(r)}\right)
       \otimes\pi_r(t')\right\|\notag\\
 &=\left\|\sum_j A_j\otimes z_j^{(r)}\right\|.
 \label{eq:nuclear-exact-matrix-norms}
\end{align}
This equality is exact at each \(r\).  Using a different carrier for
each \(j\) would not give the joint equality.

Define the \(*\)-homomorphism
\[
 \Psi:M_n(C)\longrightarrow
 \mathcal Q=
 \frac{\prod_r B(\mathcal H_r^{\oplus n})}
      {\bigoplus_r B(\mathcal H_r^{\oplus n})},
 \qquad
 \Psi(b)=[(\pi_r^{(n)}(b))_r],
\]
where the numerator is the algebra of bounded sequences and the
denominator is its ideal of norm-null sequences.  Contractivity of
the representations makes each displayed sequence bounded.  The
image \(E=\Psi(M_n(C))\) is closed and isomorphic to
\(M_n(C)/\ker\Psi\).  Since \(M_n(C)\cong M_n\otimes C\), nuclearity
of matrix algebras and permanence under tensor products make
\(M_n(C)\) nuclear \cite[II.9.4.2 and IV.3.1.1]{Blackadar}.
Its quotient \(E\) is nuclear by
\cite[Corollary~4]{ChoiEffros1977}, and hence exact by
\cite[II.9.6.6]{Blackadar}.
These permanence facts require neither units nor separability.
It is this quotient of a nuclear algebra that must be exact; the
argument requires no exactness of the ambient \(\mathcal Q\).

The norm of a class modulo norm-null sequences is the \(\limsup\) of
its coordinate norms.  Subtracting a norm-null sequence gives the
lower bound, and subtracting finitely many initial entries gives the
upper bound.  The same formula holds at each fixed matrix level under
the canonical identifications of finite matrices over a product and
its norm-null ideal.  Since \(E\subset\mathcal Q\) is completely
isometric, \eqref{eq:nuclear-exact-matrix-norms} and
\eqref{eq:tuple-approximations} give
\[
 \left\|\sum_j A_j\otimes\Psi(B_j)\right\|_{M_m(E)}
 =\limsup_r\left\|\sum_j A_j\otimes\pi_r^{(n)}(B_j)\right\|
 =\left\|\sum_j A_j\otimes z_j\right\|
\]
for every \(m\) and every \((A_j)\).  Only the fixed elements
\(B_j\) have been sent into \(E\); the varying choices in
\eqref{eq:nuclear-tuple-choices} were used in the coordinatewise
equalities.  The displayed identity contradicts
Lemma~\ref{lem:forbidden-tuple}(i).  Thus some \(r\) admits no such
representation, and its entries in
\eqref{eq:nuclear-tuple-choices} prove the proposition.
\end{proof}

\begin{proposition}[Shrinking two corners under nuclearity]
\label{prop:shrinking-nuclear}
Let \(C\) be nuclear and satisfy \(\propF\) for one fixed \(n\ge1\).
Let \(a,y\in C_+\) be contractions, with \(a\) rich, and let
\(c\in\widetilde C\) be a contraction.  There are positive
contractions \(p_1\in\Her(a)\) and \(q_1\in\Her(y)\) such that
\(p_1\) is rich,
\begin{equation}\label{eq:nuclear-shrinking-signals}
 \begin{split}
 \|\pi(a)\|>1/2&\quad\Longrightarrow\quad\|\pi(p_1)\|=1,\\
 \|\pi(y)\|>1/2&\quad\Longrightarrow\quad\|\pi(q_1)\|=1
 \end{split}
 \qquad\text{for every representation }\pi,
\end{equation}
and \(\|e_{p_1}ce_{q_1}\|_{C^{**}}<1\).
\end{proposition}

\begin{proof}
Choose \(\ell\) from Lemma~\ref{lem:forbidden-tuple}, set
\(k=\ell n^2\), and choose graph coefficients
\(\gamma_g,\sigma_g>0\) with \(\gamma_g^2+k\sigma_g^2<1\).
With these choices fixed, apply Proposition~\ref{prop:graph} to obtain
its measurements \(X,B_i\).  Group the \(B_i\) into \(\ell\) matrices
\(B_j=[B_{j,pq}]\in M_n(C)\), and choose the corresponding
\(S_{j,pq}\) by Proposition~\ref{prop:nuclear-avoidance}.
The graph gives a rich positive contraction \(p\in\Her(a)\) with
\(\|\pi(p)\|=1\) whenever \(\|\pi(a)\|>1/2\).  If
\(d=y(1-c^*p^2c)y\) vanishes in a representation where
\(\|\pi(y)\|>s\), the graph would give all of
\eqref{eq:nuclear-forbidden}, contrary to the choice of the entries.
Thus
\begin{equation}\label{eq:nonvanishing-N}
 \pi\bigl(y(1-c^*p^2c)y\bigr)=0
 \quad\Longrightarrow\quad\|\pi(y)\|\le s=1/4
\end{equation}
for every representation.  Lemma~\ref{lem:conditional-shrinking}
applies to this \(p\), its norm signal, and the displayed implication,
and gives the asserted \(p_1,q_1\).  The one fixed \(n\) controlled
both the number of entries and the algebra \(M_n(C)\) in the quotient
argument; only the finite dimensions \(N_r\) varied.
\end{proof}

\subsection{Local rich elements and their merger}

The shrinking proposition will let us join rich elements once they
are available.  For \(u\in C_+\), we first produce such an element inside \(\Her(u)\) at
each irreducible representation that detects \(u\).  This local
construction needs only the fixed-amplification hypothesis
\(\propF\); nuclearity will enter when the elements are joined.

\begin{lemma}[A rich element detected by one representation]
\label{lem:local-rich}
Let \(C\) be a \(C^*\)-algebra, and suppose that one integer
\(n\ge1\) makes \(h^{\oplus n}\) properly infinite for every
\(h\in C_+\).  If \(u\in C_+\) and an irreducible representation
\(\pi:C\to B(\mathcal H)\) satisfies \(\pi(u)\ne0\), then there is
a rich positive contraction \(p\in\Her(u)\) with \(\|\pi(p)\|=1\).
\end{lemma}

\begin{proof}
We will find \(n\) orthogonal targets in \(\Her(u)\), each detected
by \(\pi\), and one positive element comparing into every target.
Transporting \(n\) copies of its cutoff will allow \(\propF\) to
produce richness inside \(\Her(u)\).

Put \(\mathcal H_u=\overline{\pi(u)\mathcal H}\).  This space is
infinite dimensional.  Otherwise \(\pi(u)\) would have finite rank
\(r>0\).  Applying the stabilized representation to proper
infiniteness of \(u^{\oplus n}\) would approximate
\(\pi(u)^{\oplus2n}\), of rank \(2nr\), by positive operators of
rank at most \(nr\).  On the \(2nr\)-dimensional range of the target,
each such approximant has a unit vector in its kernel.  Its distance
from the target is therefore at least the least positive eigenvalue
of \(\pi(u)\), contradicting norm approximation.

The action of \(D=\Her(u)\) on \(\mathcal H_u\) is irreducible.
Indeed, \(\pi(u)\) has no kernel on its support, so
\(\pi(u)\xi\ne0\) for \(0\ne\xi\in\mathcal H_u\).
Irreducibility of \(\pi\) gives
\(\overline{\pi(C)\pi(u)\xi}=\mathcal H\).  Applying \(\pi(u)\)
and taking closure yields
\[
 \overline{\pi(uCu)\xi}=\mathcal H_u.
\]
Thus every nonzero vector is cyclic for the restricted hereditary
algebra, which proves irreducibility.

Apply the selfadjoint form of Kadison transitivity
\cite[Theorem~1]{Kadison1957} to the irreducible unitization of this
action.  Choose \(n+1\) orthonormal vectors and prescribe \(n+1\)
distinct real eigenvalues.  A selfadjoint interpolant has the form
\(h_0+\lambda1\), with \(h_0\in D_{\rm sa}\).  The corresponding
eigenvalues of \(\pi(h_0)\) remain distinct, so at most one is zero.
Choose \(n\) nonzero ones.  Nonnegative continuous functions
supported in pairwise disjoint neighborhoods of these values,
avoiding zero and positive at the chosen values, give orthogonal
elements
\[
 b_1,\ldots,b_n\in D_+,\qquad \pi(b_i)\ne0.
\]
The functions vanish at zero, so these elements belong to \(D\)
also when \(D\) is nonunital.

There is an element \(b'\in D_+\), still detected by \(\pi\), with
\(b'\precsim b_i\) for every \(i\).  Start with \(e'=b_1\).
Given the current \(e'\) and the next \(b_i\), choose a vector on
which \(\pi(e')\) is nonzero.  Cyclicity for \(D\) and nonzero
\(\pi(b_i)\) give \(h\in D\) such that \(\pi(z)\ne0\), where
\(z=b_ihe'\).  Then
\[
 z^*z\in\Her(e'),\qquad zz^*\in\Her(b_i),\qquad
 z^*z\CuEq zz^*.
\]
Replacing \(e'\) by \(z^*z\) preserves nonzero image and every
previous comparison, by transitivity.  At the end we obtain \(b'\)
and
\begin{equation}\label{eq:local-rich-comparison}
 (b')^{\oplus n}\precsim b_1\oplus\cdots\oplus b_n
 \CuEq\sum_{i=1}^{n}b_i\precsim u.
\end{equation}
The middle equivalence uses orthogonality, and the last comparison
uses \(\sum_i b_i\in\Her(u)\).

Choose \(\tau>0\) with \(2\tau<\|\pi(b')\|\).  Transport the
\(n\) source copies in \eqref{eq:local-rich-comparison}, after cutoff
\(\tau\), into \(\Her(u)\).  Let \(t\) be a ramp of \(b'\), zero
on \([0,\tau]\) and one on \([2\tau,\infty)\).  It belongs to the
common cutoff hereditary algebra and has represented norm one.
The sum \(p\) of its \(n\) orthogonal transported copies is a
positive contraction in \(\Her(u)\) with \(\|\pi(p)\|=1\).
By \eqref{eq:orthogonal-copy-cutoffs}, for every \(\varepsilon\ge0\),
\[
 (p-\varepsilon)_+\CuEq((t-\varepsilon)_+)^{\oplus n}.
\]
The right side is properly infinite by the hypothesis applied to
the scalar positive \((t-\varepsilon)_+\).  Thus \(p\) is rich.
\end{proof}

We now join two rich elements while retaining the union of their
norm signals.  At each step Proposition~\ref{prop:shrinking-nuclear}
is used with \(c=1\) in the unitization.  Its support bound is one
global norm below one for that step; the numerical bound may change
at the next step.

\begin{lemma}[Merging two rich elements]\label{lem:merger}
Let \(C\) be nuclear and satisfy \(\propF\) for a fixed \(n\ge1\).
For \(u\in C_+\) and rich positive contractions \(a,b\in\Her(u)\),
there is a rich positive contraction \(c_0\in\Her(u)\) such that
every representation \(\pi\) satisfies
\begin{equation}\label{eq:merger-signal}
 \bigl(\|\pi(a)\|>1/2\ \text{or}\ \|\pi(b)\|>1/2\bigr)
 \quad\Longrightarrow\quad\|\pi(c_0)\|=1.
\end{equation}
\end{lemma}

\begin{proof}
All comparisons remain in \(C\otimes\Kc\), and all uses of shrinking
are in \(C\).  Since \(b\) is properly infinite,
\(b^{\oplus n}\precsim b\).  Apply Lemma~\ref{lem:transport} with
cutoff \(1/8\) to obtain orthogonal copy maps
\[
 \theta_i(h)=w_i h w_i^*:
 \Her((b-1/8)_+)\longrightarrow\Her(b)\subset\Her(u)
 \quad(1\le i\le n).
\]
Let \(y_0\) be a ramp of \(b\), zero on \([0,1/4]\) and one on
\([1/2,1]\), and set \(a_0=a\).  For \(1\le j\le n\) we construct
positive contractions \(a_j,y_j\) such that \(a_j\) is rich and
belongs to \(\Her(u)\), \(y_j\) belongs to the common domain of
the copy maps, and the following invariants hold:
\begin{equation}\label{eq:merger-induction}
 \begin{gathered}
 a_j\theta_i(y_j)=0\qquad(1\le i\le j),\\
 \|\pi(a)\|>1/2\ \Longrightarrow\ \|\pi(a_j)\|=1,\\
 \|\pi(b)\|>1/2\ \Longrightarrow\ \|\pi(y_j)\|=1
 \quad\text{for every representation }\pi.
 \end{gathered}
\end{equation}

Once the induction reaches \(j=n\), we will set
\(c_0=a_n+\sum_{i=1}^n\theta_i(y_n)\).  The induction will make the
summands orthogonal, and the final use of \(\propF\) will prove the
sum of the \(n\) copies rich without requiring \(y_n\) itself to be
rich.

Suppose \(a_{j-1},y_{j-1}\) have been constructed.  At \(j=1\),
use the initial pair \(a_0,y_0\); at later stages use the preceding
invariants.  Apply
Proposition~\ref{prop:shrinking-nuclear} to
\(a_{j-1},\theta_j(y_{j-1})\), and \(1\in\widetilde C\), and call
its outputs \(p_j,q_j\).  Thus \(p_j\) is rich and belongs to
\(\Her(a_{j-1})\), while
\[
 \kappa_j=\|e_{p_j}e_{q_j}\|<1.
\]
The hereditary-image equality
\eqref{eq:transport-hereditary-image} supplies a positive contraction
\(z_j\in\Her(y_{j-1})\) such that
\begin{equation}\label{eq:merger-pullback}
 q_j=\theta_j(z_j).
\end{equation}
This pullback is an algebra element.  Its represented norms equal
those of \(q_j\) by \eqref{eq:transport-represented-norm}.

For \(0\le t\le1\), put \(f(t)=(2t-1)_+\) and
\(g(t)=\min\{1,2t\}\).  We next move \(p_j\) away from a cutoff
of \(q_j\).  The support
bound ensures that multiplication by \(1-g(q_j)\) is bounded below
on the support of \(p_j\).  The polar part of
\((1-g(q_j))p_j\) will therefore transport \(p_j\), preserving its
cutoff Cuntz classes and represented norms.
Define
\[
 y_j=f(z_j),\qquad H_j=1-g(q_j),\qquad T_j=H_jp_j.
\]
Here \(T_j=p_j-g(q_j)p_j\in\Her(u)\).  In the universal
representation, \(g(q_j)\) is supported under \(e_{q_j}\).
For \(\xi\in e_{p_j}\mathcal H\), this gives
\[
 \|H_j\xi\|^2
 \ge\|(1-e_{q_j})H_j\xi\|^2
 =\|(1-e_{q_j})\xi\|^2
 \ge(1-\kappa_j^2)\|\xi\|^2.
\]
Since \(H_j\) is a contraction, we obtain the order bounds
\begin{equation}\label{eq:merger-polar-order}
 (1-\kappa_j^2)p_j^2\le T_j^*T_j\le p_j^2.
\end{equation}
Let \(T_j=l_j|T_j|\) be its polar decomposition in the bidual.
The order bounds give \(\Her(|T_j|)=\Her(p_j)\) and
\(l_j^*l_j=e_{p_j}\).  The polar membership argument in the proof
of Lemma~\ref{lem:transport}, applied now to \(T_j\), gives
\[
 a_j=l_jp_jl_j^*\in\Her(T_jT_j^*)\subset\Her(u).
\]
It is a positive contraction.  The same argument makes each
transported cutoff Cuntz equivalent to the corresponding cutoff of
\(p_j\), so \(a_j\) is rich.  The normal extensions of the
representations preserve \(l_j^*l_j=e_{p_j}\); conjugating back by
\(l_j^*\) shows that \(a_j\) and \(p_j\) have equal norm in every
representation.

We check that this change preserves every earlier orthogonality
relation.  Since \(gf=f\), functional calculus and
\eqref{eq:merger-pullback} give
\[
 \theta_j(y_j)H_j=f(q_j)(1-g(q_j))=0,
\]
and hence \(\theta_j(y_j)T_j=0\).  For \(i<j\), the previous
orthogonality and \(p_j\in\Her(a_{j-1})\) give
\(\theta_i(y_{j-1})p_j=0\), while orthogonality of the copy ranges
gives \(\theta_i(y_{j-1})g(q_j)=0\).  Thus
\(\theta_i(y_{j-1})T_j=0\).  Each of these positive elements
annihilates \(\overline{\ran(T_j)}\), whose projection is
\(l_jl_j^*\), and so annihilates \(a_j\).  Because
\(y_j\in\Her(y_{j-1})\), the same is true of \(\theta_i(y_j)\)
for \(i<j\).  This proves the orthogonality in
\eqref{eq:merger-induction}.

The norm signals also persist simultaneously.  Wherever
\(\|\pi(a)\|>1/2\), the first shrinking step gives
\(\|\pi(p_1)\|=1\), hence \(\|\pi(a_1)\|=1\); that signal is
above \(1/2\) and survives every subsequent step.  Wherever
\(\|\pi(b)\|>1/2\), the initial ramp has \(\|\pi(y_0)\|=1\).
If \(\|\pi(y_{j-1})\|=1\), its copy under \(\theta_j\) has norm
one, shrinking gives \(\|\pi(q_j)\|=1\), and transport gives
\(\|\pi(z_j)\|=1\).  Since \(f(1)=1\), spectral mapping gives
\(\|\pi(y_j)\|=1\).

Finally set
\[
 c_0=a_n+\sum_{i=1}^{n}\theta_i(y_n).
\]
All summands are orthogonal positive contractions in \(\Her(u)\).
Their sum is a contraction, and in each representation its norm is
the maximum of the summand norms.  The preceding signals prove
\eqref{eq:merger-signal}.  For every \(\tau\ge0\), orthogonal
functional calculus and \eqref{eq:orthogonal-copy-cutoffs} give
\[
 (c_0-\tau)_+
 \CuEq(a_n-\tau)_+\oplus((y_n-\tau)_+)^{\oplus n}.
\]
The first summand is properly infinite because \(a_n\) is rich.
The second is properly infinite by \(\propF\), applied to the
scalar positive \((y_n-\tau)_+\).  Their block sum is properly
infinite, so \(c_0\) is rich.  The construction includes zero
inputs through zero copy maps and zero cutoffs.
\end{proof}

\subsection{Finite ideal assembly and absorption}

Lemma~\ref{lem:local-rich} supplies a rich element wherever \(u\) is
detected, and Lemma~\ref{lem:merger} joins any finite family while
retaining all of its norm signals.  These two facts put
\((u-\varepsilon)_+\), for every \(\varepsilon>0\), in the ideal of
a single rich element.  Ideal comparison then supplies the doubled
cutoff comparison needed for Theorem~\ref{thm:nuclear-to-individual}.

\begin{proof}[Proof of Theorem~\ref{thm:nuclear-to-individual}]
Fix \(u\in C_+\), and let \(\mathcal R_u\) be the set of rich positive
contractions in \(\Her(u)\).  Let \(J\) be the closed ideal generated
by all \((p-1/2)_+\) with \(p\in\mathcal R_u\).  We claim that
\(u\in J\).  Otherwise an irreducible representation of \(C/J\)
would detect \(u+J\).  Its composition with the quotient map is an
irreducible representation of \(C\), so Lemma~\ref{lem:local-rich}
gives \(p\in\mathcal R_u\) with represented norm one.  Then
\((p-1/2)_+\) has nonzero image, contrary to its membership in \(J\).

The ideals generated by finite subfamilies of these cutoffs form an
increasing directed family with dense union in \(J\).  For
\(\varepsilon>0\), Lemma~\ref{lem:ideals-cutoffs}(i) therefore puts
\((u-\varepsilon)_+\) in the ideal generated by
\((p_i-1/2)_+\) for finitely many \(p_1,\ldots,p_r\in\mathcal R_u\).
Starting from the rich zero element, merge them successively by
Lemma~\ref{lem:merger}.  The result is a rich \(p\in\Her(u)\) for
which every representation satisfies
\[
 \|\pi(p_i)\|>1/2\text{ for some }i
 \quad\Longrightarrow\quad\|\pi(p)\|=1.
\]
A norm-one signal survives the next merger because it is above
\(1/2\).  The primitive ideal test in
Lemma~\ref{lem:ideals-cutoffs}(ii) now gives
\((p_i-1/2)_+\in\Id(p)\) for every \(i\).  Consequently
\begin{equation}\label{eq:nuclear-final-ideal}
 (u-\varepsilon)_+\in\Id(p).
\end{equation}
If the cutoff is zero, the zero element already suffices.

Since \(p\) is properly infinite, ideal comparison with this target in
the stabilization, followed by hereditary comparison, yields
\begin{equation}\label{eq:nuclear-final-comparison}
 ((u-\varepsilon)_+)^{\oplus2}\precsim p\precsim u.
\end{equation}
The target \(u\) is fixed as \(\varepsilon\downarrow0\).  Explicitly,
choose \(\varepsilon_j\downarrow0\) and, for each \(j\), a comparison
coefficient in \eqref{eq:nuclear-final-comparison} with error less
than \(\varepsilon_j\) to the doubled cutoff.  That cutoff differs
from \(u\oplus u\) by at most \(\varepsilon_j\).  The resulting
coefficients therefore compare \(u\oplus u\) to \(u\) in norm.
Thus \(u\) is properly infinite.  The fixed \(n\) was applied only
to scalar positives and their scalar cutoffs throughout the proof.
\end{proof}

\begin{remark}[Extension of the tuple proof to exact algebras]
\label{rem:tuple-exact}
The amplification-removal argument uses nuclearity only to obtain
exactness of \(E=\Psi(M_n(C))\) in
Proposition~\ref{prop:nuclear-avoidance}.  Quotients of exact
\(C^*\)-algebras are exact, including in the nonseparable case
\cite[IV.3.4.20]{Blackadar}.  Finite matrix amplification also
preserves exactness, since its tensor products and ideal kernels
are the corresponding finite matrices over those for \(C\).
This deeper quotient-permanence theorem therefore makes the same
avoidance, shrinking, merger, and ideal assembly valid when \(C\)
is merely exact.  It gives another proof of the amplification-removal
conclusion of Theorem~\ref{thm:exact-infiniteness}.  The perturbation
and localization proof below obtains that conclusion without invoking
quotient permanence of exactness.
\end{remark}

\begin{proof}[Proof of Corollary~\ref{cor:absorption}]
Set \(C=A\), the separable nuclear algebra in the corollary, with
its stated fixed-amplification hypothesis.
Theorem~\ref{thm:nuclear-to-individual} gives \(\propP\), and
Theorem~\ref{thm:individual-to-strong} gives strong pure infiniteness.
The general nonunital absorption implication is recorded in the
paragraph following \cite[Definition~3.1]{KS}.  We spell out the
central-sequence interface that includes the nonunital case.

Let \(C_\omega\) be the norm ultrapower of \(C\) along a free
ultrafilter on \(\N\), with \(C\) embedded diagonally, and put
\[
 F(C)=\frac{C_\omega\cap C'}{\operatorname{Ann}(C,C_\omega)},
 \qquad
 \operatorname{Ann}(C,C_\omega)
 =\{x\in C_\omega:xC=Cx=0\}.
\]
Suppose first that \(C\ne0\).  Separability makes \(C\)
\(\sigma\)-unital, so \(F(C)\) is unital by
\cite[Corollary~1.10(1)]{Kirchberg2006}.  The central-sequence embedding
statement immediately before \cite[Proposition~4.4]{Kirchberg2006}
gives a unital embedding \(\Oinf\hookrightarrow F(C)\): its ambient
original algebra is the strongly purely infinite \(C\), and the
diagonal copy of \(C\) is the required separable nuclear subalgebra
of \(C_\omega\).  The algebra \(\Oinf\) is unital and separable, and
Remark~4.6 supplies self-absorption and connectedness of its unitary
group.  The ``in particular'' clause of Proposition~4.4(4) therefore
gives stable \(\Oinf\)-absorption.  Since \(U(\Oinf)=U_0(\Oinf)\), where \(U_0\)
is the identity component, the commutator subgroup of \(U(\Oinf)\)
lies in \(U_0(\Oinf)\); part (5) therefore removes stabilization.  Hence
\(C\cong C\otimes_{\min}\Oinf\).  The Proposition permits \(C\) to
be nonunital; the unital embedding is into the annihilator quotient
\(F(C)\).  If \(C=0\), the asserted isomorphism is immediate.
\end{proof}

\section{Uniform perturbation from exactness}\label{sec:exactness}

We now begin the exact proof that \propF\ implies \propP.  Its first
step is a perturbation estimate that applies simultaneously to every
primitive quotient in which a fixed positive element has norm at least
a prescribed positive threshold.  The constant in that estimate comes
from the following obstruction: there is one positive error threshold,
independent of the unital exact algebra and the matrix size,
below which a contraction cannot contain approximately reducing copies
of every finite matrix contraction.

\subsection{The universal contraction obstruction}

Recall that a $C^*$-algebra $B$ is \emph{exact} if, for every closed
two-sided ideal $I$ of every $C^*$-algebra $C$, the canonical map
\[
 (B\otimes_{\min}C)/(B\otimes_{\min}I)
 \longrightarrow B\otimes_{\min}(C/I)
\]
is an isomorphism.  Here and below, homomorphisms of $C^*$-algebras are
$*$-homomorphisms.  In the obstruction, the same finite-dimensional
copy must approximately intertwine both a contraction and its adjoint.

\begin{proposition}[Uniform gap]\label{prop:uniform-gap}
There is a constant $\gamma>0$ such that no unital exact $C^*$-algebra
$B$ contains a contraction $v$ with the following property: for every
integer $d\geq1$ and every contraction $D\in M_d(\C)$, there are a unital
representation $\pi:B\to B(H_\pi)$ and an isometry
$J:\C^d\to H_\pi$ satisfying
\begin{equation}\label{eq:forbidden-models}
 \|\pi(v)J-JD\|<\gamma,
 \qquad
 \|\pi(v)^*J-JD^*\|<\gamma.
\end{equation}
The representation and the isometry are allowed to depend on $d$ and
$D$.
\end{proposition}

The proof has two parts.  The universal unital algebra of a contraction
fails exactness, as witnessed by a tensor quotient.  If exact targets
admitted the proposed reducing copies with arbitrarily small common
error, their evaluations would give a lower bound on every matrix norm
of each fixed polynomial space.  Holding that space and one evaluation
fixed while an ideal approximate unit varies
would force the tensor quotient map to be injective.  The two relations
in \eqref{eq:forbidden-models} are what allow this argument to test words
containing both the generator and its adjoint.

Let $\U=C^*(1,u)$ be the universal
unital $C^*$-algebra generated by a contraction.  Thus every contraction
$T$ on a Hilbert space determines a unital homomorphism $u\mapsto T$.
One construction is to complete the unital $*$-algebra of polynomials
in $u,u^*$ in the supremum seminorm over all such evaluations, after
dividing out its null space.  The seminorm of each polynomial is finite
because a word in a contraction and its adjoint has norm at most one.

We shall use the following elementary quotient-norm formula.  If $D$
is unital and $(e_\lambda)$ is a positive contractive approximate unit
of an ideal $I\subset C$, then, for $z\in D\otimes_{\min}C$,
\begin{equation}\label{eq:tensor-quotient-norm}
 \|z+D\otimes_{\min}I\|
 =\lim_\lambda\|z(1\otimes(1-e_\lambda))\|.
\end{equation}
When $C$ is nonunital, the multiplication is performed in its
unitization; its result still belongs to $D\otimes_{\min}C$.
Indeed, $1\otimes e_\lambda$ is an approximate unit for
$D\otimes_{\min}I$.  The product on the right represents the same
quotient class as $z$, so its norm bounds the left side from below.
For any $y\in D\otimes_{\min}I$, contractivity gives
\[
 \limsup_\lambda\|z(1\otimes(1-e_\lambda))\|
 \leq \|z-y\|
       +\limsup_\lambda\|y(1\otimes(1-e_\lambda))\|
 =\|z-y\|.
\]
Taking the infimum over $y$ proves the formula, including existence of
the limit.

The formula also gives the form of subalgebra permanence that we need:
exactness passes to unital $C^*$-subalgebras; see also
\cite[IV.3.4.3]{Blackadar}.  For a unital inclusion $D\subset B$, the spatial
tensor norms on $D\otimes C$ and $D\otimes(C/I)$ are the norms
inherited from the corresponding tensor products with $B$.
Formula~\eqref{eq:tensor-quotient-norm} therefore gives, for
$z\in D\otimes_{\min}C$,
\[
 \|z+D\otimes_{\min}I\|
 =\|z+B\otimes_{\min}I\|.
\]
If $B$ is exact, this equals
$\|(\id_D\otimes q)z\|$, where $q:C\to C/I$ is the quotient map.
This proves exactness of $D$.

\begin{lemma}\label{lem:nonexact-universal}
The universal contraction algebra $\U$ is not exact.
\end{lemma}

\begin{proof}
We first embed the full group algebra $Q=C^*(\F_2)$ into $\U$.
The two free generators are sent to
\[
 \exp(2\pi i\operatorname{Re}u),
 \qquad \exp(2\pi i\operatorname{Im}u).
\]
To see that the resulting homomorphism is injective, take any pair of
unitaries $S,T$ on a Hilbert space.  Bounded Borel functional calculus
gives selfadjoint contractions $H,K$ such that
$S=\exp(\pi iH)$ and $T=\exp(\pi iK)$.  The contraction
$(H+iK)/2$ evaluates the two displayed unitaries to $S,T$.
Thus every unital representation of $Q$ is a restriction of an evaluation of
$\U$, and the homomorphism is isometric.  The logarithms need only
belong to the bounded operators on the representation space.
Courtney and Sherman give a related embedding with a different
normalization in \cite[Theorem~6.10]{CS}.

Nonexactness of $C^*(\F_2)$ was established by
Wassermann \cite{Wassermann1976}.  Here is an explicit proof for $Q$.
We construct a representation of $Q\otimes_{\min}Q$ that annihilates
$Q\otimes_{\min}I$, where $I$ is the kernel of the regular quotient.
Exactness would make it descend to $Q\otimes_{\min}(Q/I)$.  A concrete
tensor sum, however, has spatial norm in that algebra smaller than the
norm the proposed descended representation would assign to it.
This will exhibit the required failure of the tensor quotient map.
Put $G=\F_2$, and
write $u_s$ for its canonical unitaries in $Q$.  On $\ell^2(G)$ let
\[
 \lambda_s\delta_g=\delta_{sg},
 \qquad \rho_s\delta_g=\delta_{gs^{-1}}.
\]
There is a state $\varphi$ on $Q\otimes_{\min}Q$ with
\begin{equation}\label{eq:diagonal-state}
 \varphi(u_s\otimes u_t)=
 \begin{cases}1,&s=t,\\0,&s\ne t.\end{cases}
\end{equation}
For completeness, finite quotients supplying this state can be
constructed directly.  On the finite word ball of radius $m$ in $G$,
partially defined left translation by each of the two generators is
an injection between subsets of equal size, hence extends to a
permutation of the ball.  Freeness gives a homomorphism $q_m$ from
$G$ to the finite permutation group generated by these extensions.
The inverse of each extension agrees with inverse left translation
whenever both endpoints lie in the ball.
If a reduced word has length at most $m$, applying its permutations
to the identity follows its successive suffixes inside the ball and
ends at that word.  Consequently $q_m$ distinguishes distinct elements
of the word ball of radius $m$.

For the finite group $F_m=q_m(G)$, its commuting left and right
regular representations give a representation of
$C^*(F_m)\otimes_{\min}C^*(F_m)$: these algebras are finite direct
sums of matrix algebras, for which the spatial and maximal tensor
norms agree.  Pulling its vector state at the identity back to
$Q\otimes_{\min}Q$ gives the coefficient
$\boldsymbol{1}_{q_m(s)=q_m(t)}$ on $u_s\otimes u_t$.
These coefficients tend to those in \eqref{eq:diagonal-state}.
The group tensors span a dense $*$-subalgebra, and all the states have
norm one, so their limit on that span extends to a state on the
spatial completion.

The GNS representation of $\varphi$ is precisely
\begin{equation}\label{eq:left-right-gns}
 u_s\otimes u_t\longmapsto\lambda_s\rho_t
 \quad\hbox{on }\ell^2(G),
\end{equation}
with cyclic vector $\delta_e$.  Indeed the cyclic vector associated
with $u_s\otimes u_t$ corresponds to $\delta_{st^{-1}}$; the inner
products agree because
\[
 s^{-1}a=t^{-1}b\quad\Longleftrightarrow\quad st^{-1}=ab^{-1},
\]
and left multiplication by $u_x\otimes u_y$ becomes
$\lambda_x\rho_y$.  These vectors span $\ell^2(G)$.

Let $I=\ker(Q\to C_r^*(G))$.  The linear unitary
$R\delta_g=\delta_{g^{-1}}$ satisfies $R\lambda_sR^*=\rho_s$.
Thus the representation in \eqref{eq:left-right-gns} annihilates
$Q\otimes_{\min}I$.  Exactness of $Q$ would make it descend to a
representation of $Q\otimes_{\min}C_r^*(G)$, in which every
$u_s\otimes\lambda_s$ fixes $\delta_e$.

This contradicts a spatial norm bound.  In a faithful unital
representation $\pi:Q\to B(H)$, the unitary
\[
 V(\xi\otimes\delta_g)=\pi(u_g)\xi\otimes\delta_g
\]
satisfies
$V(1\otimes\lambda_s)V^*=\pi(u_s)\otimes\lambda_s$.
If $a,b$ are the free generators, the norm of
\[
 u_a\otimes\lambda_a+u_{a^{-1}}\otimes\lambda_{a^{-1}}
 +u_b\otimes\lambda_b+u_{b^{-1}}\otimes\lambda_{b^{-1}}
\]
in this faithful spatial representation is therefore the norm of the
adjacency operator on the $4$-regular tree.  That norm is at most
$2\sqrt3<4$.  To check the bound, orient each edge away from the
identity, and for a finitely supported vector $\xi$ bound the
contribution of a parent--child edge by
\[
 2|\xi(p)\xi(c)|
 \leq 3^{-1/2}|\xi(p)|^2+3^{1/2}|\xi(c)|^2.
\]
Every nonroot vertex has three children and one parent, and hence
total coefficient $2\sqrt3$; the root has coefficient
$4/\sqrt3\leq2\sqrt3$.  The quadratic-form bound extends by density
and proves the asserted operator norm bound.  In the purported
descended representation the same sum has eigenvalue $4$ on
$\delta_e$, which is impossible for a $*$-representation.
Thus $Q$ is not exact.  Its unital inclusion in $\U$ and the
inheritance proved above show that $\U$ is not exact either.
\end{proof}

We have obtained the algebraic obstruction.  To turn it into a fixed
error bound, we next test polynomial norms on finite contractions and
control those tests at every matrix amplification.  The latter
uniformity will allow one polynomial model to be held fixed while an
ideal approximate unit varies.

\begin{lemma}[Finite matrix tests]\label{lem:finite-matrix-tests}
For every integer $r\geq1$ and every $Z\in M_r(\U)$ whose entries are
$*$-polynomials in $u,u^*$,
\begin{equation}\label{eq:finite-matrix-tests}
 \|Z\|_{M_r(\U)}
 =\sup_{d\geq1,\ D\in M_d(\C),\ \|D\|\leq1}\|Z(D,D^*)\|.
\end{equation}
\end{lemma}

\begin{proof}
Contractivity of evaluations proves one inequality.  Every unital
representation of $\U$ is evaluation at the image of $u$, so these
evaluations also compute the matrix norm on the left.  For the other
inequality, evaluate at a contraction $T\in B(H)$, and take a unit vector
$\xi=(\xi_1,\ldots,\xi_r)\in H^r$ arbitrarily close to attaining
the norm of $Z(T,T^*)$.  If the entries of $Z$ have degree at most
$m$, let
\[
 H_0=\operatorname{span}
 \{w(T,T^*)\xi_j:|w|\leq m,\ 1\leq j\leq r\},
\]
where the empty word is included.  This is a finite-dimensional
nonzero space.  Set $D=P_{H_0}T|_{H_0}$, so that
$D^*=P_{H_0}T^*|_{H_0}$.  Applying a word successively from the right
never changes the relevant vectors under compression, since all its
suffix vectors lie in $H_0$.  Hence
\[
 w(D,D^*)\xi_j=w(T,T^*)\xi_j\quad (|w|\leq m),
\]
and $Z(D,D^*)\xi=Z(T,T^*)\xi$.  Taking the supremum over evaluations
and over these vectors proves \eqref{eq:finite-matrix-tests}.
The scalar norm-test formula is proved in \cite[Proposition~2.2]{CS};
the finite-word compression used here is the argument in the proof of
\cite[Theorem~3.1]{CS}. Including all components of the test vector
gives the matrix-valued assertion needed here.
\end{proof}

\begin{lemma}[A bound at every amplification]\label{lem:uniform-matrix-bound}
Let $E\subset\U$ be a finite-dimensional linear space consisting of
$*$-polynomials in $u,u^*$.  There is a finite constant $C_E\geq0$
with the following property.  Suppose $v$ is a contraction in a unital
$C^*$-algebra $B$, and fix one number $\theta>0$.  Assume that for every
integer $d\geq1$ and every contraction $D\in M_d(\C)$, there are a unital
representation $\pi:B\to B(H_\pi)$ and an isometry $J:\C^d\to H_\pi$
satisfying the two inequalities in \eqref{eq:forbidden-models} with
$\theta$ in place of $\gamma$.
For the homomorphism $\Phi:\U\to B$ with $\Phi(u)=v$, one then has
\begin{equation}\label{eq:all-amplification-bound}
 \|\Phi^{(r)}(Z)\|\geq(1-C_E\theta)\|Z\|
 \qquad (r\geq1,\ Z\in M_r(E)).
\end{equation}
\end{lemma}

\begin{proof}
If $E=0$, take $C_E=0$.  Otherwise choose a basis $P_1,\ldots,P_k$
of $E$, with coordinate functionals $f_1,\ldots,f_k$, and choose
polynomial representatives
$P_i=\sum_w a_{iw}w$.  Define
\begin{equation}\label{eq:explicit-CE}
 C_E=\sum_{i=1}^k\|f_i\|\sum_w|a_{iw}|\,|w|.
\end{equation}
Constants are allowed in the representatives; the empty word has
length zero.  For $Z=\sum_i B_i\otimes P_i\in M_r(E)$ and scalar
unit vectors $\alpha,\beta\in\C^r$,
\[
 |\alpha^*B_i\beta|
 =|f_i(\alpha^*Z\beta)|\leq\|f_i\|\|Z\|.
\]
It follows that $\|B_i\|\leq\|f_i\|\|Z\|$ for every matrix size $r$.

Given a finite contraction $D$, choose $\pi,J$ as in the hypothesis.
Both $\pi(v)$ and $D$ are contractions, and telescoping the two
intertwining relations gives
\[
 \|w(\pi(v),\pi(v)^*)J-Jw(D,D^*)\|\leq|w|\theta
\]
for every word $w$, including words in both letters.  More explicitly,
the difference for a word $A_1\cdots A_m$ is the sum of $m$ terms,
each consisting of contractions on either side of one difference
$\pi(v)J-JD$ or $\pi(v)^*J-JD^*$.  Put $J_r=1_{\C^r}\otimes J$.
The coefficient bound and \eqref{eq:explicit-CE} yield
\[
 \|\pi^{(r)}(\Phi^{(r)}(Z))J_r-J_r Z(D,D^*)\|
 \leq C_E\theta\|Z\|.
\]
Since $J_r$ is an isometry and $\pi^{(r)}$ is contractive,
\[
 \|Z(D,D^*)\|\leq\|\Phi^{(r)}(Z)\|+C_E\theta\|Z\|.
\]
Taking the supremum supplied by Lemma~\ref{lem:finite-matrix-tests}
proves \eqref{eq:all-amplification-bound}.  In particular, $C_E$
depends on the fixed polynomial space and its chosen representatives,
but on none of $r,d,D,B,v,\pi,J$.
\end{proof}

\begin{proof}[Proof of Proposition~\ref{prop:uniform-gap}]
If the Proposition were false, for every integer $j\geq1$ there
would be a unital exact algebra $B_j$ and a contraction $v_j\in B_j$
satisfying \eqref{eq:forbidden-models} with $1/j$ in place of $\gamma$.
Let $\Phi_j:\U\to B_j$ send $u$ to $v_j$.
Fix a finite-dimensional polynomial space $E$.  Before choosing any
matrix amplification, tensor factor, compression, or ideal, choose one
integer $j=j(E)\geq\max\{1,2C_E\}$.  Lemma~\ref{lem:uniform-matrix-bound}
then gives
\begin{equation}\label{eq:fixed-j-matrix-bound}
 \|Z\|\leq2\|\Phi_j^{(r)}(Z)\|
 \qquad(r\geq1,\ Z\in M_r(E)).
\end{equation}

This same $j$ works after tensoring with an arbitrary $C^*$-algebra
$C$: for every $z\in E\odot C$,
\begin{equation}\label{eq:fixed-j-tensor-bound}
 \|z\|\leq2\|(\Phi_j\otimes\id_C)z\|,
\end{equation}
where $\odot$ denotes the algebraic tensor product and all norms are
spatial.  The assertion is immediate if $C=0$.  Otherwise, represent
$C$ faithfully on a nonzero Hilbert space $K$, and use
faithful representations of $\U$ and $B_j$ for their respective
spatial tensor products.  Compress the $C$ representation to a
nonzero finite-dimensional subspace $K_0\subset K$.  If
$z=\sum_i P_i\otimes c_i$, its compression has coefficients
$P_{K_0}c_i|_{K_0}\in M_d(\C)$ and, after interchanging the tensor
factors, belongs to $M_d(E)$.  Inequality~\eqref{eq:fixed-j-matrix-bound}
bounds its norm by twice the norm of the corresponding compression
on the $B_j$ side, which is at most
$2\|(\Phi_j\otimes\id_C)z\|$.  The norms of these compressions
have supremum $\|z\|$: as $K_0$ increases, the projections
$1\otimes P_{K_0}$ converge strongly to the identity, and their
two-sided compressions converge strongly to $z$.  This proves
\eqref{eq:fixed-j-tensor-bound}.

Now fix any ideal $I\subset C$, with quotient map $q:C\to C/I$
and a positive contractive approximate unit $(e_\lambda)$ of $I$.
For $z\in E\odot C$, every
$z(1\otimes(1-e_\lambda))$ still belongs to $E\odot C$.
Consequently \eqref{eq:fixed-j-tensor-bound} applies to the entire
net with the single index $j(E)$.  Formula~\eqref{eq:tensor-quotient-norm}
and exactness of $B_j$ give
\begin{align*}
 \|z+\U\otimes_{\min}I\|
 &=\lim_\lambda\|z(1\otimes(1-e_\lambda))\|\\
 &\leq2\lim_\lambda
   \|(\Phi_j\otimes\id_C)(z)(1\otimes(1-e_\lambda))\|\\
 &=2\|(\Phi_j\otimes\id_{C/I})((\id_\U\otimes q)z)\|\\
 &\leq2\|(\id_\U\otimes q)z\|.
\end{align*}
Finite sums of polynomial tensors are dense in $\U\otimes_{\min}C$.
Although $j(E)$ may vary with the polynomial space, the resulting
constant $2$ does not.  The last inequality therefore extends by
continuity to every $z\in\U\otimes_{\min}C$.  It says that the
canonical surjective homomorphism
\[
 \kappa:(\U\otimes_{\min}C)/(\U\otimes_{\min}I)
 \longrightarrow\U\otimes_{\min}(C/I)
\]
has zero kernel: if $\kappa(z+\U\otimes_{\min}I)=0$, the quotient
class has norm zero.  An injective $*$-homomorphism of $C^*$-algebras
is isometric, so the fixed factor $2$ suffices to prove that $\kappa$
is an isomorphism.  Since $C,I$ were arbitrary, $\U$ would be exact,
contrary to Lemma~\ref{lem:nonexact-universal}.
\end{proof}

Proposition~\ref{prop:uniform-gap} therefore supplies a single
$\gamma>0$ for all unital exact targets and all matrix sizes.  The remaining
part of this section converts that obstruction into a perturbation
with a quotient norm lower bound.

\subsection{Uniform perturbation in primitive quotients}
\label{sec:perturbation}

We now turn the uniform gap into a perturbation that works in many
primitive quotients at once.  The fixed amplification condition
\propF{} supplies finite reducing copies on a spectral cut; the gap
then prevents one fixed element from nearly cancelling every such
copy.  We first record the two consequences of the comparison and
transport results in Section~\ref{sec:foundations} that this argument
uses.

\begin{lemma}[Fixed-amplification ideal comparison]
\label{lem:ideal-comparison}
Let \(C\) be a \(C^*\)-algebra, let \(h\in C_+\), and let \(n\geq1\)
be an integer such that \(h^{\oplus n}\) is properly infinite.
For every \(x\in\Id(h)_+\) and every integer \(k\geq1\),
\begin{equation}
 x^{\oplus k}\precsim h^{\oplus n}.
 \label{eq:ideal-multiplicity}
\end{equation}
In particular, if \(C\) satisfies \propF{}, this conclusion holds for
every \(h\in C_+\) with the same \(n\).
\end{lemma}

\begin{proof}
In the common stabilization,
\[
 \Id_{C\otimes\K}(h^{\oplus n})
   =\Id_C(h)\otimes\K.
\]
Thus \(x^{\oplus k}\) belongs to the ideal generated by
\(h^{\oplus n}\).  The proper-target ideal comparison in
Lemma~\ref{lem:comparison-rules} gives
\eqref{eq:ideal-multiplicity}.  This argument takes place in the
stabilization and therefore covers every finite rectangular corner
and nonunital \(C\).  It also covers \(h=0\), when \(x=0\).
\end{proof}

\begin{lemma}[Exact implementation of a spectral cut]
\label{lem:cut-implementation}
Let \(C\) be a \(C^*\)-algebra, let \(m,s\geq1\), and suppose that
\(x\in M_m(C)_+\) and \(h\in M_s(C)_+\) satisfy \(x\precsim h\).
If \(f\colon[0,\|x\|]\to[0,\infty)\) is continuous and vanishes
on a neighborhood of zero, then there is \(R\in M_{s,m}(C)\) such that
\begin{equation}
 R^*R=f(x),\qquad RR^*\in\Her(h),\qquad
 \|R\|^2=\|f(x)\|.
 \label{eq:exact-cut}
\end{equation}
\end{lemma}

\begin{proof}
If \(f(x)=0\), take \(R=0\).  Otherwise, choose \(\delta>0\) so
that \(f\) vanishes on \([0,\delta]\), and choose
\(0<\varepsilon<\delta\).  Then
\(f(x)^{1/2}\in\Her((x-\varepsilon)_+)\).
Apply Lemma~\ref{lem:transport} to \(x\precsim h\) at this cutoff,
and let \(w\in M_{s,m}(C^{**})\) be its rectangular partial isometry.
Its initial projection is the support of \((x-\varepsilon)_+\).
The membership and support conclusions of that lemma therefore give
\[
 R=w f(x)^{1/2}\in M_{s,m}(C),\qquad
 R^*R=f(x),\qquad RR^*=w f(x)w^*\in\Her(h).
\]
The norm identity follows from \(\|R\|^2=\|R^*R\|\).
Only the product \(R\), rather than \(w\) itself, is asserted to
belong to the rectangular algebra.
\end{proof}

For \(P\in\Prim(C)\), a subscript \(P\) denotes the image in \(C/P\),
also at matrix levels.  An irreducible representation with kernel
\(P\) faithfully represents this quotient and computes its norm.
The primitive ideal test in Lemma~\ref{lem:ideals-cutoffs} will let us
recognize ideal membership from these quotient images.

The uniform gap will be applied to the unitization of a matrix algebra
over \(C\).  We record the two permanence facts that make this a valid
exact target, also when \(C\) has no unit.

\begin{lemma}[The exact target for matrix models]
\label{lem:exact-matrix-unitization}
If \(A\) is an exact \(C^*\)-algebra and \(n\ge1\), then
\(\widetilde{M_n(A)}\) is exact.
\end{lemma}

\begin{proof}
Finite matrix exactness follows, for every \(C^*\)-algebra \(D_0\),
from the identifications
\[
 M_n(A)\otimes_{\min}D_0=M_n(A\otimes_{\min}D_0),
\]
which preserve the ideal kernels in the definition of exactness.
For the unitization step, let \(A_0\) be a nonunital exact algebra,
let \(I\triangleleft D_0\) be a closed two-sided ideal, and write
\(q_0\colon D_0\to D_0/I\).  The scalar character
\(\chi\colon\widetilde A_0\to\C\) gives, for every
\(z\in\widetilde A_0\otimes_{\min}D_0\), the bounded decomposition
\[
 z=z_0+1\otimes d,\qquad
 d=(\chi\otimes\id)(z),\quad z_0\in A_0\otimes_{\min}D_0.
\]
Indeed, both projections are bounded, and the decomposition holds
first on algebraic tensors and then on their completion.  If
\((\id\otimes q_0)(z)=0\), applying \(\chi\otimes\id\) gives
\(d\in I\).  The isometric inclusion
\[
 A_0\otimes_{\min}(D_0/I)
 \subset\widetilde A_0\otimes_{\min}(D_0/I)
\]
then shows \((\id\otimes q_0)(z_0)=0\).  Exactness of \(A_0\) gives
\(z_0\in A_0\otimes_{\min}I\), hence
\(z\in\widetilde A_0\otimes_{\min}I\).  If \(M_n(A)\) is already
unital, this second step is unnecessary.
\end{proof}

Fix, once and for all, a constant
\begin{equation}
 0<\sigma<\min\{1,\gamma/3\},
 \label{eq:perturbation-constant}
\end{equation}
where \(\gamma\) is supplied by Proposition~\ref{prop:uniform-gap}.
This \(\sigma\) is independent of the algebra and of the fixed
amplification integer.

\begin{proposition}[Uniform perturbation]\label{prop:perturbation}
Let \(C\) be an exact \(C^*\)-algebra satisfying \propF{} for one
fixed integer \(n\geq1\).  Let \(p,q\in C\) be contractions and
let \(h,k,c\in C_+\) satisfy
\begin{equation}
 p^*ph=h,\qquad qq^*k=k,\qquad
 c^{\oplus n}\precsim h,\qquad c^{\oplus n}\precsim k.
 \label{eq:perturbation-hypotheses}
\end{equation}
For every \(T\in C\) and every \(\eta,\alpha>0\), there is a
contraction \(y\in C\) such that
\begin{equation}
 \|(T+\eta pyq)_P\|\geq\sigma\eta
 \quad\text{for every }P\in\Prim(C)\text{ with }\|c_P\|\geq\alpha.
 \label{eq:perturbation-conclusion}
\end{equation}
If \(h,k\) belong to a hereditary subalgebra
\(\mathcal D\subset C\), then \(y\) may be chosen in \(\mathcal D\).
If, in addition or independently, \(p,q\) are positive and \(pq=0\),
then \(y\) may be chosen selfadjoint.  Both conclusions may be
imposed simultaneously.  The constant \(\sigma\) is independent
of all the data in the Proposition, including \(n\) and the norms
of \(T,h,k,c\).
\end{proposition}

\begin{proof}
Write \(K=\{P\in\Prim(C):\|c_P\|\geq\alpha\}\).  If \(K\) is
empty, \(y=0\) satisfies all conclusions.  Suppose henceforth that
\(K\) is nonempty.  The contradiction will have three levels of choices.
First we fix comparison rows and a single contraction in \(M_n(C)\).  Each finite matrix contraction then determines an
admissible perturbation.  If it fails, a primitive quotient witnessing
that failure supplies an approximate reducing copy for that fixed
contraction.  All these copies therefore come from one element
of the same exact algebra \(\widetilde{M_n(C)}\).

\smallskip
\noindent\emph{Fixed comparison rows.}
Choose a continuous function \(e\colon[0,\infty)\to[0,1]\) that
vanishes near zero and is one on \([\alpha/2,\infty)\).
Apply Lemma~\ref{lem:cut-implementation} to the last two comparisons
in \eqref{eq:perturbation-hypotheses}.  This gives rows
\(R,S\in M_{1,n}(C)\) with
\begin{equation}
 R^*R=S^*S=E=e(c)^{\oplus n},\qquad
 RR^*\in\Her(h),\quad SS^*\in\Her(k),\qquad
 \|R\|,\|S\|\leq1.
 \label{eq:perturbation-rows}
\end{equation}
The relation \(p^*ph=h\), together with its adjoint, shows that
left and right multiplication by \(p^*p\) are the identity on
\(\Her(h)\).  Hence
\((1-p^*p)RR^*(1-p^*p)=0\), which implies
\((1-p^*p)R=0\).  The same argument for \(S\) gives
\begin{equation}
 p^*pR=R,\qquad qq^*S=S.
 \label{eq:perturbation-row-identities}
\end{equation}
At this point fix the single element
\begin{equation}
 V=\eta^{-1}R^*p^*Tq^*S\in M_n(C).
 \label{eq:fixed-testing-element}
\end{equation}
The rows \(R,S\) and the element \(V\) will not depend on any finite
matrix model.  Fix now also the global contraction
\[
 \widehat V=V(\max\{1,|V|\})^{-1}\in M_n(C).
\]
The inverse is evaluated in the unitization, so its product with
\(V\) remains in \(M_n(C)\).

\smallskip
\noindent\emph{Reducing copies of arbitrary finite contractions.}
Let \(d\geq1\) and let \(D\in M_d(\C)\) be a contraction.
Put \(c_1=(c-\alpha/2)_+\).  Condition \propF{} and
Lemma~\ref{lem:ideal-comparison}, applied to the scalar element
\(c_1\), give
\[
 c_1^{\oplus d}\precsim c_1^{\oplus n}.
\]
Choose a continuous \(\varphi\colon[0,\infty)\to[0,1]\), vanishing
near zero and equal to one on \([\alpha/4,\infty)\), and put
\(e_1(c)=\varphi(c_1)\).  Lemma~\ref{lem:cut-implementation}
yields \(L\in M_{n,d}(C)\) such that
\begin{equation}
 L^*L=e_1(c)^{\oplus d},\qquad
 LL^*\in\Her(c_1^{\oplus n}),\qquad \|L\|\leq1.
 \label{eq:finite-model-column}
\end{equation}
Because \(e(c)c_1=c_1\), multiplication by \(E\) is the identity
on \(\Her(c_1^{\oplus n})\).  The range-square argument used for
\eqref{eq:perturbation-row-identities} gives \(EL=L\).
Define
\begin{equation}
 w=LDL^*\in M_n(C),\qquad \|w\|\leq1,\qquad EwE=w.
 \label{eq:finite-model}
\end{equation}
Scalar matrix multiplication here requires no unit in \(C\).

For each \(P\in K\), choose an irreducible representation
\(\pi_P\colon C\to B(H_P)\) with kernel \(P\), and use its matrix
amplifications.  Identify \(C/P\) with the image of the induced
faithful representation on \(H_P\), and make the same identification
at rectangular matrix levels.  The spectral subspace
\(1_{(3\alpha/4,\infty)}(\pi_P(c))H_P\) is nonzero, since
\(\|c_P\|\geq\alpha\).  Choose a unit vector \(\xi_P\) in that
subspace, so that \(\pi_P(e_1(c))\xi_P=\xi_P\).  Define
\[
 J_P\colon\C^d\longrightarrow H_P^{\oplus n},\qquad
 J_P\lambda=L_P(\lambda\otimes\xi_P).
\]
Equation~\eqref{eq:finite-model-column} proves \(J_P^*J_P=1_d\).
Both \(D\) and \(D^*\) preserve \(\C^d\otimes\xi_P\), and the
same equation gives
\begin{equation}
 w_PJ_P=J_PD,\qquad w_P^*J_P=J_PD^*.
 \label{eq:reducing-model}
\end{equation}
Thus one copy reduces both \(w_P\) and its adjoint.  No eigenvector
at the norm of \(c_P\), and no continuous choice of the vectors
\(\xi_P\), is required.

\smallskip
\noindent\emph{Admissible tests and the two optional conclusions.}
For the chosen \(d,D\), take
\[
 y_0=-RwS^*\in C,\qquad\|y_0\|\leq1.
\]
If \(h,k\in\mathcal D\), then \(RR^*,SS^*\in\mathcal D\).
For a positive contractive approximate unit \((a_\lambda)\) of
\(\mathcal D\), the identity
\[
 \|(1-a_\lambda)R\|^2
 =\|(1-a_\lambda)RR^*(1-a_\lambda)\|\longrightarrow0
\]
and its analogue for \(S\) give
\(a_\lambda RwS^*a_\lambda\to RwS^*\) in norm.
Each approximant belongs to
\(\mathcal D C\mathcal D\subset\mathcal D\), so \(y_0\in\mathcal D\).
This uses the range squares of the rows and makes no assertion that
their individual entries belong to \(\mathcal D\).

If \(p,q\geq0\) and \(pq=0\), then \(p^2R=R\) implies \(pR=R\):
multiply \((1-p^2)R=0\) by \((1+p)^{-1}\) in the unitization.
Likewise \(qS=S\).  It follows that \(pS=qR=0\) and \(S^*R=0\).
Consequently \(y_0^2=0\) and \(py_0^*q=0\).  The positive elements
\(y_0y_0^*\) and \(y_0^*y_0\) are orthogonal, whence
\[
 \|y_0+y_0^*\|^2
 =\|y_0y_0^*+y_0^*y_0\|=\|y_0\|^2\leq1.
\]
We may therefore use \(y=y_0+y_0^*\) when selfadjointness is
requested, and \(y=y_0\) otherwise.  In both cases \(pyq=py_0q\);
hereditary membership is preserved if it is also requested.

\smallskip
\noindent\emph{One global contraction and the uniform gap.}
Suppose that no admissible \(y\) satisfies
\eqref{eq:perturbation-conclusion}.  For each \(d,D\), choose a
witness \(P\in K\) for the test just constructed.  Then
\[
 \|(T-\eta pRwS^*q)_P\|<\sigma\eta.
\]
Multiplying by the contractive outer factors \(R^*p^*\) and \(q^*S\),
and using \eqref{eq:perturbation-row-identities}, gives
\begin{equation}
 \|(V-EwE)_P\|=\|(V-w)_P\|<\sigma.
 \label{eq:local-near-model}
\end{equation}
Although \(P\) may depend on \(d,D\), the element \(V\) is the one
fixed in \eqref{eq:fixed-testing-element}.

Functional calculus commutes with quotient maps for the fixed global
clipping \(\widehat V\).  Hence in each witness quotient
\begin{equation}
 \|V_P-\widehat V_P\|
 =\|(|V_P|-1)_+\|
 =(\|V_P\|-1)_+<\sigma.
 \label{eq:quotient-clipping}
\end{equation}
The last strict inequality follows from \eqref{eq:local-near-model}
and \(\|w_P\|\leq1\).  The clipping identity follows by squaring the
absolute value of
\(V_P-V_P\max\{1,|V_P|\}^{-1}\) and applying scalar functional
calculus.  Combining \eqref{eq:local-near-model},
\eqref{eq:quotient-clipping}, and \eqref{eq:reducing-model} gives
\begin{equation}
 \|\widehat V_PJ_P-J_PD\|<2\sigma<\gamma,\qquad
 \|\widehat V_P^*J_P-J_PD^*\|<2\sigma<\gamma.
 \label{eq:forbidden-intertwining}
\end{equation}

By Lemma~\ref{lem:exact-matrix-unitization},
\(B=\widetilde{M_n(C)}\) is a unital exact algebra.
The amplified representations \(\pi_P^{(n)}\) extend unitally to
\(B\).  Although the witnessing \(P\) varies with \(d,D\), the
contraction \(\widehat V\) and the exact algebra \(B\) remain fixed.
Thus \(\widehat V\in B\) has, for every \(d,D\), a unital
representation and an isometry satisfying
both inequalities in \eqref{eq:forbidden-intertwining}.  This
contradicts Proposition~\ref{prop:uniform-gap} and proves the
Proposition with both optional conclusions and the same \(\sigma\).
\end{proof}

\section{Exact localization and removal of amplification}
\label{sec:localization}

We now use Proposition~\ref{prop:perturbation} to prove the fixed-amplification
part of Theorem~\ref{thm:exact-infiniteness}.  Throughout this section,
let $C$ be an exact $C^*$-algebra satisfying \propF{} for one fixed
integer $n\geq1$, and retain the constant $0<\sigma<1$ from
\eqref{eq:perturbation-constant}.  For $P\in\Prim(C)$, write
$a_P=a+P$, with its norm computed in any irreducible representation
having kernel $P$.  A hereditary subalgebra $\mathcal D\subset C$
is nonzero at $P$ when $\mathcal D\not\subset P$.

Our immediate objective is to construct a square-zero element in a
hereditary subalgebra that is nonzero on a prescribed quasi-compact
subset of $\Prim(C)$.  Its initial and range squares are orthogonal
and Cuntz equivalent.  Repeating the construction on one fixed norm
slice will remove the amplification integer $n$ and give \propP{}.
Theorem~\ref{thm:individual-to-strong} will then supply the strong
pure infiniteness conclusion.  All finite covers in the construction
concern the primitive ideal space; the partition of unity used below
is on a real interval.

\subsection{Local elements and finite covers}

The next lemma uses only \propF{}: finite-dimensional hereditary
images are excluded, and irreducibility then supplies common positive
elements locally.  Exactness enters later through
Proposition~\ref{prop:perturbation}, when these local choices must
work simultaneously on a whole quasi-compact set.

\begin{lemma}[Local common elements]\label{lem:local-common-element}
Let $C$ be an exact $C^*$-algebra satisfying \propF{} for one fixed
integer $n\geq1$, and let $P\in\Prim(C)$.
\begin{enumerate}[label=\textup{(\roman*)}]
\item If $h\in C_+$ and $h_P\ne0$, then for every integer $m\ge1$
there are $m$ pairwise orthogonal elements of $\Her(h)_+$, each
nonzero at $P$.
\item If $h,k\in C_+$ are nonzero at $P$, there is $c\in C_+$
such that $c_P\ne0$ and
\[
 c^{\oplus n}\precsim h,\qquad c^{\oplus n}\precsim k.
\]
If $h,k$ belong to a hereditary subalgebra $\mathcal D$, $c$ can also
be chosen in $\mathcal D$.
\end{enumerate}
\end{lemma}

\begin{proof}
Fix an irreducible representation $\pi:C\to B(H_\pi)$ with kernel
$P$.  First, no nonzero $\pi(a)$, $a\in C_+$, has finite rank.
Indeed, if its rank were $d\in\N$, applying the amplified
representation to
$a^{\oplus 2n}\precsim a^{\oplus n}$ would express an operator of
rank $2nd$ as a norm limit of operators of rank at most $nd$.
This is impossible, since the set of operators of rank at most a
fixed finite integer is norm closed.

Put $H_h=\overline{\pi(h)H_\pi}$.  We claim that the restriction of
$\pi(\Her(h))$ to $H_h$ is strongly dense in $B(H_h)$.
Irreducibility and the bicommutant theorem
\cite[I.9.1.1]{Blackadar} give strong density of $\pi(C)$ in $B(H_\pi)$.  Consequently the strong closure of
$\pi(\Her(h))$ contains $\pi(h)B(H_\pi)\pi(h)$.
For $t>0$, let $E_t=1_{[t,\infty)}(\pi(h))$.  Every operator
$E_tTE_t$, with $T\in B(H_h)$ extended by zero on $H_h^\perp$, has
the form $\pi(h)S\pi(h)$: the inverse of $\pi(h)$ on $E_tH_\pi$
is bounded.  As $t\downarrow0$, $E_tTE_t\to T$ strongly on $H_h$.
This proves the claim.

The space $H_h$ is infinite dimensional by the finite-rank
obstruction.  Choose orthonormal vectors $\xi_1,\ldots,\xi_m$ in
$H_h$ and distinct numbers $\lambda_1,\ldots,\lambda_m\in(0,1)$.
Let $T\in B(H_h)$ be the positive contraction with
$T\xi_j=\lambda_j\xi_j$ and zero on their orthogonal complement.
Choose nonnegative continuous functions $f_j$ on $\R$, with disjoint
compact supports in $(0,1)$ and $f_j(\lambda_j)=1$.
Kaplansky density \cite[I.9.1.3]{Blackadar} supplies a bounded net of selfadjoint elements
$b_\lambda$ of the represented hereditary algebra converging
strongly to $T$.  Continuous functional calculus on bounded
selfadjoint sets is strongly continuous; this follows by uniform
polynomial approximation on a common bounded interval.  Hence
$f_j(b_\lambda)\xi_j\to\xi_j$.  One $b_\lambda$ makes all these
$m$ vectors nonzero.  Lift it to a selfadjoint
$b\in\Her(h)$.  The elements $f_j(b)$ lie in $\Her(h)$, are
pairwise orthogonal, and are nonzero at $P$.  The functions vanish
at zero, so this argument also applies to a nonunital hereditary
algebra.  This proves \textup{(i)}.

For two positives $u,v$ nonzero at $P$, strong density of $\pi(C)$
gives $x\in C$ with $\pi(uxv)\ne0$.  For example, choose a vector
on which $\pi(v)$ is nonzero and a rank-one operator carrying its
image to a vector not annihilated by $\pi(u)$; approximate that
operator strongly on the image under $\pi(v)$.  With $z=uxv$, the positive
$d=zz^*$ satisfies
\[
 d_P\ne0,\qquad d\in\Her(u),\qquad
 d\CuEq z^*z\in\Her(v).
\]
Here and below $\CuEq$ denotes Cuntz equivalence, and the comparison
rules are those of Lemma~\ref{lem:comparison-rules}.
For a finite list of positives nonzero at $P$, repeating this
operation produces one positive, nonzero at $P$, subequivalent to
every element of the list.  At each step the new element lies in
the hereditary algebra of the preceding one, so all previous
comparisons and any prescribed hereditary membership persist.

Apply \textup{(i)} with $m=n$ separately to $h$ and $k$, obtaining
orthogonal lists $h_1,\ldots,h_n$ and $k_1,\ldots,k_n$.
Take a common element $c$ for the resulting list of $2n$ positives.
Then
\[
 c^{\oplus n}\precsim\bigoplus_{j=1}^n h_j
 \CuEq\sum_{j=1}^n h_j\precsim h,
\]
and the same argument gives $c^{\oplus n}\precsim k$.
The asserted choice in $\mathcal D$ follows from the hereditary
membership just observed.
\end{proof}

We record the norm-attainment fact that ensures a strict inequality
in the finite-cover argument.

\begin{lemma}[Pure-state norm attainment]\label{lem:purestateattainment}
If $B$ is a $C^*$-algebra and $0\ne b\in B_+$, there is an
irreducible representation $\pi$ of $B$ with
$\|\pi(b)\|=\|b\|$.
\end{lemma}

\begin{proof}
In the state space of the unitization $\widetilde B$, the set
\[
 F_b=\{\varphi:\varphi(b)=\|b\|\}
\]
is nonempty: evaluation at the largest spectral value of $b$ on
$C^*(1,b)$ extends to a state of $\widetilde B$ by the positive
Hahn--Banach extension theorem \cite[II.6.3.1]{Blackadar}.  It is a weak-star compact face,
since every state has value at most $\|b\|$ on $b$.  By the
Krein--Milman theorem it has an extreme point, which is therefore
a pure state of $\widetilde B$; see \cite[II.6.3.2--3]{Blackadar}.
Let $\widetilde\pi$ be its irreducible GNS representation
\cite[II.6.4.8]{Blackadar}.  It does not annihilate $B$, because
its value on $b$ is positive.  The closure of
$\widetilde\pi(B)H_{\widetilde\pi}$ is a nonzero reducing subspace
for $\widetilde\pi(\widetilde B)$, hence the whole space.
If an operator commutes with $\widetilde\pi(B)$, then its commutator
with $\widetilde\pi(x)$ annihilates $\widetilde\pi(B)H_{\widetilde\pi}$
for every $x\in\widetilde B$, since $xB\subset B$.
It therefore commutes with $\widetilde\pi(\widetilde B)$.
Thus the restriction $\pi=\widetilde\pi|_B$ is nondegenerate and
irreducible.  The vector state gives
$\|\pi(b)\|\ge\|b\|$, and contractivity gives equality.
\end{proof}

\begin{lemma}[Primitive norm slices]\label{lem:primitive-compactness}
Let $C$ be a $C^*$-algebra.  For $c\in C_+$ and $\alpha>0$, the set
\[
 \Omega(c,\alpha)=\{P\in\Prim(C):\|c_P\|>\alpha\}
\]
is open, and
\[
 K(c,\alpha)=\{P\in\Prim(C):\|c_P\|\ge\alpha\}
\]
is quasi-compact.  No Hausdorff hypothesis is required.
\end{lemma}

\begin{proof}
The first set is $\{P:(c-\alpha)_+\notin P\}$, which is open in
the hull-kernel topology.  For the second assertion, extend any
relative open cover to ambient open sets.  Write those sets as
\[
 \mathcal O_\lambda=\{P:J_\lambda\not\subset P\},
 \qquad I=\overline{\sum_\lambda J_\lambda},
\]
where each $J_\lambda$ is a closed ideal of $C$.
If $\|c+I\|\ge\alpha$, Lemma~\ref{lem:purestateattainment}
applied in $C/I$ gives an irreducible representation attaining that
norm.  Its pullback has a primitive kernel containing $I$, belongs
to $K(c,\alpha)$, and lies outside every $\mathcal O_\lambda$.
This contradicts the covering property.  Thus
$\|c+I\|<\alpha$; if $C/I=0$, the same inequality is immediate.
Choose $d\in I$ with $\|c-d\|<\alpha$, and then approximate $d$
by an element
$d_F\in\sum_{\lambda\in F}J_\lambda$ for a finite index set
$F$, closely enough that $\|c-d_F\|<\alpha$ still holds.
Any primitive ideal containing every $J_\lambda$, $\lambda\in F$,
has $\|c_P\|<\alpha$.  Hence the finitely many
$\mathcal O_\lambda$, $\lambda\in F$, cover $K(c,\alpha)$.
The proof uses quasi-compactness only and does not assert that
$K(c,\alpha)$ is closed.
\end{proof}

\begin{lemma}[Finite coupling]\label{lem:finite-coupling}
Let $C$ be an exact $C^*$-algebra satisfying \propF{} for one fixed
integer $n\geq1$.  Let $K_0\subset\Prim(C)$ be quasi-compact and let
$\mathcal D\subset C$ be hereditary.  Suppose that at each
$P\in K_0$ one can choose positive contractions $p,q\in\mathcal D$
and $h,k\in\mathcal D_+$ satisfying
\[
 pq=0,\qquad ph=h,\qquad qk=k,\qquad h_P\ne0,\quad k_P\ne0.
\]
There are finitely many such choices $(p_i,q_i,h_i,k_i)$ and one
selfadjoint $U\in\mathcal D$ such that, at every $P\in K_0$,
\[
 (p_iUq_i)_P\ne0
\]
for at least one selected index $i$.  Any support conditions in
the original choices are retained in the selected ones.
\end{lemma}

\begin{proof}
If $K_0$ is empty, choose $U=0$ and the empty family.  Otherwise,
Lemma~\ref{lem:local-common-element} associates to every original
choice a positive $c$ nonzero at its specified point, with
$c^{\oplus n}\precsim h,k$.  Choose
$0<\alpha<\|c_P\|$.  The resulting open sets
$\Omega(c,\alpha)$ cover $K_0$, so choose a finite subcover
$\Omega_i=\Omega(c_i,\alpha_i)$, $1\le i\le N$.

Set $\theta=\sigma/4$, $\eta_i=\theta^{i-1}$, and $U_0=0$.
At stage $i$, apply Proposition~\ref{prop:perturbation} with
$T=p_iU_{i-1}q_i$, scale $\eta_i$, and the selected data.
The identities $p_ih_i=h_i$, $q_ik_i=k_i$ imply
$p_i^2h_i=h_i$, $q_i^2k_i=k_i$.  The simultaneous hereditary and
selfadjoint options give a selfadjoint contraction
$y_i\in\mathcal D$ such that $U_i=U_{i-1}+\eta_i y_i$ satisfies
\[
 \|(p_iU_iq_i)_P\|\ge\sigma\eta_i
 \qquad(P\in K(c_i,\alpha_i)).
\]
Let $U=U_N$.  Later perturbations change the $i$th corner by at
most
\[
 \sum_{j>i}\eta_j
 \le\eta_i\frac{\theta}{1-\theta}
 =\eta_i\frac{\sigma}{4-\sigma}
 <\frac{\sigma\eta_i}{2}.
\]
Consequently
\begin{equation}\label{eq:finite-coupling-bound}
 \|(p_iUq_i)_P\|>\frac{\sigma\eta_i}{2}
 \qquad(P\in K(c_i,\alpha_i)).
\end{equation}
In particular each selected corner is nonzero throughout its
entire open set $\Omega_i$, proving the assertion.
\end{proof}

When applying Lemma~\ref{lem:finite-coupling}, it suffices to
produce two orthogonal positives $u,v$ nonzero at the specified
point.  Indeed, choose $0<\beta<\|u_P\|$ and a continuous
$f:[0,\infty)\to[0,1]$ that is zero near zero and one on
$[\beta,\infty)$.  Then $p=f(u)$ and $h=(u-\beta)_+$ satisfy
$ph=h$ and $h_P\ne0$.  The analogous construction from $v$ gives
$q,k$.  All four elements retain the hereditary and spectral
support restrictions of $u,v$.

\subsection{A square-zero element on a quasi-compact set}

\begin{proposition}\label{prop:square-zero}
Let $C$ be an exact $C^*$-algebra satisfying \propF{} for one fixed
integer $n\geq1$.  Let $K_0\subset\Prim(C)$ be quasi-compact, and let
$\mathcal D\subset C$ be a hereditary subalgebra nonzero at every
point of $K_0$.  There exists $x\in\mathcal D$ such that
\[
 x^2=0,\qquad x_P\ne0\quad(P\in K_0).
\]
\end{proposition}

\begin{proof}
If $K_0=\varnothing$, take $x=0$.  Assume henceforth that $K_0$
is nonempty.  We first construct a selfadjoint element that is
nonscalar on every relevant hereditary support space.  Coupling its
separated spectral windows and deleting the nearby spectral blocks
will produce a selfadjoint element whose positive and negative parts
are both nonzero at every point of $K_0$.  A final coupling of those
two orthogonal parts gives the required square-zero element.

\paragraph{First coupling: a nonscalar selfadjoint element.}
At each $P\in K_0$, choose $h\in\mathcal D_+$ with $h_P\ne0$.
Lemma~\ref{lem:local-common-element}\textup{(i)}, with $m=2$,
provides two orthogonal positives in $\Her(h)\subset\mathcal D$
nonzero at $P$.  Apply the preceding cutoff construction to these
two positives and use the resulting four inputs in
Lemma~\ref{lem:finite-coupling} to obtain $U=U^*\in\mathcal D$.
For an irreducible representation $\pi$ with kernel $P$, put
$H_{\mathcal D,\pi}=\overline{\pi(\mathcal D)H_\pi}$.
This space reduces $\pi(\mathcal D)$, which acts as zero on its
orthogonal complement.  At $P$ one selected $pUq$ is nonzero.
Since $p,q\in\mathcal D$ and $pq=0$, this implies
\begin{equation}\label{eq:nonscalarU}
 \pi(U)|_{H_{\mathcal D,\pi}}\text{ is not a scalar operator}
 \qquad(P\in K_0).
\end{equation}
In particular, nonscalarity is obtained on the support space of
$\mathcal D$, without using its zero action on the complement.

\paragraph{Second coupling: separated spectral windows.}
At a specified $P\in K_0$, the selfadjoint operator in
\eqref{eq:nonscalarU} has two distinct spectral values.  Choose
disjoint closed intervals $I,J\subset\R$ at positive distance,
and continuous nonnegative functions $F,G$ supported in $I,J$,
respectively, such that both
$F(\pi(U))|_{H_{\mathcal D,\pi}}$ and
$G(\pi(U))|_{H_{\mathcal D,\pi}}$ are nonzero.
The functions may be chosen on all of $\R$ with compact support.
If a window contains zero, its function need not vanish at zero.
We evaluate scalar functions in $\mathcal D+\C1$ inside a fixed
unitization $\widetilde C$ of $C$, and take spectral projections in
$\widetilde C^{**}$.
Let $(d_\lambda)$ be a positive contractive approximate identity of
$\mathcal D$.
The positive elements
\[
 F(U)d_\lambda F(U),\qquad G(U)d_\lambda G(U)
\]
belong to $\mathcal D$.  Their images on
$H_{\mathcal D,\pi}$ converge strongly to $F(\pi(U))^2$ and
$G(\pi(U))^2$, respectively, so a choice of middle elements makes
both images nonzero.  These positives are supported, globally in
the bidual, on $1_I(U)$ and $1_J(U)$.

The same cutoff construction gives the corresponding $p,q,h,k$,
retaining these spectral supports.  Apply
Lemma~\ref{lem:finite-coupling} again to obtain
$V=V^*\in\mathcal D$ and finitely many selected pairs
$(p_i,q_i)$, supported in windows $(I_i,J_i)$, such that
\begin{equation}\label{eq:spectralwindowV}
 \begin{gathered}
 p_i=1_{I_i}(U)p_i1_{I_i}(U),\qquad
 q_i=1_{J_i}(U)q_i1_{J_i}(U),\\
 \text{for every }P\in K_0\text{ there is }i
 \text{ with }(p_iVq_i)_P\ne0.
 \end{gathered}
\end{equation}
Put $\rho=\min_i\dist(I_i,J_i)>0$.
The positive elements $p_i,q_i$ need not commute with $U$;
only the displayed support identities will be used.

\paragraph{Spectral deletion produces both signs.}
We retain the corners in \eqref{eq:spectralwindowV} while removing
all sufficiently narrow diagonal spectral blocks.  A positive operator
with all those diagonal blocks zero must vanish; the same holds for
a negative operator.  Since the retained corners prevent vanishing,
this will force both signs in every quotient over $K_0$.

Fix a compact interval $[-M,M]$ containing the spectrum of $U$
in the unitization.  Take a finite continuous nonnegative
partition of unity $f_1,\ldots,f_s$ on this interval, whose
supports $S_j$ have diameter less than $\rho/8$.
Define
\begin{equation}\label{eq:deletionH}
 H=\sum_{\dist(S_j,S_k)>\rho/2}
       f_j(U)Vf_k(U).
\end{equation}
The summation uses ordered pairs, and its condition is symmetric
in $j,k$, so $H=H^*$.  Although a factor $f_j(U)$ may have a
scalar part, it lies in $\mathcal D+\C1$; the middle
factor $V\in\mathcal D$ implies $H\in\mathcal D$.

For each selected pair $(p_i,q_i)$ we have
\begin{equation}\label{eq:window-preservation}
 p_iHq_i=p_iVq_i.
\end{equation}
Indeed, expand $p_iVq_i$ using $\sum_jf_j(U)=1$ on both sides.
For a summand $p_if_j(U)Vf_k(U)q_i$ to be nonzero, the supports
$S_j,S_k$ must meet $I_i,J_i$, respectively, by
\eqref{eq:spectralwindowV}.  Their distances then satisfy
\[
 \dist(S_j,S_k)
 \ge\dist(I_i,J_i)
       -\operatorname{diam}(S_j)-\operatorname{diam}(S_k)
 >\frac{3\rho}{4}>\frac\rho2.
\]
Every potentially nonzero summand is therefore retained in
\eqref{eq:deletionH}.  This argument requires no commutation of
$p_i,q_i$ with $U$.  Equations~\eqref{eq:spectralwindowV} and
\eqref{eq:window-preservation} show that $H_P\ne0$ for every
$P\in K_0$.

Let $B$ be any Borel interval of length less than $\rho/4$.
The two supports in a retained pair cannot both meet $B$, since
their distance exceeds $\rho/2$.  Hence
\[
 1_B(U)H1_B(U)=0.
\]
Partition $[-M,M]$ into finitely many disjoint Borel intervals
$B_j$ of length less than $\rho/4$, assigning endpoints so that
their spectral projections sum to the identity.  In any unital
extension of a representation, the resulting projections $E_j$
thus satisfy $E_j\pi(H)E_j=0$ and $\sum_jE_j=1$.
If $\pi(H)\ge0$, then
\[
 \|\pi(H)^{1/2}E_j\xi\|^2
 =\langle E_j\pi(H)E_j\xi,\xi\rangle=0
\]
for every $j,\xi$, whence $\pi(H)=0$.  Applying the same argument
to $-\pi(H)$ excludes a nonzero negative semidefinite image.
The spectral projection calculations can equivalently be made
in the bidual and passed through the normal extensions of the
representations.  We conclude that
\begin{equation}\label{eq:bothsigns}
 (H_+)_P\ne0,\qquad (H_-)_P\ne0
 \qquad(P\in K_0).
\end{equation}
Both positive parts belong to $\mathcal D$ and are orthogonal.

\paragraph{Third coupling: the square-zero product.}
At every $P\in K_0$, apply the same cutoff construction to $H_+,H_-$
to obtain $p,h\in\Her(H_+)$ and $q,k\in\Her(H_-)$ for
Lemma~\ref{lem:finite-coupling}.  The contractions $p,q$ are
orthogonal.  The resulting
$U'=U'^*\in\mathcal D$ has a selected nonzero corner at every
point of $K_0$.  Put
\[
 x=H_+U'H_-\in\mathcal D.
\]
If $x_P=0$, then $(pU'q)_P=0$ for every
$p\in\Her(H_+)$ and $q\in\Her(H_-)$: approximate them in
norm by elements $H_+aH_+$ and $H_-bH_-$, and use
\[
 (H_+aH_+)U'(H_-bH_-)=H_+a\,x\,bH_-.
\]
This contradicts the selected nonzero corner at $P$.
Thus $x_P\ne0$ throughout $K_0$, while
$x^2=H_+U'(H_-H_+)U'H_-=0$.
All constructed elements lie in $\mathcal D$; spectral
projections have only been used to verify their supports.
\end{proof}

\subsection{Removing the amplification}

We have constructed a square-zero element on every prescribed
quasi-compact set where a hereditary subalgebra is nonzero.  We now
apply this result repeatedly on the same primitive norm slice.  The
resulting number of orthogonal equivalent copies can be made at least
the single amplification integer in \propF{}.

\begin{proof}[Proof of Theorem~\ref{thm:exact-infiniteness}]
Fix $a\in C_+$ and $r>0$.  The set $K_0=K(a,r)$ is quasi-compact by
Lemma~\ref{lem:primitive-compactness}.  If it is empty, every primitive
image of $(a-r)_+$ is zero, so $(a-r)_+=0$, and
$((a-r)_+)^{\oplus2}\precsim a$ is immediate.
Suppose that $K_0$ is nonempty.  Proposition~\ref{prop:square-zero}
applied to $\Her(a)$ gives $x\in\Her(a)$ with $x^2=0$ and
$x_P\ne0$ throughout $K_0$.  Set $z_1=x^*x$.  Its Cuntz-equivalent
positive $xx^*$ is orthogonal to it, since
$(x^*x)(xx^*)=x^*x^2x^*=0$.  Therefore
\[
 z_1^{\oplus2}\CuEq z_1\oplus xx^*
 \CuEq z_1+xx^*\precsim a,
 \qquad (z_1)_P\ne0\quad(P\in K_0).
\]
Repeat the construction inside $\Her(z_1)$, and then inside successive
hereditary subalgebras, always over the same $K_0$.  It yields
$z_j\in C_+$ nonzero throughout $K_0$ and
$z_{j+1}^{\oplus2}\precsim z_j$.  Finite direct sums and transitivity
give
\[
 z_j^{\oplus 2^j}\precsim a.
\]
Choose a finite $j\ge1$ with $2^j\ge n$.  Dropping extra summands gives
$z_j^{\oplus n}\precsim a$.  If $(z_j)_P=0$, then $P\notin K_0$, so
$\|a_P\|<r$ and $((a-r)_+)_P=0$.  The primitive ideal test in
Lemma~\ref{lem:ideals-cutoffs} therefore gives
$(a-r)_+\in\Id(z_j)$.  Condition \propF{} makes $z_j^{\oplus n}$ properly infinite because
$z_j$ is a scalar positive element.
Lemma~\ref{lem:ideal-comparison}, used with two copies on the left,
then implies
\[
 ((a-r)_+)^{\oplus2}\precsim z_j^{\oplus n}\precsim a.
\]
Only scalar positives $z_j\in C$ have been tested against \propF{}.
Finally, $\|a-(a-r)_+\|\le r$.  As $r\downarrow0$, norm-limit closure
of comparison with the fixed target $a$ gives
\begin{equation}\label{eq:individual-infiniteness}
 a\oplus a\precsim a\qquad(a\in C_+).
\end{equation}
This includes $a=0$; when $n=1$, it also follows directly from the
hypothesis.  Therefore $C$ satisfies \propP{}.
Theorem~\ref{thm:individual-to-strong} now gives strong pure
infiniteness, completing the proof.
\end{proof}

The removal of amplification also gives the comparison form needed
for the direct quantitative argument in Section~\ref{sec:quantitative}.
The following conclusion concerns the entire closed ideal, rather
than only one cutoff or primitive norm slice.

\begin{corollary}[Whole-ideal comparison]\label{cor:ideal-comparison}
Let $C$ be a $C^*$-algebra satisfying \propP{}.  For every $h\in C_+$,
$c\in\Id(h)_+$, and integer $k\ge1$,
\[
 c^{\oplus k}\precsim h.
\]
If $C$ is exact, Proposition~\ref{prop:perturbation} may be applied
with multiplicity $n=1$, with the same constant $\sigma$ and all its
stated options.
\end{corollary}

\begin{proof}
Property \propP{} makes $h$ properly infinite, so
Lemma~\ref{lem:ideal-comparison} applies with multiplicity one.
If $h=0$, the ideal is zero and the assertion is immediate.
For exact $C$, property \propP{} is precisely the amplification
hypothesis of Proposition~\ref{prop:perturbation} with $n=1$ for every
scalar positive element.  Its constant was fixed independently of $n$,
so neither the constant nor the hereditary and selfadjoint options change.
\end{proof}

\section{Direct quantitative separation under exactness}\label{sec:quantitative}

We give a direct proof of strong pure infiniteness from the exact
perturbation principle.  Throughout this section, $C$ is exact and
satisfies $\propP$.  This is $\propF$ with $n=1$, so
Proposition~\ref{prop:perturbation} is available with that multiplicity
and its fixed constant $0<\sigma<1$.  Corollary~\ref{cor:ideal-comparison}
applies to the whole ideal $\Id(h)$ generated by each positive $h$.
For an exact algebra satisfying $\propF$, the scalar conclusion
\eqref{eq:individual-infiniteness} proved in Section~\ref{sec:localization}
supplies $\propP$ before that section invokes the general theorem.

The quantitative feature of this proof is a fixed lower norm bound in
primitive quotients, retained while a mixed product tends to zero.
We first construct positive elements in two prescribed hereditary
subalgebras with this property.  Spectral clipping then supplies bounded
comparison coefficients, so the small product survives Cuntz
implementation.  The resulting estimates verify the strong-infiniteness
inequalities directly.  Thus this proof uses the comparison tools and
the exact perturbation and localization results in
Sections~\ref{sec:foundations}, \ref{sec:exactness}, and~\ref{sec:localization}.

\begin{lemma}[Separation with a fixed quotient norm bound]
\label{lem:separating-positives}
Let $C$ be an exact $C^*$-algebra satisfying $\propP$, and let
$0<\sigma<1$ be the constant in Proposition~\ref{prop:perturbation}
with $n=1$.  Let $a,b\in C_+$ be contractions, and let $0<r<1$.  Set
\[
 m=\frac{r\sigma^2}{1024}.
\]
For every $0<\ell<1/10$ there are positives
$Q_a\in\Her(a)$ and $Q_b\in\Her(b)$ such that
\begin{align}
 \|(Q_a)_P\|&\ge m &&\text{if }\|a_P\|\ge r,\notag\\
 \|(Q_b)_P\|&\ge m &&\text{if }\|b_P\|\ge r,
 \label{eq:separation-quotient-bounds}\\
 \|Q_aQ_b\|&\le20\sqrt\ell.
 \label{eq:separation-product-bound}
\end{align}
Here and below $P$ ranges over $\Prim(C)$.
\end{lemma}

\begin{proof}
Write the polar decomposition in $C^{**}$ as
\[
 ab=vu,\qquad u=|ab|,\qquad u_*=|(ab)^*|.
\]
All three positive elements $a,b,u$ are contractions.  Choose continuous
functions on $[0,1]$ to separate the small singular values of $ab$
from a spectral corner on which an inverse square root is available.
The small values will give small cross products; on the remaining
corner we will preserve exact cancellation while perturbing the two
sides separately.  Use the following functions:
\begin{itemize}
\item $0\le L\le1$, $L=1$ on $[0,5\ell]$, and $L=0$ on
      $[6\ell,1]$; put $N=1-L$;
\item $0\le g\le1$, $g=0$ on $[0,\ell]$, and $g=1$ on
      $[2\ell,1]$;
\item $0\le R\le1$, $R=0$ on $[0,\ell/2]$, and
      $R(s)=\sqrt{\ell/s}$ for $s\ge\ell$.
\end{itemize}
Linear interpolation on the unspecified intervals gives such functions.
A function letter without an argument denotes evaluation at $u$.
$L$ selects the small singular values, $g$ confines the perturbations
away from zero, and $R$ normalizes $u$ on that region.
Functional calculus and the polar decomposition give
\begin{equation}\label{eq:separation-functional-identities}
 L(u_*)v=vL(u),\qquad N(u_*)=vNv^*,\qquad
 gN=N,\qquad (RuR-\ell)g=0.
\end{equation}
The second identity uses $N(0)=0$; the first remains valid with
$L(0)=1$.  Moreover
\[
 vR=ab\,\frac{R(u)}u\in C,\qquad
 vN=ab\,\frac{N(u)}u\in C,
\]
where both quotient functions are extended continuously by zero at
$s=0$, since $R$ and $N$ vanish near zero.
We use the unitization for scalar terms throughout this proof.

\medskip\noindent\textit{The four columns and their products.}
We first specify conditions under which four columns have the product
estimates we need.  Suppose that orthogonal positive contractions
$e,f\in C$ and invertibles $Z,W\in\widetilde C$ satisfy
\begin{equation}\label{eq:separation-inverse-conditions}
 ge=e,\quad gf=f,\quad Z^*W=1,\quad
 Z=1+gzg,\quad W=1+gwg,\quad
 \|Z\|,\|W\|\le2
\end{equation}
for some $z,w\in C$.  For such a choice, define
\begin{equation}\label{eq:separation-four-columns}
 \begin{array}{ll}
 X_0=aL(u_*),&X=\ell^{-1/2}avRZe,\\[2pt]
 Y_0=bL(u),&Y=\ell^{-1/2}bRWf.
 \end{array}
\end{equation}
These elements belong to $C$: the only polar factor is $vR\in C$,
and all unitized factors occur multiplied by elements of $C$.
For any such admissible choice, the proposed outputs are
\begin{equation}\label{eq:separation-final-positives}
 Q_a=(X_0X_0^*+XX^*)^{1/2},\qquad
 Q_b=(Y_0Y_0^*+YY^*)^{1/2}.
\end{equation}
The norm of each positive output dominates the norms of its two
columns in every primitive quotient.  At the end, a row polar
decomposition will bound the product of these outputs by the four
cross products estimated next.

Put $E_R=RuR-\ell$.  Since $E_Rg=0$ and $f=gf$, we have
\[
 E_RWf=E_Rf+E_Rgwgf=0.
\]
In particular, no commutation of $W$ with $g$ is needed.  Using
$v^*ab=u$ and retaining the indicated order of factors gives
\begin{equation}\label{eq:separation-high-orthogonality}
 X^*Y=\ell^{-1}eZ^*RuRWf=eZ^*Wf=ef=0.
\end{equation}
The other three products obey
\begin{equation}\label{eq:separation-mixed-bounds}
 \|X_0^*Y_0\|\le6\ell,\qquad
 \|X^*Y_0\|\le2\sqrt{6\ell},\qquad
 \|X_0^*Y\|\le2\sqrt{6\ell}.
\end{equation}
Indeed,
\begin{align*}
 X_0^*Y_0&=vL(u)uL(u),\\
 X^*Y_0&=\ell^{-1/2}eZ^*RuL(u),\\
 X_0^*Y&=\ell^{-1/2}vL(u)uRWf.
\end{align*}
Here $\|uL(u)\|\le6\ell$.  For $0\le s\le\ell$,
$\ell^{-1/2}R(s)sL(s)\le\sqrt\ell$; for $s\ge\ell$ the same
expression equals $\sqrt{s}L(s)\le\sqrt{6\ell}$.
Together with \eqref{eq:separation-inverse-conditions}, these scalar
estimates prove \eqref{eq:separation-mixed-bounds}.
We now construct $e,f,Z,W$ so that, in each prescribed primitive
quotient, at least one of the two relevant columns also has a fixed
lower norm bound.

\medskip\noindent\textit{Two orthogonal cores.}
Set $h_0=(u-2\ell)_+$.  Since $h_0\oplus h_0\precsim h_0$,
Lemma~\ref{lem:cut-implementation} supplies a row
$[r_1,r_2]\in M_{1,2}(C)$ such that
\begin{equation}\label{eq:separation-exact-cores}
 r_i^*r_j=\begin{cases}\theta(h_0),&i=j,\\0,&i\ne j,\end{cases}
 \qquad r_1r_1^*+r_2r_2^*\in\Her(h_0),
\end{equation}
where $\theta\colon[0,\infty)\to[0,1]$ is continuous, vanishes near
zero, and equals one on $[2\ell,\infty)$.  Define
\[
 e_0=r_1r_1^*,\qquad f_0=r_2r_2^*.
\]
They are orthogonal positive contractions in $\Her(h_0)$; each belongs
to this hereditary algebra because it is dominated by their sum.
In any primitive quotient with $\|u_P\|\ge4\ell$,
$\|(h_0)_P\|\ge2\ell$, and hence
\begin{equation}\label{eq:separation-core-norms}
 \|(e_0)_P\|=\|(f_0)_P\|=\|\theta((h_0)_P)\|=1.
\end{equation}
This includes equality at the threshold $4\ell$ by the chosen plateau.

Put
\[
 k_e=(e_0-1/2)_+,\qquad k_f=(f_0-1/2)_+,
\]
and let $e=\chi(e_0)$ and $f=\chi(f_0)$, where
$\chi\colon[0,1]\to[0,1]$ is continuous, vanishes near zero, and
is one on $[1/2,1]$.
Then
\begin{equation}\label{eq:separation-core-identities}
 ef=0,\qquad ek_e=k_e,\qquad fk_f=k_f,\qquad ge=e,\qquad gf=f.
\end{equation}
For the last two identities, $gh_0=h_0$ implies that multiplication
by $g$ is the identity on $\Her(h_0)$.
By \eqref{eq:separation-core-norms}, both $(k_e)_P$ and $(k_f)_P$
are nonzero whenever $\|u_P\|\ge4\ell$.

\medskip\noindent\textit{The quotient regions requiring a high column.}
Define
\begin{equation}\label{eq:separation-normalization-data}
 \tau=\frac{r\sqrt\ell}{8},\qquad
 C_1=avR,\quad C_2=bR,\qquad
 h_i=(|C_ig|-\tau)_+,\quad d_i=(|C_iN|-2\tau)_+.
\end{equation}
We first check why a lower bound on the region
$\|(d_i)_P\|\ge\tau$ will recover any signal missed by the
corresponding low column.
Suppose $\|a_P\|\ge r$ and $\|(X_0)_P\|<r/2$.
Since $a=aL(u_*)+aN(u_*)$, we have
$\|(aN(u_*))_P\|>r/2$.
The identity $N(u_*)=vNv^*$ gives $\|(avN)_P\|>r/2$.
On $[0,1]$, the scalar inequality $R^2N^2\ge\ell N^2$ holds:
where $N$ is nonzero, $s>5\ell$ and $R(s)^2=\ell/s\ge\ell$.
Consequently
\[
 (C_1N)(C_1N)^*
 =avR^2N^2v^*a
 \ge\ell avN^2v^*a,
\]
and therefore
\[
 \|(C_1N)_P\|>\frac{r\sqrt\ell}{2}=4\tau.
\]
It follows that $\|(d_1)_P\|>2\tau$.
Similarly, if $\|b_P\|\ge r$ and $\|(Y_0)_P\|<r/2$,
then $\|(bN)_P\|>r/2$.  The order inequality
$bR^2N^2b\ge\ell bN^2b$ gives $\|(d_2)_P\|>2\tau$.
Thus every prescribed quotient either has a low column of norm at
least $r/2$ or belongs to the corresponding perturbation region.

\medskip\noindent\textit{Comparison inputs and normalized left factors.}
We now verify the inputs to Proposition~\ref{prop:perturbation}
on these two regions.
If $(d_i)_P\ne0$, then $C_iN=C_igN$ implies
$\|(C_ig)_P\|>2\tau$, so $(h_i)_P\ne0$.  The same assumption
implies $N_P\ne0$, and therefore $\|u_P\|>5\ell$.
Consequently $(k_e)_P,(k_f)_P\ne0$.
The primitive ideal test and Corollary~\ref{cor:ideal-comparison} give
\begin{equation}\label{eq:separation-core-comparisons}
 d_1\precsim h_1,\quad d_1\precsim k_e,\qquad
 d_2\precsim h_2,\quad d_2\precsim k_f.
\end{equation}

Let $D_i=C_ig$, and use continuous functional calculus in the
unitization to set
\begin{equation}\label{eq:separation-left-normalization}
 F_i=\bigl(\max\{\tau,|D_i^*|\}\bigr)^{-1},\qquad
 p_i=F_iD_i.
\end{equation}
Thus $\|F_i\|\le\tau^{-1}$ and $p_i\in C$.
The functional-calculus intertwining identity
\[
 F_iD_i=D_i\bigl(\max\{\tau,|D_i|\}\bigr)^{-1}
\]
follows first for polynomials in $D_iD_i^*$ and $D_i^*D_i$, and then
by uniform approximation.  In particular,
\begin{equation}\label{eq:separation-normalized-core}
 p_i^*p_i=|D_i|^2\bigl(\max\{\tau,|D_i|\}\bigr)^{-2}\le1,
 \qquad p_i^*p_ih_i=h_i.
\end{equation}
The use of $|D_i^*|$ in the left factor $F_i$ is essential to this
identity.

\medskip\noindent\textit{Two compatible inverse perturbations.}
Fix
\[
 \eta=\frac18,\qquad \eta'=\frac{\sigma\eta}{16}
 =\frac{\sigma}{128}.
\]
Apply Proposition~\ref{prop:perturbation}, with $n=1$, to
\[
 (p,q,h,k,c,T,\alpha)
 =(p_1,e,h_1,k_e,d_1,F_1C_1e,\tau)
\]
at scale $\eta$.  Its hypotheses follow from
\eqref{eq:separation-core-identities},
\eqref{eq:separation-core-comparisons}, and
\eqref{eq:separation-normalized-core}.
For the resulting contraction $y\in C$, put
\[
 Z_1=1+\eta gyg.
\]
Since $\|Z_1-1\|\le\eta<1$, the element $Z_1$ is invertible;
define $W_1=(Z_1^*)^{-1}$.  Since $ge=e$,
the conclusion of the Proposition becomes
\begin{equation}\label{eq:separation-first-test}
 \|(F_1C_1Z_1e)_P\|\ge\sigma\eta
 \qquad\text{if }\|(d_1)_P\|\ge\tau.
\end{equation}
Next apply the same Proposition at scale $\eta'$ to
\[
 (p,q,h,k,c,T,\alpha)
 =(p_2,f,h_2,k_f,d_2,F_2C_2W_1f,\tau).
\]
The element $T$ is in $C$; the Proposition imposes no bound on its
norm.  With the resulting contraction $y'\in C$, set
\[
 W=W_1+\eta'gy'g.
\]
Since $gf=f$, the second conclusion is
\begin{equation}\label{eq:separation-second-test}
 \|(F_2C_2Wf)_P\|\ge\sigma\eta'
 \qquad\text{if }\|(d_2)_P\|\ge\tau.
\end{equation}

We record the inverse estimates, including the factorizations needed
to preserve \eqref{eq:separation-first-test}.
For $a_0\in C$ with $\|a_0\|<1$, the Neumann series gives
\begin{equation}\label{eq:separation-factored-series}
 (1+ga_0g)^{-1}=1+g\left(
 \sum_{k=1}^{\infty}(-1)^k a_0(g^2a_0)^{k-1}\right)g.
\end{equation}
The series of middle factors converges in $C$ with norm at most
$\|a_0\|/(1-\|a_0\|)$.  This formula does not require $g$ to be
a projection.  Applying it to $Z_1^*$ gives
\[
 W_1=1+gw_1g,\qquad \|w_1\|\le\frac17.
\]
Thus $W=1+gwg$ for $w=w_1+\eta'y'\in C$, with
\[
 \|w\|\le\frac17+\frac1{128}<\frac14.
\]
Applying \eqref{eq:separation-factored-series} to $W^*=1+gw^*g$ now
proves that $W$ is invertible.  Define $Z=(W^*)^{-1}$.  The same
series gives
\begin{equation}\label{eq:separation-inverse-bounds}
 Z=1+gzg,\qquad \|z\|\le\frac13,
 \qquad \|Z\|\le\frac43,
 \quad\|W\|\le\frac54,
 \quad\|Z_1\|\le\frac98.
\end{equation}
In particular, all conditions in
\eqref{eq:separation-inverse-conditions} hold.

The second perturbation changes the inverse used in the first column.
To keep the first lower bound independent of $\ell$, we must estimate
this change without the potentially large factor $\|F_1\|$.
The inverse difference has the exact order
\[
 Z-Z_1=-\eta' Zg(y')^*gZ_1.
\]
As $Zg=g(1+zg^2)$, its effect on the first normalized column is
bounded by
\begin{align}
 \|F_1C_1(Z-Z_1)e\|
 &=\eta'\|p_1(1+zg^2)(y')^*gZ_1e\|\notag\\
 &\le\eta'\left(1+\frac13\right)\frac98
 =\frac32\eta'.
 \label{eq:separation-no-normalization-loss}
\end{align}
Since $3\eta'/2\le\sigma\eta/2$,
\eqref{eq:separation-first-test} remains bounded below by
$\sigma\eta/2$ with $Z$ in place of $Z_1$.
Using $\|F_i\|\le\tau^{-1}$ in each primitive quotient, and
$\tau\ell^{-1/2}=r/8$, we obtain
\begin{align}
 \|X_P\|&\ge\frac r8\frac{\sigma\eta}{2}
 &&\text{if }\|(d_1)_P\|\ge\tau,\notag\\
 \|Y_P\|&\ge\frac r8\sigma\eta'
 &&\text{if }\|(d_2)_P\|\ge\tau.
 \label{eq:separation-high-lower-bounds}
\end{align}
The cancellation $\tau\ell^{-1/2}=r/8$ is what removes
$\ell$ from both lower bounds.

\medskip\noindent\textit{Assembling the two positive outputs.}

To check the hereditary membership of $Q_a,Q_b$,
\eqref{eq:separation-four-columns} and the bounds on $Z,W$ give
\[
 X_0X_0^*\le a^2,\quad XX^*\le4\ell^{-1}a^2,\qquad
 Y_0Y_0^*\le b^2,\quad YY^*\le4\ell^{-1}b^2.
\]
By Lemma~\ref{lem:comparison-rules}, the first two summands belong
to $\Her(a)$ and the last two to $\Her(b)$.
Thus $Q_a\in\Her(a)$ and $Q_b\in\Her(b)$.
The earlier coverage reduction and
\eqref{eq:separation-high-lower-bounds}, together with domination of
the column norms by the output norms, prove
\eqref{eq:separation-quotient-bounds} with
\[
 m=\frac r8\sigma\eta'=\frac{r\sigma^2}{1024}.
\]
Here $m\le r/2$ and $\eta'\le\eta/2$, so this single lower bound
applies to all cases.  It is independent of $\ell$.

Finally put $U_a=[X_0,X]$ and $U_b=[Y_0,Y]$ in $M_{1,2}(C)$.
Their range squares are $Q_a^2$ and $Q_b^2$.
The rectangular polar decompositions in the bidual have the form
$U_a=Q_av_a$ and $U_b=Q_bv_b$, where
$v_av_a^*$ and $v_bv_b^*$ are the support projections of $Q_a$
and $Q_b$, respectively.  Hence both identities
\[
 U_a^*U_b=v_a^*Q_aQ_bv_b,\qquad
 Q_aQ_b=v_a(U_a^*U_b)v_b^*
\]
hold.  Since partial isometries have norm at most one, they imply
\[
 \|Q_aQ_b\|=\|U_a^*U_b\|
 \le6\ell+4\sqrt{6\ell}\le20\sqrt\ell,
\]
by \eqref{eq:separation-high-orthogonality} and
\eqref{eq:separation-mixed-bounds}.
This proves \eqref{eq:separation-product-bound}.

The construction also covers zero elements and empty norm regions.
If $h_0=0$, the row in \eqref{eq:separation-exact-cores} may be zero.
Empty perturbation targets permit zero perturbations, and the coverage
argument then assigns every prescribed quotient to a low column.
All inverse operations have either the positive floor $\tau$ or the
convergent series \eqref{eq:separation-factored-series}.
\end{proof}

The norms of the preceding outputs may grow as $\ell$ decreases.
The next observation controls the comparison coefficients using only
the fixed spectral floor $m$.  This is enough for their small product
to survive Cuntz implementation.

\begin{lemma}[Comparison above a spectral floor]
\label{lem:bounded-implementation}
Let $C$ be a $C^*$-algebra, let $Q,c\in C_+$ with $\|c\|\le1$,
and let $m>0$.  If
$c\precsim(Q^2-m^2/4)_+$, then for each $0<\delta<1$ there is
$x\in C$ such that
\begin{equation}\label{eq:separation-bounded-implementation}
 \|(Qx)^*(Qx)-c\|<\delta,\qquad
 \|Qx\|\le2,\qquad \|x\|\le\frac4m.
\end{equation}
\end{lemma}

\begin{proof}
Choose a comparison implementer $v\in C$ with
\[
 \|v^*(Q^2-m^2/4)_+v-c\|<\delta.
\]
Such an implementer lies in $C$ by compression to the two scalar
matrix corners.  On $[0,\|Q\|]$ define the continuous function
\[
 j(s)=\begin{cases}s^{-1}(s^2-m^2/4)_+^{1/2},&s>0,\\0,&s=0,
 \end{cases}
 \qquad x=j(Q)v.
\]
Then $(Qx)^*(Qx)=v^*(Q^2-m^2/4)_+v$ and
$\|Qx\|^2<1+\delta<2$, which in particular gives $\|Qx\|\le2$.
Moreover the scalar inequality
\[
 j(s)^2\le\frac4{m^2}s^2j(s)^2
\]
holds because $j$ vanishes on $[0,m/2]$.
Functional calculus followed by multiplication by $v^*$ and $v$ gives
\[
 x^*x\le\frac4{m^2}(Qx)^*(Qx).
\]
Thus $\|x\|\le(2/m)\|Qx\|\le4/m$, independently of the norm
of the original comparison implementer $v$.
\end{proof}

\begin{proposition}[Direct exact proof]\label{prop:direct-strong-infiniteness}
Let $C$ be an exact $C^*$-algebra satisfying $\propP$.  For every
$a,b\in C_+$ and every $\varepsilon>0$, there are $s,t\in C$ such that
\[
 \|s^*a^2s-a^2\|<\varepsilon,\qquad
 \|t^*b^2t-b^2\|<\varepsilon,\qquad
 \|s^*abt\|<\varepsilon.
\]
Consequently $C$ is strongly purely infinite.
\end{proposition}

\begin{proof}
Fix the constant $\sigma$ in Proposition~\ref{prop:perturbation} with
$n=1$.
Let $a,b\in C_+$ be contractions and let $\varepsilon>0$.
Choose $0<r<1$ with $r^2<\varepsilon/3$ and put
$m=r\sigma^2/1024$.  Next choose $0<\ell<1/10$ so small that
\begin{equation}\label{eq:separation-final-ell}
 \frac{320\sqrt\ell}{m^2}<\frac\varepsilon3.
\end{equation}
Take $Q_a,Q_b$ from Lemma~\ref{lem:separating-positives}.
If $(a^2-r^2)_+$ has nonzero image at $P$, then $\|a_P\|>r$,
so $\|(Q_a)_P\|\ge m$.  Thus $(Q_a^2-m^2/4)_+$ has nonzero
image there.  The primitive ideal test and
Corollary~\ref{cor:ideal-comparison} imply
\begin{equation}\label{eq:separation-final-comparisons}
 (a^2-r^2)_+\precsim(Q_a^2-m^2/4)_+,\qquad
 (b^2-r^2)_+\precsim(Q_b^2-m^2/4)_+.
\end{equation}

Finally choose $0<\delta<1$ such that
\begin{equation}\label{eq:separation-final-delta}
 \delta+(4+\delta)\delta<\frac\varepsilon3.
\end{equation}
Apply Lemma~\ref{lem:bounded-implementation} to the comparisons in
\eqref{eq:separation-final-comparisons}.  We obtain $x_a,x_b\in C$
with
\begin{align*}
 \|(Q_ax_a)^*(Q_ax_a)-(a^2-r^2)_+\|&<\delta,\\
 \|(Q_bx_b)^*(Q_bx_b)-(b^2-r^2)_+\|&<\delta,
\end{align*}
and
\[
 \|Q_ax_a\|,\|Q_bx_b\|\le2,\qquad
 \|x_a\|,\|x_b\|\le\frac4m.
\]
Consequently
\begin{equation}\label{eq:separation-implemented-product}
 \|x_a^*Q_aQ_bx_b\|
 \le\frac{16}{m^2}\|Q_aQ_b\|
 \le\frac{320\sqrt\ell}{m^2}.
\end{equation}

Since $Q_a\in\Her(a)=\overline{aCa}$,
$Q_ax_a\in\overline{aC}$; similarly $Q_bx_b\in\overline{bC}$.
Choose $s,t\in C$ such that
\[
 \|as-Q_ax_a\|<\delta,\qquad
 \|bt-Q_bx_b\|<\delta.
\]
For elements $U,V$ of norm at most $2$, replacing them by elements
within $\delta$ changes $U^*U$ or $U^*V$ by at most
$(4+\delta)\delta$, by expanding the difference and using the
triangle inequality.  Also
$\|a^2-(a^2-r^2)_+\|,\|b^2-(b^2-r^2)_+\|\le r^2$.
Therefore
\begin{align*}
 \|s^*a^2s-a^2\|,
 \ \|t^*b^2t-b^2\|
 &\le r^2+\delta+(4+\delta)\delta<\varepsilon,\\
 \|s^*abt\|
 &\le\frac{320\sqrt\ell}{m^2}+(4+\delta)\delta<\varepsilon,
\end{align*}
by \eqref{eq:separation-final-ell} and
\eqref{eq:separation-final-delta}.
The choice order is $r$, then $m$, then $\ell$, then $\delta$;
the lower bound $m$ has no dependence on the later parameters.

For arbitrary positives $a,b$, put $M=\max\{1,\|a\|,\|b\|\}$.
Apply the preceding result to $a/M,b/M$ with tolerance
$\varepsilon/M^2$.  Multiplication of all three inequalities by
$M^2$ gives the required inequalities for $a,b$ with the same
$s,t\in C$.  Zero elements, including the zero algebra, are included.
These are precisely the three inequalities in \eqref{eq:SPI-definition},
so $C$ is strongly purely infinite.
\end{proof}


\bibliographystyle{plainnat}
\bibliography{references}
\end{document}
